% oss-plantilla: oss-2026.11
% Preámbulo corporativo fijo. No editar en sd-NN.tex: exportar_latex.py verificar lo compara byte a byte.
% Para cambiar el formato, editar esta plantilla u oss.sty y subir la versión de la primera línea.
\PassOptionsToPackage{unicode}{hyperref}
\PassOptionsToPackage{hyphens}{url}
\documentclass[11pt]{article}
\usepackage{amsmath,amssymb}
\usepackage{iftex}
\ifPDFTeX
  % Prism y otros editores en la nube: pdfLaTeX con Helvetica como familia por defecto.
  \usepackage[T1]{fontenc}
  \usepackage[utf8]{inputenc}
  \usepackage{textcomp}
  \usepackage[scaled=0.92]{helvet}
  \renewcommand{\familydefault}{\sfdefault}
  \DeclareUnicodeCharacter{2264}{\ensuremath{\leq}}
  \DeclareUnicodeCharacter{2265}{\ensuremath{\geq}}
  \DeclareUnicodeCharacter{2260}{\ensuremath{\neq}}
  \DeclareUnicodeCharacter{2248}{\ensuremath{\approx}}
  \DeclareUnicodeCharacter{2190}{\ensuremath{\leftarrow}}
  \DeclareUnicodeCharacter{2192}{\ensuremath{\rightarrow}}
  \DeclareUnicodeCharacter{25CF}{\ensuremath{\bullet}}
  \DeclareUnicodeCharacter{2713}{\ensuremath{\checkmark}}
\else
  % XeLaTeX/LuaLaTeX (compilación oficial local): DejaVu Sans.
  \usepackage{fontspec}
  \defaultfontfeatures{Ligatures=TeX,Scale=MatchLowercase}
  \setmainfont{DejaVu Sans}
  \setsansfont{DejaVu Sans}
  \setmonofont{DejaVu Sans Mono}
\fi
\usepackage[letterpaper,margin=20mm]{geometry}
\usepackage{microtype}
\usepackage{parskip}
\usepackage{xcolor}
\usepackage{longtable,booktabs,array}
\usepackage{calc}
\usepackage{etoolbox}
\makeatletter
\patchcmd\longtable{\par}{\if@noskipsec\mbox{}\fi\par}{}{}
\makeatother
\usepackage{footnotehyper}
\makesavenoteenv{longtable}
\usepackage{graphicx}
\makeatletter
\newsavebox\pandoc@box
\newcommand*\pandocbounded[1]{%
  \sbox\pandoc@box{#1}%
  \Gscale@div\@tempa{\textheight}{\dimexpr\ht\pandoc@box+\dp\pandoc@box\relax}%
  \Gscale@div\@tempb{\linewidth}{\wd\pandoc@box}%
  \ifdim\@tempb\p@<\@tempa\p@\let\@tempa\@tempb\fi
  \ifdim\@tempa\p@<\p@\scalebox{\@tempa}{\usebox\pandoc@box}%
  \else\usebox{\pandoc@box}%
  \fi}
\def\maxwidth{\ifdim\Gin@nat@width>\linewidth\linewidth\else\Gin@nat@width\fi}
\def\maxheight{\ifdim\Gin@nat@height>\textheight\textheight\else\Gin@nat@height\fi}
\makeatother
\setkeys{Gin}{width=\maxwidth,height=\maxheight,keepaspectratio}
\makeatletter
\def\fps@figure{htbp}
\makeatother
\usepackage[normalem]{ulem}
\providecommand{\st}[1]{\sout{#1}}
\usepackage{fancyvrb}
\setlength{\emergencystretch}{3em}
\providecommand{\tightlist}{%
  \setlength{\itemsep}{0pt}\setlength{\parskip}{0pt}}
\providecommand{\passthrough}[1]{#1}
\setcounter{secnumdepth}{-\maxdimen}
\ifPDFTeX
  \usepackage[spanish]{babel}
\else
  \usepackage[bidi=default]{babel}
  \babelprovide[main,import]{spanish}
\fi
\let\LanguageShortHands\languageshorthands
\def\languageshorthands#1{}
\usepackage{bookmark}
\usepackage{xurl}
\urlstyle{same}
\hypersetup{pdflang={es-CL},pdfcreator={Only Simple Solutions - exportar_latex.py}}
\usepackage{oss}
\begin{document}

\hypersetup{pageanchor=false}
\ossCover{Matriz de cumplimiento técnico y trazabilidad}{Subdocumento 3 · T-12}
\hypersetup{pageanchor=true}
\tableofcontents
\clearpage
\section*{Lista de tablas}
\addcontentsline{toc}{section}{Lista de tablas}

\noindent\hangindent=2em\hyperref[tab:t12-1]{Tabla T-12.1: Matriz de cumplimiento y trazabilidad}\dotfill\pageref{tab:t12-1}\par\vspace{0.4\baselineskip}
\clearpage


\renewcommand{\tablename}{Tabla}
\renewcommand{\thetable}{T-12.\arabic{table}}

\section{Formulario T-12 · Matriz de cumplimiento y trazabilidad}\label{formulario-t-12}

Esta matriz responde cada requerimiento del catálogo vigente (versión 3.1, 308 elementos: 227 funcionales, 72 no funcionales y 9 obligaciones del proponente) y cada requisito técnico de las Bases Transversales (374). La columna Cumple usa cuatro valores: Sí, cuando la propuesta lo satisface y la sección citada lo respalda; Sí, condicionado, cuando lo satisface a condición de una decisión o prueba que la sección señala; Parcial, cuando la propuesta cubre solo una parte y la fila dice cuál falta; y No, cuando aún no está cubierto y la fila indica en qué subdocumento se completará. Las referencias a tablas y secciones son las de los Subdocumentos 3, 4 y 5 vigentes. Los productos y marcas que los subdocumentos no han decidido figuran como por definir.


\begin{longtable}[]{@{}
  >{\raggedright\arraybackslash\footnotesize}p{(\columnwidth - 10\tabcolsep) * \real{0.1100}}
  >{\raggedright\arraybackslash\footnotesize}p{(\columnwidth - 10\tabcolsep) * \real{0.2900}}
  >{\raggedright\arraybackslash\footnotesize}p{(\columnwidth - 10\tabcolsep) * \real{0.1100}}
  >{\raggedright\arraybackslash\footnotesize}p{(\columnwidth - 10\tabcolsep) * \real{0.2700}}
  >{\raggedright\arraybackslash\footnotesize}p{(\columnwidth - 10\tabcolsep) * \real{0.2200}}@{}}
\caption{Matriz de cumplimiento y trazabilidad.}\label{tab:t12-1}\\
\toprule\noalign{}
\begin{minipage}[b]{\linewidth}\raggedright
\textbf{ID Requerimiento}
\end{minipage} & \begin{minipage}[b]{\linewidth}\raggedright
\textbf{Descripción}
\end{minipage} & \begin{minipage}[b]{\linewidth}\raggedright
\textbf{Cumple}
\end{minipage} & \begin{minipage}[b]{\linewidth}\raggedright
\textbf{Componente que lo satisface}
\end{minipage} & \begin{minipage}[b]{\linewidth}\raggedright
\textbf{Sección de la propuesta}
\end{minipage} \\
\midrule\noalign{}
\endfirsthead
\toprule\noalign{}
\begin{minipage}[b]{\linewidth}\raggedright
\textbf{ID Requerimiento}
\end{minipage} & \begin{minipage}[b]{\linewidth}\raggedright
\textbf{Descripción}
\end{minipage} & \begin{minipage}[b]{\linewidth}\raggedright
\textbf{Cumple}
\end{minipage} & \begin{minipage}[b]{\linewidth}\raggedright
\textbf{Componente que lo satisface}
\end{minipage} & \begin{minipage}[b]{\linewidth}\raggedright
\textbf{Sección de la propuesta}
\end{minipage} \\
\midrule\noalign{}
\endhead
\bottomrule\noalign{}
\endlastfoot
RF-001 & El sistema deberá permitir la venta física de la unidad aunque exista una reserva temporal activa en el canal digital. & Sí & Servicio de pedidos (reserva temporal / prelación de canal presencial) & Subdoc.3 §3.3.2 (Frente Retail); Anexo A.4.1 (SUP-12) \\
\addlinespace[2pt]
RF-002 & El sistema deberá exigir credenciales individuales al usuario antes de habilitar cualquier acción. & Sí & Base tecnológica (identidad y gestión de accesos) — traspaso de sesión, revocación por finiquito, aprovisionamiento B2B & Subdoc.3 §3.3.2 (plataforma común); Anexo A.4.1 (SUP-19, SUP-24) \textbar{} Subdoc.4 §4.1 \\
\addlinespace[2pt]
RF-003 & El sistema debe impedir el acceso anónimo a una terminal. & Sí & Base tecnológica (identidad y gestión de accesos) — traspaso de sesión, revocación por finiquito, aprovisionamiento B2B & Subdoc.3 §3.3.2 (plataforma común); Anexo A.4.1 (SUP-19, SUP-24) \textbar{} Subdoc.4 §4.1 \\
\addlinespace[2pt]
RF-004 & El sistema debe impedir el acceso mediante credencial compartida. & Sí & Base tecnológica (identidad y gestión de accesos) — traspaso de sesión, revocación por finiquito, aprovisionamiento B2B & Subdoc.3 §3.3.2 (plataforma común); Anexo A.4.1 (SUP-19, SUP-24) \textbar{} Subdoc.4 §4.1 \\
\addlinespace[2pt]
RF-005 & El sistema deberá permitir el traspaso de sesión nominativa entre usuarios en una terminal compartida. & Sí & Base tecnológica (identidad y gestión de accesos) — traspaso de sesión, revocación por finiquito, aprovisionamiento B2B & Subdoc.3 §3.3.2 (plataforma común); Anexo A.4.1 (SUP-19, SUP-24) \textbar{} Subdoc.4 §4.1 \\
\addlinespace[2pt]
RF-006 & El sistema deberá cerrar automáticamente la sesión por inactividad en la terminal compartida. & Sí & Base tecnológica (identidad y gestión de accesos) — traspaso de sesión, revocación por finiquito, aprovisionamiento B2B & Subdoc.3 §3.3.2 (plataforma común); Anexo A.4.1 (SUP-19, SUP-24) \textbar{} Subdoc.4 §4.1 \\
\addlinespace[2pt]
RF-007 & El sistema deberá impedir la ejecución de la función de originación a un usuario sin capacitación normativa acreditada y vigente. & Sí & Base tecnológica (identidad y gestión de accesos) — traspaso de sesión, revocación por finiquito, aprovisionamiento B2B & Subdoc.3 §3.3.2 (plataforma común); Anexo A.4.1 (SUP-19, SUP-24) \textbar{} Subdoc.4 §4.1 \\
\addlinespace[2pt]
RF-008 & El sistema deberá impedir la ejecución de la función de repactación a un usuario sin capacitación normativa acreditada y vigente. & Sí & Base tecnológica (identidad y gestión de accesos) — traspaso de sesión, revocación por finiquito, aprovisionamiento B2B & Subdoc.3 §3.3.2 (plataforma común); Anexo A.4.1 (SUP-19, SUP-24) \textbar{} Subdoc.4 §4.1 \\
\addlinespace[2pt]
RF-009 & El sistema deberá revocar automáticamente la habilitación funcional del usuario al vencer su acreditación de capacitación. & Sí & Base tecnológica (identidad y gestión de accesos) — traspaso de sesión, revocación por finiquito, aprovisionamiento B2B & Subdoc.3 §3.3.2 (plataforma común); Anexo A.4.1 (SUP-19, SUP-24) \textbar{} Subdoc.4 §4.1 \\
\addlinespace[2pt]
RF-010 & El sistema deberá permitir al proveedor patrocinar la habilitación temporal de acceso de su repositor externo, con fecha de término declarada. & Sí & Base tecnológica (identidad y gestión de accesos) — traspaso de sesión, revocación por finiquito, aprovisionamiento B2B & Subdoc.3 §3.3.2 (plataforma común); Anexo A.4.1 (SUP-19, SUP-24) \textbar{} Subdoc.4 §4.1 \\
\addlinespace[2pt]
RF-011 & El sistema deberá caducar automáticamente la habilitación de acceso del personal externo al cumplirse su vigencia. & Sí & Base tecnológica (identidad y gestión de accesos) — traspaso de sesión, revocación por finiquito, aprovisionamiento B2B & Subdoc.3 §3.3.2 (plataforma común); Anexo A.4.1 (SUP-19, SUP-24) \textbar{} Subdoc.4 §4.1 \\
\addlinespace[2pt]
RF-012 & El sistema deberá registrar de forma individualizada el acceso y la actividad del personal externo. & Sí & Base tecnológica (identidad y gestión de accesos) — traspaso de sesión, revocación por finiquito, aprovisionamiento B2B & Subdoc.3 §3.3.2 (plataforma común); Anexo A.4.1 (SUP-19, SUP-24) \textbar{} Subdoc.4 §4.1 \\
\addlinespace[2pt]
RF-013 & El sistema deberá revocar la totalidad de los accesos y credenciales del trabajador a partir del término efectivo de su vínculo. & Sí & Base tecnológica (identidad y gestión de accesos) — traspaso de sesión, revocación por finiquito, aprovisionamiento B2B & Subdoc.3 §3.3.2 (plataforma común); Anexo A.4.1 (SUP-19, SUP-24) \textbar{} Subdoc.4 §4.1 \\
\addlinespace[2pt]
RF-014 & El sistema debe conciliar los accesos vigentes contra la nómina activa. & Sí & Base tecnológica (identidad y gestión de accesos) — traspaso de sesión, revocación por finiquito, aprovisionamiento B2B & Subdoc.3 §3.3.2 (plataforma común); Anexo A.4.1 (SUP-19, SUP-24) \textbar{} Subdoc.4 §4.1 \\
\addlinespace[2pt]
RF-015 & El sistema debe informar al oficial de seguridad los accesos huérfanos detectados en la conciliación. & Sí & Base tecnológica (identidad y gestión de accesos) — traspaso de sesión, revocación por finiquito, aprovisionamiento B2B & Subdoc.3 §3.3.2 (plataforma común); Anexo A.4.1 (SUP-19, SUP-24) \textbar{} Subdoc.4 §4.1 \\
\addlinespace[2pt]
RF-016 & El sistema deberá propagar el cambio de precio a las 380 líneas de caja. & Sí & Servicio de oferta comercial & Subdoc.3 §3.3.2 (Frente Retail); Anexo A.4.1 (SUP-08, SUP-09, SUP-10) \textbar{} Subdoc.4 §4.1 (Tabla 4.2) \\
\addlinespace[2pt]
RF-017 & El sistema deberá propagar el cambio de precio al canal digital. & Sí & Servicio de oferta comercial & Subdoc.3 §3.3.2 (Frente Retail); Anexo A.4.1 (SUP-08, SUP-09, SUP-10) \textbar{} Subdoc.4 §4.1 (Tabla 4.2) \\
\addlinespace[2pt]
RF-018 & El sistema deberá propagar el cambio de precio a los puntos de exhibición física de cada tienda afectada. & Sí & Servicio de oferta comercial & Subdoc.3 §3.3.2 (Frente Retail); Anexo A.4.1 (SUP-08, SUP-09, SUP-10) \textbar{} Subdoc.4 §4.1 (Tabla 4.2) \\
\addlinespace[2pt]
RF-019 & El sistema deberá permitir al reponedor registrar la referencia y la tienda de la etiqueta cambiada, con el instante del cambio. & Sí & Servicio de oferta comercial & Subdoc.3 §3.3.2 (Frente Retail); Anexo A.4.1 (SUP-08, SUP-09, SUP-10) \textbar{} Subdoc.4 §4.1 (Tabla 4.2) \\
\addlinespace[2pt]
RF-020 & El sistema deberá registrar la identidad individual del ejecutor del cambio de etiqueta. & Sí & Servicio de oferta comercial & Subdoc.3 §3.3.2 (Frente Retail); Anexo A.4.1 (SUP-08, SUP-09, SUP-10) \textbar{} Subdoc.4 §4.1 (Tabla 4.2) \\
\addlinespace[2pt]
RF-021 & El sistema deberá permitir al ejecutivo de cumplimiento recuperar el precio publicado de una referencia para una fecha, hora y canal determinados. & Sí & Servicio de oferta comercial & Subdoc.3 §3.3.2 (Frente Retail); Anexo A.4.1 (SUP-08, SUP-09, SUP-10) \textbar{} Subdoc.4 §4.1 (Tabla 4.2) \\
\addlinespace[2pt]
RF-022 & El sistema deberá obtener el precio vigente de la referencia según el último precio propagado a la etiqueta de esa tienda en la fecha de la venta. & Sí & Servicio de oferta comercial & Subdoc.3 §3.3.2 (Frente Retail); Anexo A.4.1 (SUP-08, SUP-09, SUP-10) \textbar{} Subdoc.4 §4.1 (Tabla 4.2) \\
\addlinespace[2pt]
RF-023 & El sistema deberá cobrar el menor de ambos precios para el consumidor. & Sí & Servicio de oferta comercial & Subdoc.3 §3.3.2 (Frente Retail); Anexo A.4.1 (SUP-08, SUP-09, SUP-10) \textbar{} Subdoc.4 §4.1 (Tabla 4.2) \\
\addlinespace[2pt]
RF-024 & El sistema deberá generar un registro de incidente de discrepancia de precio consultable. & Sí & Servicio de oferta comercial & Subdoc.3 §3.3.2 (Frente Retail); Anexo A.4.1 (SUP-08, SUP-09, SUP-10) \textbar{} Subdoc.4 §4.1 (Tabla 4.2) \\
\addlinespace[2pt]
RF-028 & El sistema deberá impedir la venta de la referencia al precio nuevo mientras su punto de exhibición no confirme la actualización, salvo que el artículo se encuentre bloqueado para venta. & Sí & Servicio de oferta comercial & Subdoc.3 §3.3.2 (Frente Retail); Anexo A.4.1 (SUP-08, SUP-09, SUP-10) \textbar{} Subdoc.4 §4.1 (Tabla 4.2) \\
\addlinespace[2pt]
RF-029 & El sistema deberá generar un indicador diario de puntos de exhibición desactualizados con identificación individual de cada punto. & Sí & Servicio de oferta comercial & Subdoc.3 §3.3.2 (Frente Retail); Anexo A.4.1 (SUP-08, SUP-09, SUP-10) \textbar{} Subdoc.4 §4.1 (Tabla 4.2) \\
\addlinespace[2pt]
RF-030 & El sistema deberá agrupar los cambios de precio de sala en ventanas parametrizadas fuera del horario de atención. & Sí & Servicio de oferta comercial & Subdoc.3 §3.3.2 (Frente Retail); Anexo A.4.1 (SUP-08, SUP-09, SUP-10) \textbar{} Subdoc.4 §4.1 (Tabla 4.2) \\
\addlinespace[2pt]
RF-032 & El sistema deberá evaluar la vigencia de cada promoción al instante de emisión del documento. & Sí & Servicio de oferta comercial & Subdoc.3 §3.3.2 (Frente Retail); Anexo A.4.1 (SUP-08, SUP-09, SUP-10) \textbar{} Subdoc.4 §4.1 (Tabla 4.2) \\
\addlinespace[2pt]
RF-033 & El sistema deberá evaluar la aplicabilidad de la promoción según el canal de la venta. & Sí & Servicio de oferta comercial & Subdoc.3 §3.3.2 (Frente Retail); Anexo A.4.1 (SUP-08, SUP-09, SUP-10) \textbar{} Subdoc.4 §4.1 (Tabla 4.2) \\
\addlinespace[2pt]
RF-034 & El sistema deberá evaluar la aplicabilidad de la promoción según la tienda en que se realiza la venta. & Sí & Servicio de oferta comercial & Subdoc.3 §3.3.2 (Frente Retail); Anexo A.4.1 (SUP-08, SUP-09, SUP-10) \textbar{} Subdoc.4 §4.1 (Tabla 4.2) \\
\addlinespace[2pt]
RF-035 & El sistema deberá crear una reserva temporal para la unidad agregada al carro. & Sí & Servicio de pedidos / Servicio de comisiones / Canal digital (comercio electrónico vigente) & Subdoc.3 §3.3.2 (Frente Retail); Anexo A.4.1 (SUP-11, SUP-12, SUP-13, SUP-14) \textbar{} Subdoc.4 §4.1 (Tabla 4.2) \\
\addlinespace[2pt]
RF-036 & El sistema deberá permitir parametrizar el período de vigencia de la reserva temporal. & Sí & Servicio de pedidos / Servicio de comisiones / Canal digital (comercio electrónico vigente) & Subdoc.3 §3.3.2 (Frente Retail); Anexo A.4.1 (SUP-11, SUP-12, SUP-13, SUP-14) \textbar{} Subdoc.4 §4.1 (Tabla 4.2) \\
\addlinespace[2pt]
RF-037 & El sistema deberá mantener activa la reserva mientras el cliente mantenga actividad, conforme al período parametrizado. & Sí & Servicio de pedidos / Servicio de comisiones / Canal digital (comercio electrónico vigente) & Subdoc.3 §3.3.2 (Frente Retail); Anexo A.4.1 (SUP-11, SUP-12, SUP-13, SUP-14) \textbar{} Subdoc.4 §4.1 (Tabla 4.2) \\
\addlinespace[2pt]
RF-038 & El sistema deberá impedir que una misma unidad sea comprometida simultáneamente en más de una operación de venta digital. & Sí & Servicio de pedidos / Servicio de comisiones / Canal digital (comercio electrónico vigente) & Subdoc.3 §3.3.2 (Frente Retail); Anexo A.4.1 (SUP-11, SUP-12, SUP-13, SUP-14) \textbar{} Subdoc.4 §4.1 (Tabla 4.2) \\
\addlinespace[2pt]
RF-039 & El sistema debe liberar la reserva temporal al expirar su período de vigencia. & Sí & Servicio de pedidos / Servicio de comisiones / Canal digital (comercio electrónico vigente) & Subdoc.3 §3.3.2 (Frente Retail); Anexo A.4.1 (SUP-11, SUP-12, SUP-13, SUP-14) \textbar{} Subdoc.4 §4.1 (Tabla 4.2) \\
\addlinespace[2pt]
RF-040 & El sistema debe devolver la unidad liberada al disponible para vender. & Sí & Servicio de pedidos / Servicio de comisiones / Canal digital (comercio electrónico vigente) & Subdoc.3 §3.3.2 (Frente Retail); Anexo A.4.1 (SUP-11, SUP-12, SUP-13, SUP-14) \textbar{} Subdoc.4 §4.1 (Tabla 4.2) \\
\addlinespace[2pt]
RF-041 & El sistema deberá liberar de inmediato la reserva digital de esa unidad, sin esperar la expiración del período. & Sí & Servicio de pedidos / Servicio de comisiones / Canal digital (comercio electrónico vigente) & Subdoc.3 §3.3.2 (Frente Retail); Anexo A.4.1 (SUP-11, SUP-12, SUP-13, SUP-14) \textbar{} Subdoc.4 §4.1 (Tabla 4.2) \\
\addlinespace[2pt]
RF-042 & El sistema deberá aplicar el período de vigencia de reserva reducido declarado como parámetro para el evento anual. & Sí & Servicio de pedidos / Servicio de comisiones / Canal digital (comercio electrónico vigente) & Subdoc.3 §3.3.2 (Frente Retail); Anexo A.4.1 (SUP-11, SUP-12, SUP-13, SUP-14) \textbar{} Subdoc.4 §4.1 (Tabla 4.2) \\
\addlinespace[2pt]
RF-043 & El sistema deberá preautorizar el medio de pago al aceptar el pedido, sin capturar el cobro. & Sí & Servicio de pedidos / Servicio de comisiones / Canal digital (comercio electrónico vigente) & Subdoc.3 §3.3.2 (Frente Retail); Anexo A.4.1 (SUP-11, SUP-12, SUP-13, SUP-14) \textbar{} Subdoc.4 §4.1 (Tabla 4.2) \\
\addlinespace[2pt]
RF-044 & El sistema deberá verificar la existencia física real de la unidad en el punto asignado antes de habilitar la captura del cobro. & Sí & Servicio de pedidos / Servicio de comisiones / Canal digital (comercio electrónico vigente) & Subdoc.3 §3.3.2 (Frente Retail); Anexo A.4.1 (SUP-11, SUP-12, SUP-13, SUP-14) \textbar{} Subdoc.4 §4.1 (Tabla 4.2) \\
\addlinespace[2pt]
RF-045 & El sistema deberá capturar el cobro únicamente al registrarse el evento de confirmación de la preparacion física. & Sí & Servicio de pedidos / Servicio de comisiones / Canal digital (comercio electrónico vigente) & Subdoc.3 §3.3.2 (Frente Retail); Anexo A.4.1 (SUP-11, SUP-12, SUP-13, SUP-14) \textbar{} Subdoc.4 §4.1 (Tabla 4.2) \\
\addlinespace[2pt]
RF-046 & El sistema deberá determinar automáticamente la alternativa de resolución aplicable según el motor de reglas (cancelar, sustituir, derivar a tercero o entrega diferida). & Sí & Servicio de pedidos / Servicio de comisiones / Canal digital (comercio electrónico vigente) & Subdoc.3 §3.3.2 (Frente Retail); Anexo A.4.1 (SUP-11, SUP-12, SUP-13, SUP-14) \textbar{} Subdoc.4 §4.1 (Tabla 4.2) \\
\addlinespace[2pt]
RF-047 & El sistema deberá requerir la aceptación explícita del cliente antes de capturar el cobro. & Sí & Servicio de pedidos / Servicio de comisiones / Canal digital (comercio electrónico vigente) & Subdoc.3 §3.3.2 (Frente Retail); Anexo A.4.1 (SUP-11, SUP-12, SUP-13, SUP-14) \textbar{} Subdoc.4 §4.1 (Tabla 4.2) \\
\addlinespace[2pt]
RF-048 & El sistema deberá anular la preautorización de pago del pedido. & Sí & Servicio de pedidos / Servicio de comisiones / Canal digital (comercio electrónico vigente) & Subdoc.3 §3.3.2 (Frente Retail); Anexo A.4.1 (SUP-11, SUP-12, SUP-13, SUP-14) \textbar{} Subdoc.4 §4.1 (Tabla 4.2) \\
\addlinespace[2pt]
RF-049 & El sistema deberá cancelar el pedido registrando el motivo de la cancelación. & Sí & Servicio de pedidos / Servicio de comisiones / Canal digital (comercio electrónico vigente) & Subdoc.3 §3.3.2 (Frente Retail); Anexo A.4.1 (SUP-11, SUP-12, SUP-13, SUP-14) \textbar{} Subdoc.4 §4.1 (Tabla 4.2) \\
\addlinespace[2pt]
RF-050 & El sistema deberá generar la conciliación diaria de preautorizaciones vencidas sin captura. & Sí & Servicio de pedidos / Servicio de comisiones / Canal digital (comercio electrónico vigente) & Subdoc.3 §3.3.2 (Frente Retail); Anexo A.4.1 (SUP-11, SUP-12, SUP-13, SUP-14) \textbar{} Subdoc.4 §4.1 (Tabla 4.2) \\
\addlinespace[2pt]
RF-051 & El sistema deberá notificar al cliente el cambio de estado de su pedido antes de efectuar cualquier cobro definitivo. & Sí & Servicio de pedidos / Servicio de comisiones / Canal digital (comercio electrónico vigente) & Subdoc.3 §3.3.2 (Frente Retail); Anexo A.4.1 (SUP-11, SUP-12, SUP-13, SUP-14) \textbar{} Subdoc.4 §4.1 (Tabla 4.2) \\
\addlinespace[2pt]
RF-052 & El sistema deberá permitir a todos los actores consultar el estado del pedido desde una única fuente de verdad. & Sí & Servicio de pedidos / Servicio de comisiones / Canal digital (comercio electrónico vigente) & Subdoc.3 §3.3.2 (Frente Retail); Anexo A.4.1 (SUP-11, SUP-12, SUP-13, SUP-14) \textbar{} Subdoc.4 §4.1 (Tabla 4.2) \\
\addlinespace[2pt]
RF-053 & El sistema deberá permitir marcar como no elegible para cumplimiento digital el stock de exhibición de las categorías declaradas. & Sí & Servicio de pedidos / Servicio de comisiones / Canal digital (comercio electrónico vigente) & Subdoc.3 §3.3.2 (Frente Retail); Anexo A.4.1 (SUP-11, SUP-12, SUP-13, SUP-14) \textbar{} Subdoc.4 §4.1 (Tabla 4.2) \\
\addlinespace[2pt]
RF-054 & El sistema debe calcular la base de comisión reconociendo al vendedor y a la tienda de origen de la unidad. & Sí & Servicio de pedidos / Servicio de comisiones / Canal digital (comercio electrónico vigente) & Subdoc.3 §3.3.2 (Frente Retail); Anexo A.4.1 (SUP-11, SUP-12, SUP-13, SUP-14) \textbar{} Subdoc.4 §4.1 (Tabla 4.2) \\
\addlinespace[2pt]
RF-055 & El sistema debe condicionar la transmisión de la base de comisión al movimiento real de inventario verificado en bodega. & Sí & Servicio de pedidos / Servicio de comisiones / Canal digital (comercio electrónico vigente) & Subdoc.3 §3.3.2 (Frente Retail); Anexo A.4.1 (SUP-11, SUP-12, SUP-13, SUP-14) \textbar{} Subdoc.4 §4.1 (Tabla 4.2) \\
\addlinespace[2pt]
RF-056 & El sistema deberá reasignar el pedido a otro punto de despacho con existencia disponible. & Sí & Servicio de pedidos / Servicio de comisiones / Canal digital (comercio electrónico vigente) & Subdoc.3 §3.3.2 (Frente Retail); Anexo A.4.1 (SUP-11, SUP-12, SUP-13, SUP-14) \textbar{} Subdoc.4 §4.1 (Tabla 4.2) \\
\addlinespace[2pt]
RF-057 & El sistema deberá impedir una segunda reasignación del mismo pedido. & Sí & Servicio de pedidos / Servicio de comisiones / Canal digital (comercio electrónico vigente) & Subdoc.3 §3.3.2 (Frente Retail); Anexo A.4.1 (SUP-11, SUP-12, SUP-13, SUP-14) \textbar{} Subdoc.4 §4.1 (Tabla 4.2) \\
\addlinespace[2pt]
RF-058 & El sistema deberá acotar la reasignación a la ventana de tiempo parametrizada. & Sí & Servicio de pedidos / Servicio de comisiones / Canal digital (comercio electrónico vigente) & Subdoc.3 §3.3.2 (Frente Retail); Anexo A.4.1 (SUP-11, SUP-12, SUP-13, SUP-14) \textbar{} Subdoc.4 §4.1 (Tabla 4.2) \\
\addlinespace[2pt]
RF-059 & El sistema deberá ofrecer al cliente un producto equivalente. & Sí & Servicio de pedidos / Servicio de comisiones / Canal digital (comercio electrónico vigente) & Subdoc.3 §3.3.2 (Frente Retail); Anexo A.4.1 (SUP-11, SUP-12, SUP-13, SUP-14) \textbar{} Subdoc.4 §4.1 (Tabla 4.2) \\
\addlinespace[2pt]
RF-060 & El sistema deberá ofrecer al cliente la espera compensada con la compensación parametrizada. & Sí & Servicio de pedidos / Servicio de comisiones / Canal digital (comercio electrónico vigente) & Subdoc.3 §3.3.2 (Frente Retail); Anexo A.4.1 (SUP-11, SUP-12, SUP-13, SUP-14) \textbar{} Subdoc.4 §4.1 (Tabla 4.2) \\
\addlinespace[2pt]
RF-061 & El sistema deberá ofrecer al cliente la liberacion del pedido con anulación de la preautorización. & Sí & Servicio de pedidos / Servicio de comisiones / Canal digital (comercio electrónico vigente) & Subdoc.3 §3.3.2 (Frente Retail); Anexo A.4.1 (SUP-11, SUP-12, SUP-13, SUP-14) \textbar{} Subdoc.4 §4.1 (Tabla 4.2) \\
\addlinespace[2pt]
RF-062 & El sistema debe monitorear el tiempo restante de cada pedido respecto de su fecha prometida de entrega. & Sí & Servicio de pedidos / Servicio de comisiones / Canal digital (comercio electrónico vigente) & Subdoc.3 §3.3.2 (Frente Retail); Anexo A.4.1 (SUP-11, SUP-12, SUP-13, SUP-14) \textbar{} Subdoc.4 §4.1 (Tabla 4.2) \\
\addlinespace[2pt]
RF-063 & El sistema debe priorizar la preparación del pedido al alcanzar el umbral parametrizado de tiempo restante. & Sí & Servicio de pedidos / Servicio de comisiones / Canal digital (comercio electrónico vigente) & Subdoc.3 §3.3.2 (Frente Retail); Anexo A.4.1 (SUP-11, SUP-12, SUP-13, SUP-14) \textbar{} Subdoc.4 §4.1 (Tabla 4.2) \\
\addlinespace[2pt]
RF-064 & El sistema debe permitir al cliente no autenticado consultar el precio vigente de una referencia. & Sí & Servicio de pedidos / Servicio de comisiones / Canal digital (comercio electrónico vigente) & Subdoc.3 §3.3.2 (Frente Retail); Anexo A.4.1 (SUP-11, SUP-12, SUP-13, SUP-14) \textbar{} Subdoc.4 §4.1 (Tabla 4.2) \\
\addlinespace[2pt]
RF-065 & El sistema debe permitir al cliente no autenticado consultar la disponibilidad de la referencia por tienda. & Sí & Servicio de pedidos / Servicio de comisiones / Canal digital (comercio electrónico vigente) & Subdoc.3 §3.3.2 (Frente Retail); Anexo A.4.1 (SUP-11, SUP-12, SUP-13, SUP-14) \textbar{} Subdoc.4 §4.1 (Tabla 4.2) \\
\addlinespace[2pt]
RF-066 & El sistema debe permitir al cliente no autenticado consultar la disponibilidad de la referencia para despacho. & Sí & Servicio de pedidos / Servicio de comisiones / Canal digital (comercio electrónico vigente) & Subdoc.3 §3.3.2 (Frente Retail); Anexo A.4.1 (SUP-11, SUP-12, SUP-13, SUP-14) \textbar{} Subdoc.4 §4.1 (Tabla 4.2) \\
\addlinespace[2pt]
RF-067 & El sistema debe permitir al cliente autenticado consultar sus compras. & Sí & Servicio de pedidos / Servicio de comisiones / Canal digital (comercio electrónico vigente) & Subdoc.3 §3.3.2 (Frente Retail); Anexo A.4.1 (SUP-11, SUP-12, SUP-13, SUP-14) \textbar{} Subdoc.4 §4.1 (Tabla 4.2) \\
\addlinespace[2pt]
RF-068 & El sistema debe permitir al cliente autenticado consultar sus devoluciones. & Sí & Servicio de pedidos / Servicio de comisiones / Canal digital (comercio electrónico vigente) & Subdoc.3 §3.3.2 (Frente Retail); Anexo A.4.1 (SUP-11, SUP-12, SUP-13, SUP-14) \textbar{} Subdoc.4 §4.1 (Tabla 4.2) \\
\addlinespace[2pt]
RF-069 & El sistema debe permitir al cliente autenticado consultar su estado de cuenta. & Sí & Servicio de pedidos / Servicio de comisiones / Canal digital (comercio electrónico vigente) & Subdoc.3 §3.3.2 (Frente Retail); Anexo A.4.1 (SUP-11, SUP-12, SUP-13, SUP-14) \textbar{} Subdoc.4 §4.1 (Tabla 4.2) \\
\addlinespace[2pt]
RF-070 & El sistema debe permitir al cliente autenticado consultar sus documentos. & Sí & Servicio de pedidos / Servicio de comisiones / Canal digital (comercio electrónico vigente) & Subdoc.3 §3.3.2 (Frente Retail); Anexo A.4.1 (SUP-11, SUP-12, SUP-13, SUP-14) \textbar{} Subdoc.4 §4.1 (Tabla 4.2) \\
\addlinespace[2pt]
RF-071 & El sistema deberá calcular la base de comisión considerando el canal de origen y el canal de cumplimiento cuando estos difieran. & Sí & Servicio de pedidos / Servicio de comisiones / Canal digital (comercio electrónico vigente) & Subdoc.3 §3.3.2 (Frente Retail); Anexo A.4.1 (SUP-11, SUP-12, SUP-13, SUP-14) \textbar{} Subdoc.4 §4.1 (Tabla 4.2) \\
\addlinespace[2pt]
RF-072 & El sistema debe permitir al ejecutivo de mesón consultar el estado real del pedido. & Sí & Servicio de pedidos / Servicio de comisiones / Canal digital (comercio electrónico vigente) & Subdoc.3 §3.3.2 (Frente Retail); Anexo A.4.1 (SUP-11, SUP-12, SUP-13, SUP-14) \textbar{} Subdoc.4 §4.1 (Tabla 4.2) \\
\addlinespace[2pt]
RF-073 & El sistema debe permitir al ejecutivo de mesón consultar el precio aplicado al pedido. & Sí & Servicio de pedidos / Servicio de comisiones / Canal digital (comercio electrónico vigente) & Subdoc.3 §3.3.2 (Frente Retail); Anexo A.4.1 (SUP-11, SUP-12, SUP-13, SUP-14) \textbar{} Subdoc.4 §4.1 (Tabla 4.2) \\
\addlinespace[2pt]
RF-074 & El sistema debe permitir al ejecutivo de mesón ofrecer al cliente las opciones de resolución definidas en RF-059, RF-060 y RF-061. & Sí & Servicio de pedidos / Servicio de comisiones / Canal digital (comercio electrónico vigente) & Subdoc.3 §3.3.2 (Frente Retail); Anexo A.4.1 (SUP-11, SUP-12, SUP-13, SUP-14) \textbar{} Subdoc.4 §4.1 (Tabla 4.2) \\
\addlinespace[2pt]
RF-075 & El sistema deberá excluir al centro de distribución de Concepción como origen de promesa de entrega en el canal digital mientras no este acreditado su sistema de gestión de almacenes. & Sí & Servicio de pedidos / Servicio de comisiones / Canal digital (comercio electrónico vigente) & Subdoc.3 §3.3.2 (Frente Retail); Anexo A.4.1 (SUP-11, SUP-12, SUP-13, SUP-14) \textbar{} Subdoc.4 §4.1 (Tabla 4.2) \\
\addlinespace[2pt]
RF-076 & El sistema deberá calcular el costo total de servir de cada punto de despacho candidato, incorporando costo logístico, costo de oportunidad de sala y plazo comprometido. & Sí & Servicio de pedidos / Servicio de comisiones / Canal digital (comercio electrónico vigente) & Subdoc.3 §3.3.2 (Frente Retail); Anexo A.4.1 (SUP-11, SUP-12, SUP-13, SUP-14) \textbar{} Subdoc.4 §4.1 (Tabla 4.2) \\
\addlinespace[2pt]
RF-077 & El sistema deberá seleccionar como punto de despacho aquel de menor costo total de servir. & Sí & Servicio de pedidos / Servicio de comisiones / Canal digital (comercio electrónico vigente) & Subdoc.3 §3.3.2 (Frente Retail); Anexo A.4.1 (SUP-11, SUP-12, SUP-13, SUP-14) \textbar{} Subdoc.4 §4.1 (Tabla 4.2) \\
\addlinespace[2pt]
RF-078 & El sistema deberá registrar el valor de cada término del costo total de servir que fundamento la seleccion. & Sí & Servicio de pedidos / Servicio de comisiones / Canal digital (comercio electrónico vigente) & Subdoc.3 §3.3.2 (Frente Retail); Anexo A.4.1 (SUP-11, SUP-12, SUP-13, SUP-14) \textbar{} Subdoc.4 §4.1 (Tabla 4.2) \\
\addlinespace[2pt]
RF-079 & El sistema deberá calcular la fecha prometida de entrega en función del punto de despacho, la capacidad de preparacion y el transportista asignado. & Sí & Servicio de pedidos / Servicio de comisiones / Canal digital (comercio electrónico vigente) & Subdoc.3 §3.3.2 (Frente Retail); Anexo A.4.1 (SUP-11, SUP-12, SUP-13, SUP-14) \textbar{} Subdoc.4 §4.1 (Tabla 4.2) \\
\addlinespace[2pt]
RF-080 & El sistema deberá impedir la publicación de un plazo fijo de catálogo como fecha prometida de entrega. & Sí & Servicio de pedidos / Servicio de comisiones / Canal digital (comercio electrónico vigente) & Subdoc.3 §3.3.2 (Frente Retail); Anexo A.4.1 (SUP-11, SUP-12, SUP-13, SUP-14) \textbar{} Subdoc.4 §4.1 (Tabla 4.2) \\
\addlinespace[2pt]
RF-081 & El sistema deberá calcular el cumplimiento de la promesa de entrega por pedido individual. & Sí & Servicio de pedidos / Servicio de comisiones / Canal digital (comercio electrónico vigente) & Subdoc.3 §3.3.2 (Frente Retail); Anexo A.4.1 (SUP-11, SUP-12, SUP-13, SUP-14) \textbar{} Subdoc.4 §4.1 (Tabla 4.2) \\
\addlinespace[2pt]
RF-082 & El sistema deberá detectar la pérdida del enlace externo y conmutar automáticamente la tienda a modo desconectado. & Sí & Punto de venta con operación sin conexión / Servicio de ventas / Servicio de originación de crédito & Subdoc.3 §3.3.2 (Frente Retail, Frente Filial emisora); Anexo A.4.1 (SUP-07) \textbar{} Subdoc.4 §4.2 \\
\addlinespace[2pt]
RF-083 & El sistema deberá permitir al cajero registrar una venta en modo desconectado. & Sí & Punto de venta con operación sin conexión / Servicio de ventas / Servicio de originación de crédito & Subdoc.3 §3.3.2 (Frente Retail, Frente Filial emisora); Anexo A.4.1 (SUP-07) \textbar{} Subdoc.4 §4.2 \\
\addlinespace[2pt]
RF-084 & El sistema deberá permitir al cajero cobrar la venta en modo desconectado. & Sí & Punto de venta con operación sin conexión / Servicio de ventas / Servicio de originación de crédito & Subdoc.3 §3.3.2 (Frente Retail, Frente Filial emisora); Anexo A.4.1 (SUP-07) \textbar{} Subdoc.4 §4.2 \\
\addlinespace[2pt]
RF-085 & El sistema deberá aplicar las promociones vigentes en modo desconectado, evaluadas según RF-032, RF-033 y RF-034. & Sí & Punto de venta con operación sin conexión / Servicio de ventas / Servicio de originación de crédito & Subdoc.3 §3.3.2 (Frente Retail, Frente Filial emisora); Anexo A.4.1 (SUP-07) \textbar{} Subdoc.4 §4.2 \\
\addlinespace[2pt]
RF-086 & El sistema deberá emitir el documento de venta en contingencia utilizando folios previamente asignados por el sistema de gestión empresarial. & Sí & Punto de venta con operación sin conexión / Servicio de ventas / Servicio de originación de crédito & Subdoc.3 §3.3.2 (Frente Retail, Frente Filial emisora); Anexo A.4.1 (SUP-07) \textbar{} Subdoc.4 §4.2 \\
\addlinespace[2pt]
RF-087 & El sistema deberá permitir al vendedor consultar la existencia local de la tienda en modo desconectado. & Sí & Punto de venta con operación sin conexión / Servicio de ventas / Servicio de originación de crédito & Subdoc.3 §3.3.2 (Frente Retail, Frente Filial emisora); Anexo A.4.1 (SUP-07) \textbar{} Subdoc.4 §4.2 \\
\addlinespace[2pt]
RF-088 & El sistema deberá detectar el restablecimiento del enlace y conmutar la tienda a modo conectado. & Sí & Punto de venta con operación sin conexión / Servicio de ventas / Servicio de originación de crédito & Subdoc.3 §3.3.2 (Frente Retail, Frente Filial emisora); Anexo A.4.1 (SUP-07) \textbar{} Subdoc.4 §4.2 \\
\addlinespace[2pt]
RF-089 & El sistema deberá permitir el otorgamiento de crédito en modo desconectado exclusivamente contra cupo preaprobado vigente almacenado en caché local. & Sí & Punto de venta con operación sin conexión / Servicio de ventas / Servicio de originación de crédito & Subdoc.3 §3.3.2 (Frente Retail, Frente Filial emisora); Anexo A.4.1 (SUP-07) \textbar{} Subdoc.4 §4.2 \\
\addlinespace[2pt]
RF-090 & El sistema deberá aplicar el tope de monto parametrizado a cada operación de crédito cursada en modo desconectado. & Sí & Punto de venta con operación sin conexión / Servicio de ventas / Servicio de originación de crédito & Subdoc.3 §3.3.2 (Frente Retail, Frente Filial emisora); Anexo A.4.1 (SUP-07) \textbar{} Subdoc.4 §4.2 \\
\addlinespace[2pt]
RF-091 & El sistema deberá aplicar el tope parametrizado de número de operaciones de crédito en modo desconectado. & Sí & Punto de venta con operación sin conexión / Servicio de ventas / Servicio de originación de crédito & Subdoc.3 §3.3.2 (Frente Retail, Frente Filial emisora); Anexo A.4.1 (SUP-07) \textbar{} Subdoc.4 §4.2 \\
\addlinespace[2pt]
RF-092 & El sistema deberá impedir la apertura de una tarjeta nueva en modo desconectado. & Sí & Punto de venta con operación sin conexión / Servicio de ventas / Servicio de originación de crédito & Subdoc.3 §3.3.2 (Frente Retail, Frente Filial emisora); Anexo A.4.1 (SUP-07) \textbar{} Subdoc.4 §4.2 \\
\addlinespace[2pt]
RF-093 & El sistema deberá impedir la ampliación de cupo en modo desconectado. & Sí & Punto de venta con operación sin conexión / Servicio de ventas / Servicio de originación de crédito & Subdoc.3 §3.3.2 (Frente Retail, Frente Filial emisora); Anexo A.4.1 (SUP-07) \textbar{} Subdoc.4 §4.2 \\
\addlinespace[2pt]
RF-094 & El sistema deberá marcar la operación como cursada en modo desconectado para su validación posterior. & Sí & Punto de venta con operación sin conexión / Servicio de ventas / Servicio de originación de crédito & Subdoc.3 §3.3.2 (Frente Retail, Frente Filial emisora); Anexo A.4.1 (SUP-07) \textbar{} Subdoc.4 §4.2 \\
\addlinespace[2pt]
RF-095 & El sistema debe reconciliar hacia los sistemas centrales la totalidad de las ventas registradas en modo desconectado. & Sí & Punto de venta con operación sin conexión / Servicio de ventas / Servicio de originación de crédito & Subdoc.3 §3.3.2 (Frente Retail, Frente Filial emisora); Anexo A.4.1 (SUP-07) \textbar{} Subdoc.4 §4.2 \\
\addlinespace[2pt]
RF-096 & El sistema debe reconciliar hacia los sistemas centrales la totalidad de los documentos emitidos en modo desconectado. & Sí & Punto de venta con operación sin conexión / Servicio de ventas / Servicio de originación de crédito & Subdoc.3 §3.3.2 (Frente Retail, Frente Filial emisora); Anexo A.4.1 (SUP-07) \textbar{} Subdoc.4 §4.2 \\
\addlinespace[2pt]
RF-097 & El sistema deberá procesar la reconciliación de forma idempotente, impidiendo la duplicación de ventas o documentos ante reintentos. & Sí & Punto de venta con operación sin conexión / Servicio de ventas / Servicio de originación de crédito & Subdoc.3 §3.3.2 (Frente Retail, Frente Filial emisora); Anexo A.4.1 (SUP-07) \textbar{} Subdoc.4 §4.2 \\
\addlinespace[2pt]
RF-098 & El sistema deberá resolver el conflicto de existencia comprometida aplicando la regla de precedencia declarada, produciendo el mismo resultado ante repetición. & Sí & Punto de venta con operación sin conexión / Servicio de ventas / Servicio de originación de crédito & Subdoc.3 §3.3.2 (Frente Retail, Frente Filial emisora); Anexo A.4.1 (SUP-07) \textbar{} Subdoc.4 §4.2 \\
\addlinespace[2pt]
RF-099 & El sistema deberá generar un informe de excepciones de la reconciliación, identificando cada conflicto resuelto y su regla aplicada. & Sí & Punto de venta con operación sin conexión / Servicio de ventas / Servicio de originación de crédito & Subdoc.3 §3.3.2 (Frente Retail, Frente Filial emisora); Anexo A.4.1 (SUP-07) \textbar{} Subdoc.4 §4.2 \\
\addlinespace[2pt]
RF-100 & El sistema deberá enrutar la emisión de todo documento tributario hacia el sistema de gestión empresarial como único emisor. & Sí & Punto de venta con operación sin conexión / Servicio de ventas / Servicio de originación de crédito & Subdoc.3 §3.3.2 (Frente Retail, Frente Filial emisora); Anexo A.4.1 (SUP-07) \textbar{} Subdoc.4 §4.2 \\
\addlinespace[2pt]
RF-101 & El sistema deberá permitir al vendedor de piso consultar la disponibilidad para vender de una referencia. & Sí & Punto de venta con operación sin conexión / Servicio de ventas / Servicio de originación de crédito & Subdoc.3 §3.3.2 (Frente Retail, Frente Filial emisora); Anexo A.4.1 (SUP-07) \textbar{} Subdoc.4 §4.2 \\
\addlinespace[2pt]
RF-102 & El sistema debe permitir al vendedor de marketplace declarar la existencia de sus referencias. & Sí & Servicio de marketplace & Subdoc.3 §3.3.2 (Frente Retail); Anexo A.4.1 (SUP-15, SUP-16) \textbar{} Subdoc.4 §4.1 (Tabla 4.2), §4.2 \\
\addlinespace[2pt]
RF-103 & El sistema debe permitir al vendedor de marketplace actualizar la existencia previamente declarada. & Sí & Servicio de marketplace & Subdoc.3 §3.3.2 (Frente Retail); Anexo A.4.1 (SUP-15, SUP-16) \textbar{} Subdoc.4 §4.1 (Tabla 4.2), §4.2 \\
\addlinespace[2pt]
RF-104 & El sistema deberá registrar cada actualización de existencia declarada con su fecha y hora. & Sí & Servicio de marketplace & Subdoc.3 §3.3.2 (Frente Retail); Anexo A.4.1 (SUP-15, SUP-16) \textbar{} Subdoc.4 §4.1 (Tabla 4.2), §4.2 \\
\addlinespace[2pt]
RF-105 & El sistema deberá calcular los indicadores de nivel de servicio por vendedor externo según las reglas publicadas. & Sí & Servicio de marketplace & Subdoc.3 §3.3.2 (Frente Retail); Anexo A.4.1 (SUP-15, SUP-16) \textbar{} Subdoc.4 §4.1 (Tabla 4.2), §4.2 \\
\addlinespace[2pt]
RF-106 & El sistema deberá notificar al vendedor de marketplace la recepción de la devolución en el momento en que se registra. & Sí & Servicio de marketplace & Subdoc.3 §3.3.2 (Frente Retail); Anexo A.4.1 (SUP-15, SUP-16) \textbar{} Subdoc.4 §4.1 (Tabla 4.2), §4.2 \\
\addlinespace[2pt]
RF-107 & El sistema deberá permitir al ejecutivo registrar que parte asume el costo de la devolución. & Sí & Servicio de marketplace & Subdoc.3 §3.3.2 (Frente Retail); Anexo A.4.1 (SUP-15, SUP-16) \textbar{} Subdoc.4 §4.1 (Tabla 4.2), §4.2 \\
\addlinespace[2pt]
RF-108 & El sistema deberá permitir al ejecutivo registrar la prestación entregada al cliente (reparación, reposición o devolución del precio). & Sí & Servicio de marketplace & Subdoc.3 §3.3.2 (Frente Retail); Anexo A.4.1 (SUP-15, SUP-16) \textbar{} Subdoc.4 §4.1 (Tabla 4.2), §4.2 \\
\addlinespace[2pt]
RF-109 & El sistema deberá publicar únicamente la existencia declarada en la última actualización vigente del vendedor. & Sí & Servicio de marketplace & Subdoc.3 §3.3.2 (Frente Retail); Anexo A.4.1 (SUP-15, SUP-16) \textbar{} Subdoc.4 §4.1 (Tabla 4.2), §4.2 \\
\addlinespace[2pt]
RF-110 & El sistema debe permitir al vendedor de marketplace consultar el estado de cada pedido intermediado. & Sí & Servicio de marketplace & Subdoc.3 §3.3.2 (Frente Retail); Anexo A.4.1 (SUP-15, SUP-16) \textbar{} Subdoc.4 §4.1 (Tabla 4.2), §4.2 \\
\addlinespace[2pt]
RF-111 & El sistema debe permitir al vendedor de marketplace consultar las devoluciones de sus pedidos. & Sí & Servicio de marketplace & Subdoc.3 §3.3.2 (Frente Retail); Anexo A.4.1 (SUP-15, SUP-16) \textbar{} Subdoc.4 §4.1 (Tabla 4.2), §4.2 \\
\addlinespace[2pt]
RF-112 & El sistema debe permitir al vendedor de marketplace consultar su evaluación de desempeño vigente. & Sí & Servicio de marketplace & Subdoc.3 §3.3.2 (Frente Retail); Anexo A.4.1 (SUP-15, SUP-16) \textbar{} Subdoc.4 §4.1 (Tabla 4.2), §4.2 \\
\addlinespace[2pt]
RF-113 & El sistema debe identificar visiblemente al vendedor de la unidad en el catálogo, la ficha de producto y el proceso de compra. & Sí & Servicio de marketplace & Subdoc.3 §3.3.2 (Frente Retail); Anexo A.4.1 (SUP-15, SUP-16) \textbar{} Subdoc.4 §4.1 (Tabla 4.2), §4.2 \\
\addlinespace[2pt]
RF-114 & El sistema debe identificar visiblemente a quien despacha la unidad en el catálogo, la ficha de producto y el proceso de compra. & Sí & Servicio de marketplace & Subdoc.3 §3.3.2 (Frente Retail); Anexo A.4.1 (SUP-15, SUP-16) \textbar{} Subdoc.4 §4.1 (Tabla 4.2), §4.2 \\
\addlinespace[2pt]
RF-115 & El sistema debe mostrar las condiciones de devolución aplicables a la unidad. & Sí & Servicio de marketplace & Subdoc.3 §3.3.2 (Frente Retail); Anexo A.4.1 (SUP-15, SUP-16) \textbar{} Subdoc.4 §4.1 (Tabla 4.2), §4.2 \\
\addlinespace[2pt]
RF-116 & El sistema debe mostrar las condiciones de garantía legal aplicables a la unidad. & Sí & Servicio de marketplace & Subdoc.3 §3.3.2 (Frente Retail); Anexo A.4.1 (SUP-15, SUP-16) \textbar{} Subdoc.4 §4.1 (Tabla 4.2), §4.2 \\
\addlinespace[2pt]
RF-117 & El sistema deberá impedir que un pedido intermediado comprometa existencia propia de la compañía. & Sí & Servicio de marketplace & Subdoc.3 §3.3.2 (Frente Retail); Anexo A.4.1 (SUP-15, SUP-16) \textbar{} Subdoc.4 §4.1 (Tabla 4.2), §4.2 \\
\addlinespace[2pt]
RF-118 & El sistema deberá impedir que una unidad de inventario propio sea asignada al cumplimiento de un pedido de marketplace. & Sí & Servicio de marketplace & Subdoc.3 §3.3.2 (Frente Retail); Anexo A.4.1 (SUP-15, SUP-16) \textbar{} Subdoc.4 §4.1 (Tabla 4.2), §4.2 \\
\addlinespace[2pt]
RF-119 & El sistema deberá registrar el acuse de conocimiento de las reglas de evaluación por parte de cada vendedor externo. & Sí & Servicio de marketplace & Subdoc.3 §3.3.2 (Frente Retail); Anexo A.4.1 (SUP-15, SUP-16) \textbar{} Subdoc.4 §4.1 (Tabla 4.2), §4.2 \\
\addlinespace[2pt]
RF-120 & El sistema deberá impedir la aplicación de una regla de evaluación a un vendedor que no tenga acuse de conocimiento previo de esa regla. & Sí & Servicio de marketplace & Subdoc.3 §3.3.2 (Frente Retail); Anexo A.4.1 (SUP-15, SUP-16) \textbar{} Subdoc.4 §4.1 (Tabla 4.2), §4.2 \\
\addlinespace[2pt]
RF-121 & El sistema deberá determinar la consecuencia escalonada aplicable según la matriz de escalamiento parametrizada (advertencia, restricción de publicación, suspensión). & Sí & Servicio de marketplace & Subdoc.3 §3.3.2 (Frente Retail); Anexo A.4.1 (SUP-15, SUP-16) \textbar{} Subdoc.4 §4.1 (Tabla 4.2), §4.2 \\
\addlinespace[2pt]
RF-122 & El sistema deberá requerir la ejecución de la sanción por el rol nominado en la matriz de escalamiento. & Sí & Servicio de marketplace & Subdoc.3 §3.3.2 (Frente Retail); Anexo A.4.1 (SUP-15, SUP-16) \textbar{} Subdoc.4 §4.1 (Tabla 4.2), §4.2 \\
\addlinespace[2pt]
RF-123 & El sistema deberá registrar la trazabilidad de la sanción aplicada, indicando regla incumplida, consecuencia, rol ejecutor e instante. & Sí & Servicio de marketplace & Subdoc.3 §3.3.2 (Frente Retail); Anexo A.4.1 (SUP-15, SUP-16) \textbar{} Subdoc.4 §4.1 (Tabla 4.2), §4.2 \\
\addlinespace[2pt]
RF-124 & El sistema deberá despublicar automáticamente la oferta cuyo stock declarado haya superado el plazo de vigencia sin actualización. & Sí & Servicio de marketplace & Subdoc.3 §3.3.2 (Frente Retail); Anexo A.4.1 (SUP-15, SUP-16) \textbar{} Subdoc.4 §4.1 (Tabla 4.2), §4.2 \\
\addlinespace[2pt]
RF-125 & El sistema deberá calcular e informar al sistema de gestión empresarial la base de comisión de marketplace sobre la venta efectivamente cumplida & Sí & Servicio de marketplace & Subdoc.3 §3.3.2 (Frente Retail); Anexo A.4.1 (SUP-15, SUP-16) \textbar{} Subdoc.4 §4.1 (Tabla 4.2), §4.2 \\
\addlinespace[2pt]
RF-126 & El sistema deberá notificar al sistema de gestión empresarial la anulación de la base de comisión calculada cuando ocurra una devolución o cancelación. & Sí & Servicio de marketplace & Subdoc.3 §3.3.2 (Frente Retail); Anexo A.4.1 (SUP-15, SUP-16) \textbar{} Subdoc.4 §4.1 (Tabla 4.2), §4.2 \\
\addlinespace[2pt]
RF-127 & El sistema deberá calcular la existencia disponible para vender restando de la existencia registrada las reservas vigentes, el comprometido no despachado y el colchón de confianza aplicable. & Sí & Servicio de existencias / Servicio de oferta comercial / Servicio de abastecimiento & Subdoc.3 §3.3.2 (Frente Retail); Anexo A.4.1 (SUP-01, SUP-17, SUP-18, SUP-20) \textbar{} Subdoc.4 §4.2 \\
\addlinespace[2pt]
RF-128 & El sistema deberá determinar el valor del colchón de confianza aplicable en función de la categoría del artículo. & Sí & Servicio de existencias / Servicio de oferta comercial / Servicio de abastecimiento & Subdoc.3 §3.3.2 (Frente Retail); Anexo A.4.1 (SUP-01, SUP-17, SUP-18, SUP-20) \textbar{} Subdoc.4 §4.2 \\
\addlinespace[2pt]
RF-129 & El sistema deberá determinar el valor del colchón de confianza aplicable en función del punto de existencia. & Sí & Servicio de existencias / Servicio de oferta comercial / Servicio de abastecimiento & Subdoc.3 §3.3.2 (Frente Retail); Anexo A.4.1 (SUP-01, SUP-17, SUP-18, SUP-20) \textbar{} Subdoc.4 §4.2 \\
\addlinespace[2pt]
RF-130 & El sistema deberá registrar la traza del cálculo del disponible, incluyendo el valor de cada término de la fórmula y la versión de parámetros aplicada. & Sí & Servicio de existencias / Servicio de oferta comercial / Servicio de abastecimiento & Subdoc.3 §3.3.2 (Frente Retail); Anexo A.4.1 (SUP-01, SUP-17, SUP-18, SUP-20) \textbar{} Subdoc.4 §4.2 \\
\addlinespace[2pt]
RF-131 & El sistema deberá permitir parametrizar la frecuencia de conteo cíclico por categoría. & Sí & Servicio de existencias / Servicio de oferta comercial / Servicio de abastecimiento & Subdoc.3 §3.3.2 (Frente Retail); Anexo A.4.1 (SUP-01, SUP-17, SUP-18, SUP-20) \textbar{} Subdoc.4 §4.2 \\
\addlinespace[2pt]
RF-132 & El sistema deberá permitir parametrizar el método de conteo por categoría. & Sí & Servicio de existencias / Servicio de oferta comercial / Servicio de abastecimiento & Subdoc.3 §3.3.2 (Frente Retail); Anexo A.4.1 (SUP-01, SUP-17, SUP-18, SUP-20) \textbar{} Subdoc.4 §4.2 \\
\addlinespace[2pt]
RF-133 & El sistema deberá permitir parametrizar el criterio de gatillo de recuento extraordinario por categoría. & Sí & Servicio de existencias / Servicio de oferta comercial / Servicio de abastecimiento & Subdoc.3 §3.3.2 (Frente Retail); Anexo A.4.1 (SUP-01, SUP-17, SUP-18, SUP-20) \textbar{} Subdoc.4 §4.2 \\
\addlinespace[2pt]
RF-134 & El sistema deberá generar la programación de conteos cíclicos según la frecuencia parametrizada de cada categoría. & Sí & Servicio de existencias / Servicio de oferta comercial / Servicio de abastecimiento & Subdoc.3 §3.3.2 (Frente Retail); Anexo A.4.1 (SUP-01, SUP-17, SUP-18, SUP-20) \textbar{} Subdoc.4 §4.2 \\
\addlinespace[2pt]
RF-135 & El sistema debe permitir registrar el resultado del conteo cíclico. & Sí & Servicio de existencias / Servicio de oferta comercial / Servicio de abastecimiento & Subdoc.3 §3.3.2 (Frente Retail); Anexo A.4.1 (SUP-01, SUP-17, SUP-18, SUP-20) \textbar{} Subdoc.4 §4.2 \\
\addlinespace[2pt]
RF-136 & El sistema deberá señalar automáticamente los pasillos o referencias que requieren recuento extraordinario. & Sí & Servicio de existencias / Servicio de oferta comercial / Servicio de abastecimiento & Subdoc.3 §3.3.2 (Frente Retail); Anexo A.4.1 (SUP-01, SUP-17, SUP-18, SUP-20) \textbar{} Subdoc.4 §4.2 \\
\addlinespace[2pt]
RF-137 & El sistema deberá calcular la exactitud de inventario resultante por categoría. & Sí & Servicio de existencias / Servicio de oferta comercial / Servicio de abastecimiento & Subdoc.3 §3.3.2 (Frente Retail); Anexo A.4.1 (SUP-01, SUP-17, SUP-18, SUP-20) \textbar{} Subdoc.4 §4.2 \\
\addlinespace[2pt]
RF-138 & El sistema deberá permitir clasificar cada diferencia detectada en uno de los componentes de merma definidos (pérdida física, daño no dado de baja, error de recepción, unidad mal ubicada, devolución mal reintegrada, error de digitación). & Sí & Servicio de existencias / Servicio de oferta comercial / Servicio de abastecimiento & Subdoc.3 §3.3.2 (Frente Retail); Anexo A.4.1 (SUP-01, SUP-17, SUP-18, SUP-20) \textbar{} Subdoc.4 §4.2 \\
\addlinespace[2pt]
RF-139 & El sistema deberá impedir el cierre de un ajuste de inventario que no tenga asignado un componente de merma. & Sí & Servicio de existencias / Servicio de oferta comercial / Servicio de abastecimiento & Subdoc.3 §3.3.2 (Frente Retail); Anexo A.4.1 (SUP-01, SUP-17, SUP-18, SUP-20) \textbar{} Subdoc.4 §4.2 \\
\addlinespace[2pt]
RF-140 & El sistema deberá cuantificar el monto y la proporción de cada componente de merma sobre el total del período. & Sí & Servicio de existencias / Servicio de oferta comercial / Servicio de abastecimiento & Subdoc.3 §3.3.2 (Frente Retail); Anexo A.4.1 (SUP-01, SUP-17, SUP-18, SUP-20) \textbar{} Subdoc.4 §4.2 \\
\addlinespace[2pt]
RF-141 & El sistema deberá generar el informe mensual de merma cuya suma de componentes sea igual al total registrado en el período. & Sí & Servicio de existencias / Servicio de oferta comercial / Servicio de abastecimiento & Subdoc.3 §3.3.2 (Frente Retail); Anexo A.4.1 (SUP-01, SUP-17, SUP-18, SUP-20) \textbar{} Subdoc.4 §4.2 \\
\addlinespace[2pt]
RF-142 & El sistema deberá exigir al analista comercial completar los atributos obligatorios de la referencia (categoría, dimensiones, fragilidad, volumen, restricción de despacho) antes de guardar el registro como apto para uso operativo. & Sí & Servicio de existencias / Servicio de oferta comercial / Servicio de abastecimiento & Subdoc.3 §3.3.2 (Frente Retail); Anexo A.4.1 (SUP-01, SUP-17, SUP-18, SUP-20) \textbar{} Subdoc.4 §4.2 \\
\addlinespace[2pt]
RF-143 & El sistema deberá impedir la publicación de esa referencia en el canal digital mientras no cuente con los atributos obligatorios completos. & Sí & Servicio de existencias / Servicio de oferta comercial / Servicio de abastecimiento & Subdoc.3 §3.3.2 (Frente Retail); Anexo A.4.1 (SUP-01, SUP-17, SUP-18, SUP-20) \textbar{} Subdoc.4 §4.2 \\
\addlinespace[2pt]
RF-144 & El sistema deberá generar automáticamente la propuesta diaria de reposición utilizando el disponible para vender calculado en RF-127. & Sí & Servicio de existencias / Servicio de oferta comercial / Servicio de abastecimiento & Subdoc.3 §3.3.2 (Frente Retail); Anexo A.4.1 (SUP-01, SUP-17, SUP-18, SUP-20) \textbar{} Subdoc.4 §4.2 \\
\addlinespace[2pt]
RF-145 & El sistema deberá permitir al planificador de logística ajustar la propuesta de reposición antes de confirmar el envío. & Sí & Servicio de existencias / Servicio de oferta comercial / Servicio de abastecimiento & Subdoc.3 §3.3.2 (Frente Retail); Anexo A.4.1 (SUP-01, SUP-17, SUP-18, SUP-20) \textbar{} Subdoc.4 §4.2 \\
\addlinespace[2pt]
RF-146 & El sistema deberá registrar el ajuste realizado indicando usuario, valor propuesto, valor confirmado e instante. & Sí & Servicio de existencias / Servicio de oferta comercial / Servicio de abastecimiento & Subdoc.3 §3.3.2 (Frente Retail); Anexo A.4.1 (SUP-01, SUP-17, SUP-18, SUP-20) \textbar{} Subdoc.4 §4.2 \\
\addlinespace[2pt]
RF-147 & El sistema debe generar un reporte periódico de referencias con atributos obligatorios incompletos. & Sí & Servicio de existencias / Servicio de oferta comercial / Servicio de abastecimiento & Subdoc.3 §3.3.2 (Frente Retail); Anexo A.4.1 (SUP-01, SUP-17, SUP-18, SUP-20) \textbar{} Subdoc.4 §4.2 \\
\addlinespace[2pt]
RF-148 & El sistema debe generar un reporte periódico de publicaciones fallidas indicando la causa del error. & Sí & Servicio de existencias / Servicio de oferta comercial / Servicio de abastecimiento & Subdoc.3 §3.3.2 (Frente Retail); Anexo A.4.1 (SUP-01, SUP-17, SUP-18, SUP-20) \textbar{} Subdoc.4 §4.2 \\
\addlinespace[2pt]
RF-149 & El sistema debe permitir al gerente de logística definir el valor del colchón de confianza por categoría. & Sí & Servicio de existencias / Servicio de oferta comercial / Servicio de abastecimiento & Subdoc.3 §3.3.2 (Frente Retail); Anexo A.4.1 (SUP-01, SUP-17, SUP-18, SUP-20) \textbar{} Subdoc.4 §4.2 \\
\addlinespace[2pt]
RF-150 & El sistema debe permitir al gerente de logística definir el valor del colchón de confianza por punto de existencia. & Sí & Servicio de existencias / Servicio de oferta comercial / Servicio de abastecimiento & Subdoc.3 §3.3.2 (Frente Retail); Anexo A.4.1 (SUP-01, SUP-17, SUP-18, SUP-20) \textbar{} Subdoc.4 §4.2 \\
\addlinespace[2pt]
RF-151 & El sistema deberá registrar quien definio o modifico el colchón, cuando, el valor anterior y el valor nuevo. & Sí & Servicio de existencias / Servicio de oferta comercial / Servicio de abastecimiento & Subdoc.3 §3.3.2 (Frente Retail); Anexo A.4.1 (SUP-01, SUP-17, SUP-18, SUP-20) \textbar{} Subdoc.4 §4.2 \\
\addlinespace[2pt]
RF-152 & El sistema deberá sugerir un valor de referencia del colchón calculado a partir del histórico de ventas y quiebres de la categoría. & Sí & Servicio de existencias / Servicio de oferta comercial / Servicio de abastecimiento & Subdoc.3 §3.3.2 (Frente Retail); Anexo A.4.1 (SUP-01, SUP-17, SUP-18, SUP-20) \textbar{} Subdoc.4 §4.2 \\
\addlinespace[2pt]
RF-153 & El sistema deberá permitir al vendedor de piso registrar las unidades que ingresan al probador. & Sí & Servicio de existencias / Servicio de oferta comercial / Servicio de abastecimiento & Subdoc.3 §3.3.2 (Frente Retail); Anexo A.4.1 (SUP-01, SUP-17, SUP-18, SUP-20) \textbar{} Subdoc.4 §4.2 \\
\addlinespace[2pt]
RF-154 & El sistema deberá excluir del disponible para vender las unidades registradas en probador. & Sí & Servicio de existencias / Servicio de oferta comercial / Servicio de abastecimiento & Subdoc.3 §3.3.2 (Frente Retail); Anexo A.4.1 (SUP-01, SUP-17, SUP-18, SUP-20) \textbar{} Subdoc.4 §4.2 \\
\addlinespace[2pt]
RF-155 & El sistema debe permitir al vendedor de piso confirmar el reingreso a sala de la unidad registrada en probador. & Sí & Servicio de existencias / Servicio de oferta comercial / Servicio de abastecimiento & Subdoc.3 §3.3.2 (Frente Retail); Anexo A.4.1 (SUP-01, SUP-17, SUP-18, SUP-20) \textbar{} Subdoc.4 §4.2 \\
\addlinespace[2pt]
RF-156 & El sistema debe permitir al vendedor de piso registrar la venta de la unidad registrada en probador. & Sí & Servicio de existencias / Servicio de oferta comercial / Servicio de abastecimiento & Subdoc.3 §3.3.2 (Frente Retail); Anexo A.4.1 (SUP-01, SUP-17, SUP-18, SUP-20) \textbar{} Subdoc.4 §4.2 \\
\addlinespace[2pt]
RF-157 & El sistema deberá impedir que cualquier canal de venta consuma el saldo bruto de inventario para publicar o comprometer existencia. & Sí & Servicio de existencias / Servicio de oferta comercial / Servicio de abastecimiento & Subdoc.3 §3.3.2 (Frente Retail); Anexo A.4.1 (SUP-01, SUP-17, SUP-18, SUP-20) \textbar{} Subdoc.4 §4.2 \\
\addlinespace[2pt]
RF-158 & El sistema debe calcular el porcentaje de error probable de la referencia en función de su categoría. & Sí & Servicio de existencias / Servicio de oferta comercial / Servicio de abastecimiento & Subdoc.3 §3.3.2 (Frente Retail); Anexo A.4.1 (SUP-01, SUP-17, SUP-18, SUP-20) \textbar{} Subdoc.4 §4.2 \\
\addlinespace[2pt]
RF-159 & El sistema debe calcular el porcentaje de error probable de la referencia en función de su punto de existencia. & Sí & Servicio de existencias / Servicio de oferta comercial / Servicio de abastecimiento & Subdoc.3 §3.3.2 (Frente Retail); Anexo A.4.1 (SUP-01, SUP-17, SUP-18, SUP-20) \textbar{} Subdoc.4 §4.2 \\
\addlinespace[2pt]
RF-160 & El sistema deberá mostrar al cliente el porcentaje de error probable calculado en RF-058 junto a la disponibilidad publicada. & Sí & Servicio de existencias / Servicio de oferta comercial / Servicio de abastecimiento & Subdoc.3 §3.3.2 (Frente Retail); Anexo A.4.1 (SUP-01, SUP-17, SUP-18, SUP-20) \textbar{} Subdoc.4 §4.2 \\
\addlinespace[2pt]
RF-161 & El sistema deberá mostrar al vendedor de piso el porcentaje de error probable calculado en RF-058. & Sí & Servicio de existencias / Servicio de oferta comercial / Servicio de abastecimiento & Subdoc.3 §3.3.2 (Frente Retail); Anexo A.4.1 (SUP-01, SUP-17, SUP-18, SUP-20) \textbar{} Subdoc.4 §4.2 \\
\addlinespace[2pt]
RF-162 & El sistema deberá generar una alerta interna dirigida al rol facultado cuando la exactitud de una categoría caiga bajo el umbral definido. & Sí & Servicio de existencias / Servicio de oferta comercial / Servicio de abastecimiento & Subdoc.3 §3.3.2 (Frente Retail); Anexo A.4.1 (SUP-01, SUP-17, SUP-18, SUP-20) \textbar{} Subdoc.4 §4.2 \\
\addlinespace[2pt]
RF-163 & El sistema deberá permitir exclusivamente al rol facultado suspender manualmente la publicación de una categoría. & Sí & Servicio de existencias / Servicio de oferta comercial / Servicio de abastecimiento & Subdoc.3 §3.3.2 (Frente Retail); Anexo A.4.1 (SUP-01, SUP-17, SUP-18, SUP-20) \textbar{} Subdoc.4 §4.2 \\
\addlinespace[2pt]
RF-164 & El sistema deberá abstenerse de suspender automáticamente la publicación de una categoría que no tenga una regla de degradación declarada y aprobada previamente. & Sí & Servicio de existencias / Servicio de oferta comercial / Servicio de abastecimiento & Subdoc.3 §3.3.2 (Frente Retail); Anexo A.4.1 (SUP-01, SUP-17, SUP-18, SUP-20) \textbar{} Subdoc.4 §4.2 \\
\addlinespace[2pt]
RF-165 & El sistema deberá registrar cada cruce de información ejecutado entre ámbitos, indicando el dato cruzado, la finalidad, la base de licitud, la autorización nominada y el instante. & Sí & Servicio de control de cruces (X-01) & Subdoc.3 §3.3.2 (Frente Frontera); Anexo A.4.1 (SUP-03, SUP-04) \textbar{} Subdoc.4 §4.1 (Tabla 4.2) \\
\addlinespace[2pt]
RF-166 & El sistema deberá bloquear todo intento de cruce de información que no corresponda a una interfaz declarada en el inventario de flujos autorizados. & Sí & Servicio de control de cruces (X-01) & Subdoc.3 §3.3.2 (Frente Frontera); Anexo A.4.1 (SUP-03, SUP-04) \textbar{} Subdoc.4 §4.1 (Tabla 4.2) \\
\addlinespace[2pt]
RF-167 & El sistema deberá registrar el intento de cruce bloqueado con el componente origen, el dato solicitado y el instante. & Sí & Servicio de control de cruces (X-01) & Subdoc.3 §3.3.2 (Frente Frontera); Anexo A.4.1 (SUP-03, SUP-04) \textbar{} Subdoc.4 §4.1 (Tabla 4.2) \\
\addlinespace[2pt]
RF-168 & El sistema debe permitir al oficial de cumplimiento mantener el inventario de interfaces de cruce declaradas. & Sí & Servicio de control de cruces (X-01) & Subdoc.3 §3.3.2 (Frente Frontera); Anexo A.4.1 (SUP-03, SUP-04) \textbar{} Subdoc.4 §4.1 (Tabla 4.2) \\
\addlinespace[2pt]
RF-169 & El sistema debe impedir el registro de una interfaz de cruce que no declare finalidad, base de licitud y autorización nominada. & Sí & Servicio de control de cruces (X-01) & Subdoc.3 §3.3.2 (Frente Frontera); Anexo A.4.1 (SUP-03, SUP-04) \textbar{} Subdoc.4 §4.1 (Tabla 4.2) \\
\addlinespace[2pt]
RF-170 & El sistema deberá excluir del catálogo de atributos disponibles en el motor de campañas todo atributo de origen financiero (mora, deuda, cupo utilizado, comportamiento de pago). & Sí & Servicio de control de cruces (X-01) & Subdoc.3 §3.3.2 (Frente Frontera); Anexo A.4.1 (SUP-03, SUP-04) \textbar{} Subdoc.4 §4.1 (Tabla 4.2) \\
\addlinespace[2pt]
RF-171 & El sistema deberá rechazar la ejecución de toda facilidad comercial que invoque un atributo de origen financiero sin un registro de cruce autorizado asociado. & Sí & Servicio de control de cruces (X-01) & Subdoc.3 §3.3.2 (Frente Frontera); Anexo A.4.1 (SUP-03, SUP-04) \textbar{} Subdoc.4 §4.1 (Tabla 4.2) \\
\addlinespace[2pt]
RF-172 & El sistema deberá rechazar la ejecución de todo proceso crediticio que invoque un atributo de origen retail sin un registro de cruce autorizado asociado. & Sí & Servicio de control de cruces (X-01) & Subdoc.3 §3.3.2 (Frente Frontera); Anexo A.4.1 (SUP-03, SUP-04) \textbar{} Subdoc.4 §4.1 (Tabla 4.2) \\
\addlinespace[2pt]
RF-173 & El sistema deberá asignar a la misma persona un identificador distinto en el ámbito retail y en el ámbito fiscalizado. & Sí & Servicio de clientes Retail / Servicio de originación de crédito / Servicio de control de cruces (X-01) & Subdoc.3 §3.3.2 (Frente Retail, Frente Filial emisora, Frente Frontera); Anexo A.4.1 (SUP-03, SUP-04) \textbar{} Subdoc.4 §4.1 (Tabla 4.2) \\
\addlinespace[2pt]
RF-174 & El sistema deberá resolver la correspondencia entre identificadores únicamente durante un cruce autorizado, sin conservarla en una tabla ni en ningún atributo de ambos ámbitos, y con acceso nominado y registrado. & Sí & Servicio de control de cruces (X-01): la correspondencia se resuelve solo durante un cruce autorizado y no se conserva & Subdoc.3 §3.3.2 (Frente Frontera); Anexo A.4.1 (SUP-03, SUP-04) \textbar{} Subdoc.4 §4.1 (Tabla 4.2) \\
\addlinespace[2pt]
RF-175 & El sistema deberá rechazar la persistencia de la clave del ámbito financiero como columna o atributo en cualquier entidad del ámbito retail. & Sí & Servicio de control de cruces (X-01) & Subdoc.3 §3.3.2 (Frente Frontera); Anexo A.4.1 (SUP-03, SUP-04) \textbar{} Subdoc.4 §4.1 (Tabla 4.2) \\
\addlinespace[2pt]
RF-176 & El sistema deberá exigir el registro de la evaluación de impacto sobre la frontera de datos antes de habilitar el estado 'aprobada' de una iniciativa. & Sí & Servicio de control de cruces (X-01) & Subdoc.3 §3.3.2 (Frente Frontera); Anexo A.4.1 (SUP-03, SUP-04) \textbar{} Subdoc.4 §4.1 (Tabla 4.2) \\
\addlinespace[2pt]
RF-177 & El sistema debe permitir declarar previamente el orden de degradación de los servicios. & Sí & Base tecnológica (observabilidad y degradación controlada) & Subdoc.3 §3.3.2 (plataforma común); Anexo A.4.1 (SUP-23) \textbar{} Subdoc.4 §4.1 \\
\addlinespace[2pt]
RF-178 & El sistema debe permitir parametrizar los criterios que activan cada nivel de degradación declarado. & Sí & Base tecnológica (observabilidad y degradación controlada) & Subdoc.3 §3.3.2 (plataforma común); Anexo A.4.1 (SUP-23) \textbar{} Subdoc.4 §4.1 \\
\addlinespace[2pt]
RF-179 & El sistema deberá suspender la publicación de la categoría indicada por el orden de degradación declarado. & Sí & Base tecnológica (observabilidad y degradación controlada) & Subdoc.3 §3.3.2 (plataforma común); Anexo A.4.1 (SUP-23) \textbar{} Subdoc.4 §4.1 \\
\addlinespace[2pt]
RF-180 & El sistema deberá reducir el límite de unidades por cliente al valor declarado. & Sí & Base tecnológica (observabilidad y degradación controlada) & Subdoc.3 §3.3.2 (plataforma común); Anexo A.4.1 (SUP-23) \textbar{} Subdoc.4 §4.1 \\
\addlinespace[2pt]
RF-181 & El sistema deberá desactivar temporalmente los medios de pago de mayor fricción operativa declarados. & Sí & Base tecnológica (observabilidad y degradación controlada) & Subdoc.3 §3.3.2 (plataforma común); Anexo A.4.1 (SUP-23) \textbar{} Subdoc.4 §4.1 \\
\addlinespace[2pt]
RF-182 & El sistema deberá monitorear en tiempo real la tasa de cancelación por falta de existencia por categoría. & Sí & Base tecnológica (observabilidad y degradación controlada) & Subdoc.3 §3.3.2 (plataforma común); Anexo A.4.1 (SUP-23) \textbar{} Subdoc.4 §4.1 \\
\addlinespace[2pt]
RF-183 & El sistema deberá activar la acción de degradación declarada para esa categoría en RF-177. & Sí & Base tecnológica (observabilidad y degradación controlada) & Subdoc.3 §3.3.2 (plataforma común); Anexo A.4.1 (SUP-23) \textbar{} Subdoc.4 §4.1 \\
\addlinespace[2pt]
RF-184 & El sistema deberá permitir parametrizar las cinco ventanas de congelamiento (1 nov - 6 ene; evento anual y su semana previa; evento de noviembre; semana del Día de la Madre; última semana de enero a primera de marzo). & Sí & Base tecnológica (observabilidad y degradación controlada) & Subdoc.3 §3.3.2 (plataforma común); Anexo A.4.1 (SUP-23) \textbar{} Subdoc.4 §4.1 \\
\addlinespace[2pt]
RF-185 & El sistema deberá impedir la ejecución del despliegue o intervención en producción durante la ventana de congelamiento. & Sí & Base tecnológica (observabilidad y degradación controlada) & Subdoc.3 §3.3.2 (plataforma común); Anexo A.4.1 (SUP-23) \textbar{} Subdoc.4 §4.1 \\
\addlinespace[2pt]
RF-186 & El sistema deberá registrar el intento de despliegue bloqueado, indicando componente, solicitante e instante. & Sí & Base tecnológica (observabilidad y degradación controlada) & Subdoc.3 §3.3.2 (plataforma común); Anexo A.4.1 (SUP-23) \textbar{} Subdoc.4 §4.1 \\
\addlinespace[2pt]
RF-187 & El sistema debe permitir al ejecutivo de mesón atender íntegramente el caso de garantía legal en la tienda. & Sí & Servicio de posventa & Subdoc.3 §3.3.2 (Frente Retail); Anexo A.4.1 (SUP-21, SUP-22) \\
\addlinespace[2pt]
RF-188 & El sistema debe impedir que el flujo de atención exija la derivación del consumidor al fabricante, al servicio técnico o al vendedor de marketplace como condición para ser atendido. & Sí & Servicio de posventa & Subdoc.3 §3.3.2 (Frente Retail); Anexo A.4.1 (SUP-21, SUP-22) \\
\addlinespace[2pt]
RF-189 & El sistema deberá permitir al ejecutivo registrar el motivo de la devolución (talla, color, expectativa, falla). & Sí & Servicio de posventa & Subdoc.3 §3.3.2 (Frente Retail); Anexo A.4.1 (SUP-21, SUP-22) \\
\addlinespace[2pt]
RF-190 & El sistema deberá impedir el reingreso de una unidad devuelta al inventario disponible mientras no exista una decisión de aptitud registrada. & Sí & Servicio de posventa & Subdoc.3 §3.3.2 (Frente Retail); Anexo A.4.1 (SUP-21, SUP-22) \\
\addlinespace[2pt]
RF-191 & El sistema deberá permitir al ejecutivo registrar la decisión de aptitud de la unidad devuelta, identificando al responsable de la decisión y su instante. & Sí & Servicio de posventa & Subdoc.3 §3.3.2 (Frente Retail); Anexo A.4.1 (SUP-21, SUP-22) \\
\addlinespace[2pt]
RF-192 & El sistema deberá presentar al ejecutivo las tres opciones de garantía legal (reparación, reposición y devolución del precio) para que el consumidor elija. & Sí & Servicio de posventa & Subdoc.3 §3.3.2 (Frente Retail); Anexo A.4.1 (SUP-21, SUP-22) \\
\addlinespace[2pt]
RF-193 & El sistema debe registrar la opción de garantía legal ofrecida al consumidor. & Sí & Servicio de posventa & Subdoc.3 §3.3.2 (Frente Retail); Anexo A.4.1 (SUP-21, SUP-22) \\
\addlinespace[2pt]
RF-194 & El sistema debe registrar la opción de garantía legal efectivamente elegida por el consumidor. & Sí & Servicio de posventa & Subdoc.3 §3.3.2 (Frente Retail); Anexo A.4.1 (SUP-21, SUP-22) \\
\addlinespace[2pt]
RF-195 & El sistema deberá permitir parametrizar el plazo de garantía legal por tipo de producto, sin requerir modificacion de código. & Sí & Servicio de posventa & Subdoc.3 §3.3.2 (Frente Retail); Anexo A.4.1 (SUP-21, SUP-22) \\
\addlinespace[2pt]
RF-196 & El sistema deberá registrar el hito de resolución al consumidor con su fecha propia. & Sí & Servicio de posventa & Subdoc.3 §3.3.2 (Frente Retail); Anexo A.4.1 (SUP-21, SUP-22) \\
\addlinespace[2pt]
RF-197 & El sistema deberá registrar el hito de recuperación contra el tercero responsable con una fecha independiente de la resolución al consumidor. & Sí & Servicio de posventa & Subdoc.3 §3.3.2 (Frente Retail); Anexo A.4.1 (SUP-21, SUP-22) \\
\addlinespace[2pt]
RF-198 & El sistema deberá impedir que el estado de la recuperación contra el tercero condicione el cierre de la resolución al consumidor. & Sí & Servicio de posventa & Subdoc.3 §3.3.2 (Frente Retail); Anexo A.4.1 (SUP-21, SUP-22) \\
\addlinespace[2pt]
RF-199 & El sistema deberá exigir la entrega completa de la información precontractual antes de habilitar la evaluación de la solicitud de crédito. & Sí & Servicio de originación de crédito / Servicio de evidencia financiera / Servicio de cartera de crédito & Subdoc.3 §3.3.2 (Frente Filial emisora); Anexo A.4.1 (SUP-25, SUP-05, SUP-06) \textbar{} Subdoc.4 §4.1 (Tabla 4.2) \\
\addlinespace[2pt]
RF-200 & El sistema deberá registrar la versión del documento precontractual entregado. & Sí & Servicio de originación de crédito / Servicio de evidencia financiera / Servicio de cartera de crédito & Subdoc.3 §3.3.2 (Frente Filial emisora); Anexo A.4.1 (SUP-25, SUP-05, SUP-06) \textbar{} Subdoc.4 §4.1 (Tabla 4.2) \\
\addlinespace[2pt]
RF-201 & El sistema deberá registrar el instante de entrega de la información precontractual. & Sí & Servicio de originación de crédito / Servicio de evidencia financiera / Servicio de cartera de crédito & Subdoc.3 §3.3.2 (Frente Filial emisora); Anexo A.4.1 (SUP-25, SUP-05, SUP-06) \textbar{} Subdoc.4 §4.1 (Tabla 4.2) \\
\addlinespace[2pt]
RF-202 & El sistema deberá registrar el instante de aceptación del cliente. & Sí & Servicio de originación de crédito / Servicio de evidencia financiera / Servicio de cartera de crédito & Subdoc.3 §3.3.2 (Frente Filial emisora); Anexo A.4.1 (SUP-25, SUP-05, SUP-06) \textbar{} Subdoc.4 §4.1 (Tabla 4.2) \\
\addlinespace[2pt]
RF-203 & El sistema deberá registrar el medio por el cual se entrego y se aceptó la información precontractual. & Sí & Servicio de originación de crédito / Servicio de evidencia financiera / Servicio de cartera de crédito & Subdoc.3 §3.3.2 (Frente Filial emisora); Anexo A.4.1 (SUP-25, SUP-05, SUP-06) \textbar{} Subdoc.4 §4.1 (Tabla 4.2) \\
\addlinespace[2pt]
RF-204 & El sistema deberá registrar el contenido exacto aceptado por el cliente o su huella de integridad verificable. & Sí & Servicio de originación de crédito / Servicio de evidencia financiera / Servicio de cartera de crédito & Subdoc.3 §3.3.2 (Frente Filial emisora); Anexo A.4.1 (SUP-25, SUP-05, SUP-06) \textbar{} Subdoc.4 §4.1 (Tabla 4.2) \\
\addlinespace[2pt]
RF-205 & El sistema deberá dejar constancia expresa de que el cliente recibió la información antes de aceptar, mediante un mecanismo distinto de la sola firma en papel archivado. & Sí & Servicio de originación de crédito / Servicio de evidencia financiera / Servicio de cartera de crédito & Subdoc.3 §3.3.2 (Frente Filial emisora); Anexo A.4.1 (SUP-25, SUP-05, SUP-06) \textbar{} Subdoc.4 §4.1 (Tabla 4.2) \\
\addlinespace[2pt]
RF-206 & El sistema deberá impedir el registro de la aceptación del crédito mientras no exista acreditación de entrega previa de la información precontractual completa. & Sí & Servicio de originación de crédito / Servicio de evidencia financiera / Servicio de cartera de crédito & Subdoc.3 §3.3.2 (Frente Filial emisora); Anexo A.4.1 (SUP-25, SUP-05, SUP-06) \textbar{} Subdoc.4 §4.1 (Tabla 4.2) \\
\addlinespace[2pt]
RF-207 & El sistema deberá registrar la evidencia del consentimiento expreso e informado del cliente mediante un mecanismo verificable de integridad. & Sí & Servicio de originación de crédito / Servicio de evidencia financiera / Servicio de cartera de crédito & Subdoc.3 §3.3.2 (Frente Filial emisora); Anexo A.4.1 (SUP-25, SUP-05, SUP-06) \textbar{} Subdoc.4 §4.1 (Tabla 4.2) \\
\addlinespace[2pt]
RF-208 & El sistema deberá impedir el registro de una modificacion de condiciones del crédito que no tenga evidencia de consentimiento asociada. & Sí & Servicio de originación de crédito / Servicio de evidencia financiera / Servicio de cartera de crédito & Subdoc.3 §3.3.2 (Frente Filial emisora); Anexo A.4.1 (SUP-25, SUP-05, SUP-06) \textbar{} Subdoc.4 §4.1 (Tabla 4.2) \\
\addlinespace[2pt]
RF-209 & El sistema deberá reconstruir el acto de consentimiento presentando que se informó, en que versión y que aceptó el cliente. & Sí & Servicio de originación de crédito / Servicio de evidencia financiera / Servicio de cartera de crédito & Subdoc.3 §3.3.2 (Frente Filial emisora); Anexo A.4.1 (SUP-25, SUP-05, SUP-06) \textbar{} Subdoc.4 §4.1 (Tabla 4.2) \\
\addlinespace[2pt]
RF-210 & El sistema deberá restaurar desde archivo frío los antecedentes de una operación de crédito de cualquier cohorte dentro del plazo de conservación. & Sí & Servicio de originación de crédito / Servicio de evidencia financiera / Servicio de cartera de crédito & Subdoc.3 §3.3.2 (Frente Filial emisora); Anexo A.4.1 (SUP-25, SUP-05, SUP-06) \textbar{} Subdoc.4 §4.1 (Tabla 4.2) \\
\addlinespace[2pt]
RF-211 & El sistema debe permitir al ejecutivo de cobranza iniciar una gestión de cobranza. & Sí & Servicio de originación de crédito / Servicio de evidencia financiera / Servicio de cartera de crédito & Subdoc.3 §3.3.2 (Frente Filial emisora); Anexo A.4.1 (SUP-25, SUP-05, SUP-06) \textbar{} Subdoc.4 §4.1 (Tabla 4.2) \\
\addlinespace[2pt]
RF-212 & El sistema debe registrar la gestión de cobranza ejecutada con su medio, su instante y su ejecutor. & Sí & Servicio de originación de crédito / Servicio de evidencia financiera / Servicio de cartera de crédito & Subdoc.3 §3.3.2 (Frente Filial emisora); Anexo A.4.1 (SUP-25, SUP-05, SUP-06) \textbar{} Subdoc.4 §4.1 (Tabla 4.2) \\
\addlinespace[2pt]
RF-213 & El sistema debe impedir la ejecución de una gestión de cobranza fuera de los límites normativos de horario y de medio de contacto. & Sí & Servicio de originación de crédito / Servicio de evidencia financiera / Servicio de cartera de crédito & Subdoc.3 §3.3.2 (Frente Filial emisora); Anexo A.4.1 (SUP-25, SUP-05, SUP-06) \textbar{} Subdoc.4 §4.1 (Tabla 4.2) \\
\addlinespace[2pt]
RF-214 & El sistema debe enlazar el expediente de la repactación con la gestión de cobranza que la originó. & Sí & Servicio de originación de crédito / Servicio de evidencia financiera / Servicio de cartera de crédito & Subdoc.3 §3.3.2 (Frente Filial emisora); Anexo A.4.1 (SUP-25, SUP-05, SUP-06) \textbar{} Subdoc.4 §4.1 (Tabla 4.2) \\
\addlinespace[2pt]
RF-215 & El sistema debe enlazar el expediente de la repactación con la evidencia de consentimiento registrada en RF-207. & Sí & Servicio de originación de crédito / Servicio de evidencia financiera / Servicio de cartera de crédito & Subdoc.3 §3.3.2 (Frente Filial emisora); Anexo A.4.1 (SUP-25, SUP-05, SUP-06) \textbar{} Subdoc.4 §4.1 (Tabla 4.2) \\
\addlinespace[2pt]
RF-216 & El sistema deberá generar un reporte de conciliación diaria de saldos durante todo el proceso de migración. & Sí & Servicio de originación de crédito / Servicio de evidencia financiera / Servicio de cartera de crédito & Subdoc.3 §3.3.2 (Frente Filial emisora); Anexo A.4.1 (SUP-25, SUP-05, SUP-06) \textbar{} Subdoc.4 §4.1 (Tabla 4.2) \\
\addlinespace[2pt]
RF-217 & El sistema deberá permitir al cliente no autenticado consultar la información precontractual del crédito. & Sí & Servicio de originación de crédito / Servicio de evidencia financiera / Servicio de cartera de crédito & Subdoc.3 §3.3.2 (Frente Filial emisora); Anexo A.4.1 (SUP-25, SUP-05, SUP-06) \textbar{} Subdoc.4 §4.1 (Tabla 4.2) \\
\addlinespace[2pt]
RF-218 & El sistema deberá permitir al cliente no autenticado utilizar un simulador de costo total del crédito. & Sí & Servicio de originación de crédito / Servicio de evidencia financiera / Servicio de cartera de crédito & Subdoc.3 §3.3.2 (Frente Filial emisora); Anexo A.4.1 (SUP-25, SUP-05, SUP-06) \textbar{} Subdoc.4 §4.1 (Tabla 4.2) \\
\addlinespace[2pt]
RF-219 & El sistema deberá mantener cargada la Tasa Máxima Convencional vigente por tipo y tramo de operación, con su fecha de vigencia. & Sí & Servicio de originación de crédito / Servicio de evidencia financiera / Servicio de cartera de crédito & Subdoc.3 §3.3.2 (Frente Filial emisora); Anexo A.4.1 (SUP-25, SUP-05, SUP-06) \textbar{} Subdoc.4 §4.1 (Tabla 4.2) \\
\addlinespace[2pt]
RF-220 & El sistema deberá impedir la originación de una operación cuya tasa supere la Tasa Máxima Convencional vigente a la fecha de la operación. & Sí & Servicio de originación de crédito / Servicio de evidencia financiera / Servicio de cartera de crédito & Subdoc.3 §3.3.2 (Frente Filial emisora); Anexo A.4.1 (SUP-25, SUP-05, SUP-06) \textbar{} Subdoc.4 §4.1 (Tabla 4.2) \\
\addlinespace[2pt]
RF-221 & El sistema deberá determinar el cupo exclusivamente a partir de las variables de la evaluación de capacidad de pago documentada. & Sí & Servicio de originación de crédito / Servicio de evidencia financiera / Servicio de cartera de crédito & Subdoc.3 §3.3.2 (Frente Filial emisora); Anexo A.4.1 (SUP-25, SUP-05, SUP-06) \textbar{} Subdoc.4 §4.1 (Tabla 4.2) \\
\addlinespace[2pt]
RF-222 & El sistema deberá excluir el monto de la venta en curso del conjunto de variables de asignación de cupo. & Sí & Servicio de originación de crédito / Servicio de evidencia financiera / Servicio de cartera de crédito & Subdoc.3 §3.3.2 (Frente Filial emisora); Anexo A.4.1 (SUP-25, SUP-05, SUP-06) \textbar{} Subdoc.4 §4.1 (Tabla 4.2) \\
\addlinespace[2pt]
RF-223 & El sistema deberá excluir la solicitud del vendedor del conjunto de variables de asignación de cupo. & Sí & Servicio de originación de crédito / Servicio de evidencia financiera / Servicio de cartera de crédito & Subdoc.3 §3.3.2 (Frente Filial emisora); Anexo A.4.1 (SUP-25, SUP-05, SUP-06) \textbar{} Subdoc.4 §4.1 (Tabla 4.2) \\
\addlinespace[2pt]
RF-224 & El sistema deberá impedir al vendedor alterar el resultado de una evaluación crediticia. & Sí & Servicio de originación de crédito / Servicio de evidencia financiera / Servicio de cartera de crédito & Subdoc.3 §3.3.2 (Frente Filial emisora); Anexo A.4.1 (SUP-25, SUP-05, SUP-06) \textbar{} Subdoc.4 §4.1 (Tabla 4.2) \\
\addlinespace[2pt]
RF-225 & El sistema deberá impedir la repetición de la evaluación crediticia del mismo cliente dentro de la ventana de enfriamiento parametrizada. & Sí & Servicio de originación de crédito / Servicio de evidencia financiera / Servicio de cartera de crédito & Subdoc.3 §3.3.2 (Frente Filial emisora); Anexo A.4.1 (SUP-25, SUP-05, SUP-06) \textbar{} Subdoc.4 §4.1 (Tabla 4.2) \\
\addlinespace[2pt]
RF-226 & El sistema deberá registrar cada intento de evaluación con la identidad del solicitante y el resultado obtenido. & Sí & Servicio de originación de crédito / Servicio de evidencia financiera / Servicio de cartera de crédito & Subdoc.3 §3.3.2 (Frente Filial emisora); Anexo A.4.1 (SUP-25, SUP-05, SUP-06) \textbar{} Subdoc.4 §4.1 (Tabla 4.2) \\
\addlinespace[2pt]
RF-227 & El sistema deberá consolidar en un único registro retail los registros duplicados de un mismo cliente. & Sí & R:CL-02 Servicio de clientes Retail (Etapa 2) & Subdoc.3 §3.3.2 (Frente Retail) \textbar{} Subdoc.4 §4.1 (Tabla 4.2) \\
\addlinespace[2pt]
RF-228 & El sistema deberá registrar cada movimiento de puntos de fidelización de un cliente. & Sí & R:CL-02 Servicio de clientes Retail (Etapa 2) & Subdoc.3 §3.3.2 (Frente Retail) \textbar{} Subdoc.4 §4.1 (Tabla 4.2) \\
\addlinespace[2pt]
RF-229 & El sistema deberá construir los segmentos de clientes del retail exclusivamente con atributos comerciales. & Sí & R:CL-02 Servicio de clientes Retail (Etapa 2) & Subdoc.3 §3.3.2 (Frente Retail) \textbar{} Subdoc.4 §4.1 (Tabla 4.2) \\
\addlinespace[2pt]
RF-230 & El sistema deberá ejecutar campañas sobre los segmentos de clientes del retail. & Sí & R:CL-02 Servicio de clientes Retail (Etapa 2) & Subdoc.3 §3.3.2 (Frente Retail) \textbar{} Subdoc.4 §4.1 (Tabla 4.2) \\
\addlinespace[2pt]
RF-231 & El sistema deberá sincronizar con el sistema de fidelización conservado los datos de fidelización del cliente. & Sí & R:CL-02 Servicio de clientes Retail (Etapa 2) & Subdoc.3 §3.3.2 (Frente Retail) \textbar{} Subdoc.4 §4.1 (Tabla 4.2) \\
\addlinespace[2pt]
RNF-02 & La interfaz de declaración de existencia del vendedor de marketplace debe estar disponible de forma continua. & Sí & Nodo local de tienda (SIT-01) + Azure Chile Central en tres zonas de disponibilidad (CLD-01) con recuperación en Azure Brazil South (CLD-02) + centro de datos de casa matriz (DC-01) con sala de respaldo (DC-02), con meta de 99,9 \% mensual para los servicios críticos (RT-10.01) & Subdoc.4 §4.2 \\
\addlinespace[2pt]
RNF-03 & El registro del conteo cíclico no debe bloquear ni degradar la operación de venta de la sala. & Sí & Arquitectura orientada a eventos con presupuesto de latencia por componente (Servicio de existencias, Servicio de oferta comercial, punto de venta, Servicio de pedidos, canal digital) & Subdoc.4 §4.1, §4.2 \\
\addlinespace[2pt]
RNF-04 & El estado de un pedido no debe presentar discrepancias entre canales de consulta. & Sí & Canal digital / Servicio de pedidos — estado único del pedido consultado por todos los canales & Subdoc.3 §3.3.2 (Frente Retail) \\
\addlinespace[2pt]
RNF-05 & La actualización del estado de un pedido debe reflejarse en todos los puntos de consulta dentro del plazo máximo definido. & Sí & Arquitectura orientada a eventos con presupuesto de latencia por componente (Servicio de existencias, Servicio de oferta comercial, punto de venta, Servicio de pedidos, canal digital) & Subdoc.4 §4.1, §4.2 \\
\addlinespace[2pt]
RNF-06 & El tiempo de evaluación de una solicitud de crédito en el punto de venta físico no debe exceder el umbral objetivo. & Sí & Servicio de originación de crédito: evaluación en el flujo del punto de venta con umbral de 8 s y presupuesto de latencia por componente & Subdoc.4 §4.1, §4.2 \\
\addlinespace[2pt]
RNF-07 & El tiempo total de originación de crédito en el mesón, con información precontractual acreditada, no debe superar el tiempo del proceso actual. & Sí & Servicio de originación de crédito: tiempo total de originación en el mesón con información precontractual acreditada, umbral de 3 minutos & Subdoc.3 §3.3.2 (Frente Filial emisora) \textbar{} Subdoc.4 §4.1 (Tabla 4.2), §4.2 \\
\addlinespace[2pt]
RNF-08 & La migración de la cartera de crédito viva no debe generar pérdida de datos. & Sí & Servicio de cartera de crédito (olas de coexistencia, circuit breaker) & Subdoc.3 §3.3.2 (Frente Filial emisora); Anexo A.4.1 (SUP-25) \textbar{} Subdoc.4 §4.1 (Tabla 4.2) \\
\addlinespace[2pt]
RNF-09 & La migración de la cartera de crédito viva no debe interrumpir el servicio de cobro. & Sí & Servicio de cartera de crédito (olas de coexistencia, circuit breaker) & Subdoc.3 §3.3.2 (Frente Filial emisora); Anexo A.4.1 (SUP-25) \textbar{} Subdoc.4 §4.1 (Tabla 4.2) \\
\addlinespace[2pt]
RNF-10 & La migración de la cartera de crédito viva no debe producir divergencia de saldos. & Sí & Servicio de cartera de crédito (olas de coexistencia, circuit breaker) & Subdoc.3 §3.3.2 (Frente Filial emisora); Anexo A.4.1 (SUP-25) \textbar{} Subdoc.4 §4.1 (Tabla 4.2) \\
\addlinespace[2pt]
RNF-11 & La separación lógica entre los datos del retail y los del negocio financiero debe estar implementada. & Sí & Límites de contexto (DDD) y ámbitos Retail y Emisor con datos, identidades, redes y claves propios en CLD-01 (ADR-02 y ADR-05) & Subdoc.4 §4.1, §4.2 \\
\addlinespace[2pt]
RNF-12 & La separación lógica entre ambos ámbitos debe estar documentada. & Sí & Límites de contexto (DDD) y ámbitos Retail y Emisor con datos, identidades, redes y claves propios en CLD-01 (ADR-02 y ADR-05) & Subdoc.4 §4.1, §4.2 \\
\addlinespace[2pt]
RNF-13 & La separación lógica entre ambos ámbitos debe ser verificada por auditoría independiente. & Sí & Límites de contexto (DDD) y ámbitos Retail y Emisor con datos, identidades, redes y claves propios en CLD-01 (ADR-02 y ADR-05) & Subdoc.4 §4.1, §4.2 \\
\addlinespace[2pt]
RNF-14 & La separación física entre la infraestructura del retail y la de la filial emisora debe estar implementada y acreditada. & Sí & Límites de contexto (DDD) y ámbitos Retail y Emisor con datos, identidades, redes y claves propios en CLD-01 (ADR-02 y ADR-05) & Subdoc.4 §4.1, §4.2 \\
\addlinespace[2pt]
RNF-15 & Toda acción relevante ejecutada en el sistema debe quedar registrada indicando usuario individual, acción e instante. & Sí & Servicio de control de cruces (X-01) / Base tecnológica (identidad y gestión de accesos) — registro inalterable de acciones & Subdoc.3 §3.3.2 (Frente Frontera, plataforma común) \textbar{} Subdoc.4 §4.1 (Tabla 4.2) \\
\addlinespace[2pt]
RNF-16 & La consulta de disponibilidad en la ficha de producto del canal digital debe responder dentro del umbral definido. & Sí & Arquitectura orientada a eventos con presupuesto de latencia por componente (Servicio de existencias, Servicio de oferta comercial, punto de venta, Servicio de pedidos, canal digital) & Subdoc.4 §4.1, §4.2 \\
\addlinespace[2pt]
RNF-17 & La confirmación de un pedido durante el evento anual debe completarse dentro del umbral definido. & Sí & Arquitectura orientada a eventos con presupuesto de latencia por componente (Servicio de existencias, Servicio de oferta comercial, punto de venta, Servicio de pedidos, canal digital) & Subdoc.4 §4.1, §4.2 \\
\addlinespace[2pt]
RNF-18 & La venta completa en caja con medio de pago externo debe completarse dentro del umbral definido. & Sí & Arquitectura orientada a eventos con presupuesto de latencia por componente (Servicio de existencias, Servicio de oferta comercial, punto de venta, Servicio de pedidos, canal digital) & Subdoc.4 §4.1, §4.2 \\
\addlinespace[2pt]
RNF-19 & La propagación de un cambio de precio a las 380 líneas de caja y al canal digital debe completarse dentro del umbral definido. & Sí & Arquitectura orientada a eventos con presupuesto de latencia por componente (Servicio de existencias, Servicio de oferta comercial, punto de venta, Servicio de pedidos, canal digital) & Subdoc.4 §4.1, §4.2 \\
\addlinespace[2pt]
RNF-20 & El registro de una devolución en el mesón debe completarse dentro del umbral definido. & Sí & Arquitectura orientada a eventos con presupuesto de latencia por componente (Servicio de existencias, Servicio de oferta comercial, punto de venta, Servicio de pedidos, canal digital) & Subdoc.4 §4.1, §4.2 \\
\addlinespace[2pt]
RNF-21 & La consulta de disponibilidad desde una terminal compartida del piso de venta debe responder dentro del umbral definido. & Sí & Arquitectura orientada a eventos con presupuesto de latencia por componente (Servicio de existencias, Servicio de oferta comercial, punto de venta, Servicio de pedidos, canal digital) & Subdoc.4 §4.1, §4.2 \\
\addlinespace[2pt]
RNF-22 & La plataforma debe soportar la concurrencia y el volumen del peak digital del evento anual sin degradar los umbrales de NFR-06.x. & Sí & Escalamiento elástico en CLD-01 y capacidad fija en SIT-01 y DC-01; prueba de carga a 1,5 veces el peak declarado (RT-09.06) & Subdoc.4 §4.2, §4.3 \\
\addlinespace[2pt]
RNF-23 & La plataforma debe soportar la concurrencia y el volumen del peak presencial de la campaña de noviembre y diciembre sin degradar los umbrales de NFR-06.x. & Sí & Escalamiento elástico en CLD-01 y capacidad fija en SIT-01 y DC-01; prueba de carga a 1,5 veces el peak declarado (RT-09.06) & Subdoc.4 §4.2, §4.3 \\
\addlinespace[2pt]
RNF-24 & Una tienda debe poder operar sin enlace externo durante el período mínimo definido. & Sí & Nodo local de tienda (SIT-01) + Azure Chile Central en tres zonas de disponibilidad (CLD-01) con recuperación en Azure Brazil South (CLD-02) + centro de datos de casa matriz (DC-01) con sala de respaldo (DC-02) & Subdoc.4 §4.2 \\
\addlinespace[2pt]
RNF-25 & El centro de distribución principal debe poder operar sin enlace externo durante el período mínimo definido. & Parcial & Nodo local de tienda con autonomía de 24 h; para el centro de distribución principal el Subdocumento 4 compromete 4 h, por debajo de las 24 h exigidas, y debe corregirse o declarar el diseño que lo logre & Subdoc.4 §4.2 (Alta disponibilidad, respaldos, conexiones y puntos de falla) \\
\addlinespace[2pt]
RNF-26 & La sincronización tras la reconexión no debe superar el plazo definido. & Sí & Nodo local de tienda (SIT-01) + Azure Chile Central en tres zonas de disponibilidad (CLD-01) con recuperación en Azure Brazil South (CLD-02) + centro de datos de casa matriz (DC-01) con sala de respaldo (DC-02) & Subdoc.4 §4.2 \\
\addlinespace[2pt]
RNF-27 & La sincronización tras la reconexión no debe producir pérdida de ventas ni de documentos. & Sí & Nodo local de tienda (SIT-01) + Azure Chile Central en tres zonas de disponibilidad (CLD-01) con recuperación en Azure Brazil South (CLD-02) + centro de datos de casa matriz (DC-01) con sala de respaldo (DC-02) & Subdoc.4 §4.2 \\
\addlinespace[2pt]
RNF-28 & La operación desconectada autónoma on-premise debe sostenerse durante el período mínimo exigido por las bases técnicas. & Parcial & Operación desconectada de 24 h en el nodo local de tienda; centro de distribución principal por alinear con el mínimo exigido (hoy 4 h en el Subdocumento 4) & Subdoc.4 §4.2 (Alta disponibilidad, respaldos, conexiones y puntos de falla) \\
\addlinespace[2pt]
RNF-29 & La operación de tiendas y centros de distribución debe estar disponible en la ventana horaria definida. & Sí & Nodo local de tienda (SIT-01) + Azure Chile Central en tres zonas de disponibilidad (CLD-01) con recuperación en Azure Brazil South (CLD-02) + centro de datos de casa matriz (DC-01) con sala de respaldo (DC-02) & Subdoc.4 §4.2 \\
\addlinespace[2pt]
RNF-30 & El canal digital debe estar disponible de forma continua. & Sí & Nodo local de tienda (SIT-01) + Azure Chile Central en tres zonas de disponibilidad (CLD-01) con recuperación en Azure Brazil South (CLD-02) + centro de datos de casa matriz (DC-01) con sala de respaldo (DC-02), con meta de 99,9 \% mensual para los servicios críticos (RT-10.01) & Subdoc.4 §4.2 \\
\addlinespace[2pt]
RNF-31 & Los servicios financieros que afectan pagos, estados de cuenta y bloqueo de tarjetas deben estar disponibles de forma continua. & Sí & Nodo local de tienda (SIT-01) + Azure Chile Central en tres zonas de disponibilidad (CLD-01) con recuperación en Azure Brazil South (CLD-02) + centro de datos de casa matriz (DC-01) con sala de respaldo (DC-02), con meta de 99,9 \% mensual para los servicios críticos (RT-10.01) & Subdoc.4 §4.2 \\
\addlinespace[2pt]
RNF-32 & La solución debe cumplir los objetivos de recuperación ante desastre exigidos. & Sí & CLD-01 con recuperación en CLD-02 (RTO≤4h, RPO≤15min) — Servicio de cartera de crédito: restauración desde archivo frío & Subdoc.4 §4.2 \\
\addlinespace[2pt]
RNF-33 & Los datos de identificación de clientes, la cartera de crédito, el comportamiento de pago y los antecedentes de evaluación crediticia deben cifrarse a nivel de campo. & Sí, condicionado & Zero Trust (NIST SP 800-207) con el Servicio de control de cruces (X-01) y la identidad de cada ámbito; cifrado de campo y tokenización por declarar en el diseño detallado; Managed HSM del Emisor por probar y sujeto a la región & Subdoc.4 §4.1 \\
\addlinespace[2pt]
RNF-34 & Los datos de medios de pago no deben almacenarse en claro, sino tokenizarse. & Sí, condicionado & Zero Trust (NIST SP 800-207) con el Servicio de control de cruces (X-01) y la identidad de cada ámbito; cifrado de campo y tokenización por declarar en el diseño detallado; Managed HSM del Emisor por probar y sujeto a la región & Subdoc.4 §4.1 \\
\addlinespace[2pt]
RNF-35 & La red de cajas, la administrativa, la de videovigilancia y la inalámbrica de clientes deben estar segmentadas en las 22 tiendas. & Sí & Arquitectura Zero Trust (NIST SP 800-207): Servicio de control de cruces (X-01) y base tecnológica (identidad y gestión de accesos), cifrado de campo, tokenización, HSM & Subdoc.4 §4.1 \\
\addlinespace[2pt]
RNF-36 & La red y el ámbito de sistemas de la filial emisora deben estar separados de forma acreditada respecto del retail. & Sí & Arquitectura Zero Trust (NIST SP 800-207): Servicio de control de cruces (X-01) y base tecnológica (identidad y gestión de accesos), cifrado de campo, tokenización, HSM & Subdoc.4 §4.1 \\
\addlinespace[2pt]
RNF-37 & Debe existir segregación de funciones verificable entre originación, aprobación, modificacion de condiciones y cobranza. & Sí, condicionado & Zero Trust (NIST SP 800-207) con el Servicio de control de cruces (X-01) y la identidad de cada ámbito; cifrado de campo y tokenización por declarar en el diseño detallado; Managed HSM del Emisor por probar y sujeto a la región & Subdoc.4 §4.1 \\
\addlinespace[2pt]
RNF-38 & Debe registrarse todo acceso a la cartera de crédito, al comportamiento de pago y a los antecedentes de evaluación crediticia. & Sí & Servicio de control de cruces (X-01) / Base tecnológica (identidad y gestión de accesos) — registro inalterable de acciones & Subdoc.3 §3.3.2 (Frente Frontera, plataforma común) \textbar{} Subdoc.4 §4.1 (Tabla 4.2) \\
\addlinespace[2pt]
RNF-39 & Los actos de crédito de apertura de tarjeta y de repactación deben contar con un mecanismo de consentimiento de integridad verificable. & Sí, condicionado & Servicio de evidencia financiera (WORM) con firma electrónica avanzada (Ley 19.799), componente por acreditar & Subdoc.4 §4.1 (Tabla 4.2) \\
\addlinespace[2pt]
RNF-40 & La conformidad de recepción de mercadería y la constancia de devolución deben contar con un mecanismo de consentimiento de integridad verificable. & Sí, condicionado & Servicio de evidencia financiera (WORM) con firma electrónica avanzada (Ley 19.799), componente por acreditar & Subdoc.4 §4.1 (Tabla 4.2) \\
\addlinespace[2pt]
RNF-41 & Los documentos tributarios y antecedentes de venta deben conservarse por el plazo definido. & Sí & Servicio de evidencia financiera (WORM) / Servicio de posventa (garantía legal) / Servicio de control de cruces (X-01) & Subdoc.4 §4.1 (Tabla 4.2) \\
\addlinespace[2pt]
RNF-42 & Los antecedentes del crédito deben conservarse por el plazo del crédito más el período posterior exigido. & Sí & Servicio de evidencia financiera (WORM) / Servicio de posventa (garantía legal) / Servicio de control de cruces (X-01) & Subdoc.4 §4.1 (Tabla 4.2) \\
\addlinespace[2pt]
RNF-43 & La evidencia de consentimiento de originación y repactación debe conservarse por el mismo plazo que el crédito. & Sí & Servicio de evidencia financiera (WORM) / Servicio de posventa (garantía legal) / Servicio de control de cruces (X-01) & Subdoc.4 §4.1 (Tabla 4.2) \\
\addlinespace[2pt]
RNF-44 & La trazabilidad del precio publicado debe conservarse por el plazo definido. & Sí & Servicio de evidencia financiera (WORM) / Servicio de posventa (garantía legal) / Servicio de control de cruces (X-01) & Subdoc.4 §4.1 (Tabla 4.2) \\
\addlinespace[2pt]
RNF-45 & Los movimientos de inventario y conteos cíclicos deben conservarse por el plazo definido. & Sí & Servicio de evidencia financiera (WORM) / Servicio de posventa (garantía legal) / Servicio de control de cruces (X-01) & Subdoc.4 §4.1 (Tabla 4.2) \\
\addlinespace[2pt]
RNF-46 & Los datos de fidelización deben conservarse por el plazo definido. & Sí & Servicio de evidencia financiera (WORM) / Servicio de posventa (garantía legal) / Servicio de control de cruces (X-01) & Subdoc.4 §4.1 (Tabla 4.2) \\
\addlinespace[2pt]
RNF-47 & Los registros de videovigilancia deben conservarse por el plazo definido. & Sí & Servicio de evidencia financiera (WORM) / Servicio de posventa (garantía legal) / Servicio de control de cruces (X-01) & Subdoc.4 §4.1 (Tabla 4.2) \\
\addlinespace[2pt]
RNF-48 & La solución debe admitir la apertura de una tienda nueva por parametrización, sin desarrollo a medida. & Sí & Nodo de tienda como unidad estandarizada replicable (RT-02.12-C) & Subdoc.4 §4.1 \\
\addlinespace[2pt]
RNF-49 & Las interfaces de sala, caja y mesón financiero deben ser operables por personal recién incorporado con capacitación mínima. & Sí & Punto de venta con operación sin conexión (PWA) — interfaces operables con capacitación mínima & Subdoc.4 §4.2 \\
\addlinespace[2pt]
RNF-50 & Toda comunicación relativa al crédito debe cumplir criterios verificables de lenguaje claro, validados con personas usuarias reales. & Sí & Servicio de originación de crédito / Servicio de evidencia financiera — simulador de costo total en lenguaje normalizado & Subdoc.3 §3.3.2 (Frente Filial emisora) \\
\addlinespace[2pt]
RNF-51 & Las notificaciones definidas funcionalmente deben emitirse y registrarse en el 100\% de los eventos que las gatillan. & Sí & Canal digital (notificaciones transaccionales) / Servicio de marketplace & Subdoc.3 §3.3.2 (Frente Retail) \textbar{} Subdoc.4 §4.1 (Tabla 4.2) \\
\addlinespace[2pt]
RNF-52 & La solución debe adoptar un estándar identificado y justificado para el catálogo e identificación de producto. & Sí & Base tecnológica (plataforma de integración) — catálogo de contratos versionados (OpenAPI/AsyncAPI) & Subdoc.4 §4.1 \\
\addlinespace[2pt]
RNF-53 & La solución debe adoptar un estándar identificado y justificado para el intercambio de órdenes y avisos de despacho con los proveedores. & Sí & Base tecnológica (plataforma de integración) — catálogo de contratos versionados (OpenAPI/AsyncAPI) & Subdoc.4 §4.1 \\
\addlinespace[2pt]
RNF-54 & La solución debe adoptar un estándar identificado y justificado para el intercambio con los transportistas de última milla. & Sí & Base tecnológica (plataforma de integración) — catálogo de contratos versionados (OpenAPI/AsyncAPI) & Subdoc.4 §4.1 \\
\addlinespace[2pt]
RNF-55 & La solución debe adoptar un estándar identificado y justificado para la sincronización de catálogo, existencia y pedido con los vendedores de marketplace. & Sí & Base tecnológica (plataforma de integración) — catálogo de contratos versionados (OpenAPI/AsyncAPI) & Subdoc.4 §4.1 \\
\addlinespace[2pt]
RNF-56 & La solución debe adoptar el formato normativo vigente de documento tributario electrónico. & Sí & Base tecnológica (plataforma de integración) — catálogo de contratos versionados (OpenAPI/AsyncAPI) & Subdoc.4 §4.1 \\
\addlinespace[2pt]
RNF-57 & La evidencia completa de consentimiento de una operación arbitraria debe recuperarse dentro del plazo objetivo, a diez años. & Sí & Servicio de evidencia financiera (WORM) — recuperación de evidencia de consentimiento & Subdoc.3 §3.3.2 (Frente Filial emisora); Anexo A.4.1 (SUP-05) \textbar{} Subdoc.4 §4.2 \\
\addlinespace[2pt]
RNF-58 & La recuperación del precio publicado de una referencia en una fecha, hora y canal arbitrarios debe completarse dentro del plazo objetivo. & Sí & Arquitectura orientada a eventos con presupuesto de latencia por componente (Servicio de existencias, Servicio de oferta comercial, punto de venta, Servicio de pedidos, canal digital) & Subdoc.4 §4.1, §4.2 \\
\addlinespace[2pt]
RNF-59 & La revocación total de accesos de un trabajador debe completarse dentro del plazo máximo de la política interna. & Sí, condicionado & Zero Trust (NIST SP 800-207) con el Servicio de control de cruces (X-01) y la identidad de cada ámbito; cifrado de campo y tokenización por declarar en el diseño detallado; Managed HSM del Emisor por probar y sujeto a la región & Subdoc.4 §4.1 \\
\addlinespace[2pt]
RNF-60 & El traspaso de sesión nominativa en una terminal compartida debe completarse dentro del plazo objetivo. & Sí & Base tecnológica (identidad y gestión de accesos) — traspaso rápido de sesión (hot-swapping) & Subdoc.3 §3.3.2 (plataforma común); Anexo A.4.1 (SUP-24) \\
\addlinespace[2pt]
RNF-64 & El tiempo entre la detección de un quiebre y el aviso al cliente no debe superar el objetivo declarado. & Sí & Canal digital (notificaciones transaccionales) / Servicio de marketplace & Subdoc.3 §3.3.2 (Frente Retail) \textbar{} Subdoc.4 §4.1 (Tabla 4.2) \\
\addlinespace[2pt]
RNF-65 & No debe mantenerse un cobro capturado sobre un pedido declarado no cumplible. & Sí & Canal digital (notificaciones transaccionales) / Servicio de marketplace & Subdoc.3 §3.3.2 (Frente Retail) \textbar{} Subdoc.4 §4.1 (Tabla 4.2) \\
\addlinespace[2pt]
RNF-66 & La latencia entre la recepción física de una devolución en tienda y la notificación al vendedor de marketplace no debe superar el objetivo declarado. & Sí & Canal digital (notificaciones transaccionales) / Servicio de marketplace & Subdoc.3 §3.3.2 (Frente Retail) \textbar{} Subdoc.4 §4.1 (Tabla 4.2) \\
\addlinespace[2pt]
RNF-67 & No debe permanecer mercadería de marketplace en bodega sin aviso a su propietario. & Sí & Canal digital (notificaciones transaccionales) / Servicio de marketplace & Subdoc.3 §3.3.2 (Frente Retail) \textbar{} Subdoc.4 §4.1 (Tabla 4.2) \\
\addlinespace[2pt]
RNF-68 & La restauración desde archivo frío de una operación de crédito de la cohorte más antigua debe completarse dentro del plazo objetivo. & Sí & CLD-01 con recuperación en CLD-02 (RTO≤4h, RPO≤15min) — Servicio de cartera de crédito: restauración desde archivo frío & Subdoc.4 §4.2 \\
\addlinespace[2pt]
RNF-69 & La atención de garantía legal en mesón no debe derivar al consumidor a un tercero como condición de atención, en ninguna tienda. & Sí & Servicio de evidencia financiera (WORM) / Servicio de posventa (garantía legal) / Servicio de control de cruces (X-01) & Subdoc.4 §4.1 (Tabla 4.2) \\
\addlinespace[2pt]
RNF-70 & El ciclo de desarrollo debe producir un inventario de componentes (SBOM) por artefacto desplegado. & Sí, condicionado & Cadena DevSecOps (SBOM, SLSA, OWASP ASVS) con herramientas por fijar; Zero Trust según el apartado de seguridad del Subdocumento 4 & Subdoc.4 §4.1, §4.2 \\
\addlinespace[2pt]
RNF-71 & La cadena de suministro de software debe alcanzar el nivel SLSA 3 o superior. & Sí, condicionado & Cadena DevSecOps (SBOM, SLSA, OWASP ASVS) con herramientas por fijar; Zero Trust según el apartado de seguridad del Subdocumento 4 & Subdoc.4 §4.1, §4.2 \\
\addlinespace[2pt]
RNF-72 & La capa expuesta debe cumplir el estándar OWASP ASVS en el nivel comprometido. & Sí, condicionado & Cadena DevSecOps (SBOM, SLSA, OWASP ASVS) con herramientas por fijar; Zero Trust según el apartado de seguridad del Subdocumento 4 & Subdoc.4 §4.1, §4.2 \\
\addlinespace[2pt]
RNF-73 & La arquitectura debe implementar un modelo Zero Trust conforme a NIST SP 800-207. & Sí, condicionado & Cadena DevSecOps (SBOM, SLSA, OWASP ASVS) con herramientas por fijar; Zero Trust según el apartado de seguridad del Subdocumento 4 & Subdoc.4 §4.1, §4.2 \\
\addlinespace[2pt]
RNF-74 & Todo tratamiento de datos personales debe declarar su finalidad. & Sí & Servicio de evidencia financiera (WORM) / Servicio de posventa (garantía legal) / Servicio de control de cruces (X-01) & Subdoc.4 §4.1 (Tabla 4.2) \\
\addlinespace[2pt]
RNF-75 & Todo tratamiento de datos personales debe invocar una de las bases de licitud de la Ley N° 21.719. & Sí & Servicio de evidencia financiera (WORM) / Servicio de posventa (garantía legal) / Servicio de control de cruces (X-01) & Subdoc.4 §4.1 (Tabla 4.2) \\
\addlinespace[2pt]
RNF-76 & La compañía debe mantener el registro de actividades de tratamiento de datos personales. & Sí & Servicio de evidencia financiera (WORM) / Servicio de posventa (garantía legal) / Servicio de control de cruces (X-01) & Subdoc.4 §4.1 (Tabla 4.2) \\
\addlinespace[2pt]
OP-01 & El PROPONENTE debe especificar la alternativa de etiquetas electrónicas frente a las demás opciones de resolución de la discrepancia de precio. & Sí & Obligación del proponente (gestión del proyecto) & Subdocumento 2, Anexo A.3 \textbar{} Formulario T-9 \\
\addlinespace[2pt]
OP-02 & El PROPONENTE debe costear la alternativa de etiquetas electrónicas sin asumir su adquisición dentro del alcance. & Sí & Obligación del proponente (gestión del proyecto) & Subdocumento 2, Anexo A.3 \textbar{} Formulario T-9 \\
\addlinespace[2pt]
OP-03 & El PROPONENTE debe especificar la cobertura de despliegue de las etiquetas electrónicas. & Sí & Obligación del proponente (gestión del proyecto) & Subdocumento 2, Anexo A.3 \textbar{} Formulario T-9 \\
\addlinespace[2pt]
OP-04 & El PROPONENTE debe especificar el hardware requerido por la alternativa de etiquetas electrónicas. & Sí & Obligación del proponente (gestión del proyecto) & Subdocumento 2, Anexo A.3 \textbar{} Formulario T-9 \\
\addlinespace[2pt]
OP-05 & El PROPONENTE debe presentar el costo de la alternativa separadamente del resto de la solución. & Sí & Obligación del proponente (gestión del proyecto) & Subdocumento 2, Anexo A.3 \textbar{} Formulario T-9 \\
\addlinespace[2pt]
OP-06 & El PROPONENTE debe especificar la cantidad de dispositivos móviles requeridos por tienda y por departamento. & Sí & Obligación del proponente (gestión del proyecto) & Subdocumento 2, Anexo A.3 \textbar{} Formulario T-9 \\
\addlinespace[2pt]
OP-07 & El PROPONENTE debe especificar las características técnicas de los dispositivos móviles requeridos. & Sí & Obligación del proponente (gestión del proyecto) & Subdocumento 2, Anexo A.3 \textbar{} Formulario T-9 \\
\addlinespace[2pt]
OP-08 & El PROPONENTE debe presentar el análisis de hacer o comprar del centro de distribución de Concepción con sus tres alternativas. & Sí & Obligación del proponente (gestión del proyecto) & Subdocumento 2, Anexo A.3 \textbar{} Formulario T-9 \\
\addlinespace[2pt]
OP-09 & El PROPONENTE debe declarar el impacto de la alternativa seleccionada sobre el cumplimiento de RN-15. & Sí & Obligación del proponente (gestión del proyecto) & Subdocumento 2, Anexo A.3 \textbar{} Formulario T-9 \\
\addlinespace[2pt]
RT-02.01 & La solución se organizará en las ocho capas del numeral 2.1. El PROPONENTE presentará el diagrama de la arquitectura lógica identificando cada capa, sus componentes y las interfaces entre ellas. & Sí & Estructura por Capas (8 capas), con módulos Ancoa y emplazamiento híbrido asignado a cada una & Subdoc.4 §4.1, §4.2 \\
\addlinespace[2pt]
RT-02.02 & La arquitectura será modular, con límites de contexto explícitos y acoplamiento débil. Se rechazará toda arquitectura monolítica que no permita desplegar de forma independiente sus componentes críticos. & Sí, condicionado & Arquitectura modular con límites de contexto (DDD): diez de los trece servicios como microservicio, dos como módulo y uno por decidir; ADR-05 condicionada & Subdoc.4 §4.1 (Estilo y límites de contexto) \textbar{} Subdoc.3 §3.3.2 \\
\addlinespace[2pt]
RT-02.03 & La descripción de la arquitectura se ajustará a ISO/IEC/IEEE 42010, con vistas lógica, de procesos, de despliegue, de datos y de seguridad. & Sí & Cinco vistas de arquitectura conforme a ISO/IEC/IEEE 42010 (lógica, procesos, despliegue, datos, seguridad) & Subdoc.4 §4.1, §4.2 \\
\addlinespace[2pt]
RT-02.04 & El PROPONENTE mantendrá un registro de decisiones de arquitectura (ADR) fechado, con la alternativa escogida, las alternativas descartadas y el criterio de decisión. El registro es entregable contractual y se actualizará durante toda la ejecución. & Sí & Registro de Decisiones de Arquitectura (ADR-01 a ADR-07, Tabla 4.5 del Subdocumento 4), fechado, con alternativa descartada y criterio & Subdoc.4 §4.1 \\
\addlinespace[2pt]
RT-02.05 & La capa de servicios de negocio será sin estado. El estado de sesión y el estado de proceso residirán en almacenes externos con alta disponibilidad. & Parcial & Servicios de negocio sin estado (Java 25 y Spring Boot 4); estado de sesión en Azure Managed Redis para Retail; almacén de estado del Emisor por definir en el diseño detallado & Subdoc.4 §4.1.1 (Tabla 4.6) \\
\addlinespace[2pt]
RT-02.06 & Toda operación de escritura expuesta a reintentos será idempotente, con clave de idempotencia declarada por el cliente y ventana de deduplicación documentada. & Parcial & Clave de idempotencia por operación y deduplicación en el consumidor; la ventana de deduplicación se documenta en el diseño detallado & Subdoc.4 §4.1 (Recorridos y consistencia) \\
\addlinespace[2pt]
RT-02.07 & Los flujos de eventos garantizarán entrega al menos una vez, con deduplicación en el consumidor y orden garantizado dentro de la partición o del agregado cuando el proceso lo exija. & Sí, condicionado & Bus de eventos con Azure Event Hubs compatible con Kafka (ADR-06, condicionada): entrega al menos una vez, deduplicación en el consumidor y orden por partición & Subdoc.4 §4.1.1 (Tabla 4.6) \\
\addlinespace[2pt]
RT-02.08 & La solución implementará patrones de resiliencia demostrables: reintento con retroceso exponencial y variación aleatoria, cortacircuitos, mamparos de aislamiento, límites de tasa y tiempo de espera explícito en toda llamada remota. No se admiten llamadas remotas sin tiempo de espera. & Sí & Patrones de resiliencia (límites de tasa, tiempos de espera, reintento con retroceso exponencial, mamparos, cortacircuitos) & Subdoc.4 §4.1, §4.2 \\
\addlinespace[2pt]
RT-02.09 & La solución degradará de forma elegante: ante la indisponibilidad de un componente no crítico deberá continuar operando en modo reducido, informando la degradación a la persona usuaria, y nunca fallar de forma total. & Sí & Régimen de degradación elegante técnica y restrictiva comercial, gobernado por la base tecnológica (observabilidad y degradación controlada) & Subdoc.4 §4.1 \\
\addlinespace[2pt]
RT-02.10 & Las capas de aplicación e integración escalarán horizontalmente de forma automática, con umbrales, límites superiores y costo asociado declarados en la oferta. & Sí & Escalamiento automático en AKS con umbral de expansión/contracción y límite superior por servicio & Subdoc.4 §4.1, §4.2 \\
\addlinespace[2pt]
RT-02.11 & El PROPONENTE declarará explícitamente los puntos únicos de falla que subsistan en su arquitectura y justificará por qué son aceptables. Omitir esta declaración cuando existan puntos únicos de falla se evaluará como observación grave. & Sí & Declaración de ocho puntos únicos de falla (enlace de tienda, enlace único de CD-02, ExpressRoute a DC-01, zona de Chile Central, región de Chile Central, ERP/DTE, WMS y bus de eventos) con justificación y mitigación & Subdoc.4 §4.2 (Tabla 4.11) \\
\addlinespace[2pt]
RT-02.12 & La solución admitirá su replicación a nuevas unidades, sitios, sucursales o filiales del CLIENTE sin rediseño arquitectónico, mediante parametrización o multi-tenencia. & Sí & Nodo de tienda como unidad empaquetada y estandarizada, replicable por aprovisionamiento desde imagen & Subdoc.4 §4.1 \\
\addlinespace[2pt]
RT-02.13 & El PROPONENTE presentará un modelo de dominio del negocio del caso, con las entidades principales, sus relaciones y los eventos de negocio que las modifican. & Sí & Modelo de dominio por Límite de Contexto (retail y financiero), con catálogo de eventos de negocio & Subdoc.4 §4.1, §4.2 \\
\addlinespace[2pt]
RT-02.14 & Se valorará la aplicación documentada de patrones de arquitectura evolutiva que permitan sustituir un componente sin reescribir la solución: capa anticorrupción frente a sistemas heredados, estrangulamiento progresivo y abstracción de proveedores. & Sí & Patrón Estrangulador (Strangler Fig) con capa anticorrupción frente a sistemas heredados & Subdoc.4 §4.1 \\
\addlinespace[2pt]
RT-03.01 & El PROPONENTE declarará el proveedor de nube pública, la región primaria y la región secundaria utilizadas. El proveedor deberá contar con presencia de región o zona en Chile o en Sudamérica. & Sí, condicionado & Microsoft Azure: Chile Central como región primaria y Brazil South como secundaria, ambas con tres zonas; ADR-07 condicionada a la carta del socio y a la residencia de datos & Subdoc.4 §4.1.1 (Tabla 4.6) y Tabla 4.5 \\
\addlinespace[2pt]
RT-03.02 & Todos los componentes con requisito de alta disponibilidad se desplegarán en al menos dos zonas de disponibilidad. No se aceptará un diseño en una sola zona. & Sí & Despliegue en tres zonas de disponibilidad en CLD-01 (nube pública) para borde, servicios de negocio retail, bus de eventos y bases de datos & Subdoc.4 §4.1, §4.2 \\
\addlinespace[2pt]
RT-03.03 & La totalidad de la infraestructura se definirá como código, versionada en el repositorio del CLIENTE, revisable y reproducible. No se admite infraestructura creada manualmente por consola, salvo la cuenta raíz inicial, cuya creación deberá documentarse. & Sí, condicionado & Infraestructura como código con Terraform en CLD-01, CLD-02 y SIT-01 & Subdoc.4 §4.1.1 (Tabla 4.6) \\
\addlinespace[2pt]
RT-03.04 & La red se segmentará por capas, con subredes privadas para aplicación y datos, y exposición pública restringida a la capa de borde. Ningún componente de datos será alcanzable desde Internet. & Sí & Segmentación de red en concentrador y radios (Retail, Emisor y plataforma compartida), sin publicación de SIT-01 ni de DC-01 hacia Internet; solo la capa de borde (Front Door con firewall de aplicaciones) publica hacia Internet & Subdoc.4 §4.1, §4.2 \\
\addlinespace[2pt]
RT-03.05 & El PROPONENTE privilegiará servicios administrados por sobre servicios autoadministrados cuando ello reduzca el riesgo operacional, y justificará cada excepción. & Sí & Criterio de selección N°2: ``servicio gestionado por sobre autoadministrado'', con excepciones justificadas (centro de datos DC-01) & Subdoc.4 §4.2 \\
\addlinespace[2pt]
RT-03.06 & Se aplicarán prácticas FinOps: etiquetado obligatorio de todos los recursos por ambiente, módulo y centro de costo; presupuestos con alertas de desviación; y reporte mensual de consumo desglosado entregado al CLIENTE. & Parcial & Gestión de costos de Azure contratada como servicio; el etiquetado por ambiente y ámbito, los presupuestos y las alertas de desviación se declaran en el diseño detallado & Subdoc.4 §4.2 (Tabla 4.9) \\
\addlinespace[2pt]
RT-03.07 & El PROPONENTE declarará su estrategia de reversibilidad y de mitigación del bloqueo por proveedor, identificando qué componentes son portables, cuáles no lo son y cuál sería el esfuerzo estimado de una migración. & Sí & Declaración de modelo de consumo y portabilidad por componente (Alta/Media-Alta/Total/Baja) & Subdoc.4 §4.2 \\
\addlinespace[2pt]
RT-03.08 & La solución empleará instancias reservadas, planes de ahorro o capacidad comprometida cuando el perfil de carga lo justifique, y lo reflejará en la estructura de costos. & No & Pendiente: se declara en el diseño detallado de la Etapa 1 & Subdoc.4 §4.2 \\
\addlinespace[2pt]
RT-03.09 & Se valorará el uso de cómputo sin servidor o de contenedores administrados para las cargas de perfil variable, con el análisis comparativo de costo frente a instancias permanentes. & Parcial & Contenedores administrados (Azure Kubernetes Service) para los servicios propios; el análisis de costo frente a cómputo sin servidor queda por declarar & Subdoc.4 §4.2 (Servicios contratados en Azure) \\
\addlinespace[2pt]
RT-03.10 & El componente on-premise operará de forma autónoma y degradada ante la pérdida total del enlace con la nube, durante un período mínimo de 24 horas continuas o el mayor que fije el caso. & Parcial & Nodo local de tienda con autonomía de 24 h; centro de distribución con 4 h en el Subdocumento 4, por debajo de las 24 h exigidas, por alinear & Subdoc.4 §4.2 (Alta disponibilidad, respaldos, conexiones y puntos de falla) \textbar{} Subdoc.3 §3.4.5 \\
\addlinespace[2pt]
RT-03.11 & Durante la operación desconectada, la solución continuará registrando las transacciones operacionales críticas de forma local, con integridad garantizada y sin pérdida de datos. & Sí, condicionado & Registro local en el diario del nodo de tienda (L-02) con estados pendiente, enviado, confirmado y en conflicto; motor local por decidir & Subdoc.4 §4.1 (Recorridos y consistencia) \textbar{} §4.1.1 (Tabla 4.6) \\
\addlinespace[2pt]
RT-03.12 & Restablecido el enlace, la sincronización será automática, con reconciliación determinista de conflictos, regla de resolución documentada y bitácora auditable de las decisiones aplicadas. & Sí & Drenaje idempotente por identificador de transacción local, con resolución determinista y bitácora auditable, con conciliación en 30 minutos o menos tras una desconexión (ADR-04) & Subdoc.4 §4.1, §4.2 \\
\addlinespace[2pt]
RT-03.13 & El PROPONENTE declarará qué funciones NO estarán disponibles en modo desconectado y qué procedimiento manual las suple. La ausencia de esta declaración se evaluará como observación grave. & Parcial & Sin enlace no se originan tarjetas y la compra con tarjeta propia puede declararse no disponible; la tabla de funciones no disponibles por función y su procedimiento manual se completan en el diseño detallado & Subdoc.4 §4.1 (Seguridad, continuidad y observabilidad) \textbar{} Subdoc.3 §3.4.5 \\
\addlinespace[2pt]
RT-03.14 & Los equipos on-premise críticos serán redundantes. El almacenamiento local tolerará la falla de al menos un disco; el PROPONENTE declarará el nivel RAID escogido y lo justificará frente a las alternativas. & Parcial & Discos espejados (RAID 1) en el nodo local de tienda; nivel de RAID y medio de DC-01 por declarar con el inventario del Formulario T-11 & Subdoc.4 §4.2.1 (Tabla 4.15) \textbar{} Formulario T-11 \\
\addlinespace[2pt]
RT-03.15 & Los sistemas on-premise se endurecerán conforme a los CIS Benchmarks aplicables, con gestión centralizada de parches y ventana de aplicación acordada con el CLIENTE. & No & Pendiente: se declara en el diseño detallado de la Etapa 1 & Subdoc.4 §4.2 \\
\addlinespace[2pt]
RT-03.16 & El monitoreo del componente on-premise se integrará a la misma plataforma de observabilidad que la nube, con alertamiento unificado. & Sí & Observabilidad y degradación controlada (base tecnológica), OpenTelemetry unificado en nube y nodos de tienda & Subdoc.4 §4.1, §4.2 \\
\addlinespace[2pt]
RT-03.17 & El enlace entre el sitio on-premise y la nube será redundante, con caminos físicos y proveedores distintos, y conmutación automática con tiempo de conmutación declarado. & Sí & Enlaces primario y de respaldo por sitio, de proveedores distintos; conmutación automática SD-WAN en Coyhaique & Subdoc.4 §4.2 \\
\addlinespace[2pt]
RT-03.18 & Los dispositivos de borde y de terreno se administrarán de forma remota y centralizada: inventario, configuración, actualización de firmware y de aplicación, bloqueo y borrado remoto. & Sí & Plataforma de gestión centralizada de flota, unificando dispositivos móviles, de terreno y estaciones de trabajo & Subdoc.4 §4.2 \\
\addlinespace[2pt]
RT-03.19 & Se valorará el procesamiento en el borde de las cargas que lo admitan —filtrado, agregación previa, inferencia local— reduciendo el volumen transferido y la dependencia del enlace. & No & Pendiente: se declara en el diseño detallado de la Etapa 1 & Subdoc.4 §4.2 \\
\addlinespace[2pt]
RT-03.20 & El PROPONENTE dimensionará el ancho de banda requerido por sitio, en régimen normal y en peak, y lo justificará con el cálculo de volumen de transacciones y de datos. & Sí & Dimensionamiento de ancho de banda por enlace, derivado de la volumetría de sistema & Subdoc.4 §4.2 \\
\addlinespace[2pt]
RT-03.21 & La conexión entre la red del CLIENTE y la nube se establecerá mediante enlace privado dedicado o red privada virtual con cifrado, según lo que el volumen y la criticidad justifiquen. & Sí & VPN IPsec para tiendas; enlace privado dedicado (Azure ExpressRoute) para DC-01 & Subdoc.4 §4.2 \\
\addlinespace[2pt]
RT-03.22 & El acceso remoto de las personas trabajadoras del CLIENTE, incluido el trabajo desde el hogar, se resolverá con acceso a la red de confianza cero, con verificación de postura del dispositivo. No se admite exponer servicios internos directamente a Internet. & Sí & Acceso a red de confianza cero (ZTNA) con verificación de postura del dispositivo & Subdoc.4 §4.1, §4.2 \\
\addlinespace[2pt]
RT-03.23 & La red inalámbrica de los sitios operacionales, cuando el caso la requiera, contará con segmentación por tipo de dispositivo, autenticación por certificado o credencial de empresa y cobertura verificada mediante estudio de sitio. & Sí & Punto de acceso corporativo con autenticación por certificado y red inalámbrica de clientes aislada de la red operacional & Subdoc.4 §4.2.1 \\
\addlinespace[2pt]
RT-03.24 & El PROPONENTE declarará la calidad de servicio y la priorización de tráfico aplicada a las transacciones operacionales críticas frente al tráfico administrativo. & Sí & Clasificación de tráfico por marcado DSCP (reenvío acelerado para transacción crítica) & Subdoc.4 §4.2 \\
\addlinespace[2pt]
RT-04.01 & Los cinco ambientes del numeral 4.1 estarán habilitados y operativos como condición del hito H3 del Formulario E-25. & Sí & Los cinco ambientes obligatorios (Desarrollo, QA, Preproducción, Producción, DR) declarados con su emplazamiento & Subdoc.4 §4.2 \\
\addlinespace[2pt]
RT-04.02 & Preproducción será equivalente a producción en topología, versiones de componentes y configuración. Las diferencias que subsistan por costo se declararán expresamente y se justificarán. & Sí & Topología equivalente a Producción, con diferencia declarada y justificada (separación física reducida en Preproducción) & Subdoc.4 §4.2 \\
\addlinespace[2pt]
RT-04.03 & El código residirá en un sistema de control de versiones con ramas protegidas, revisión obligatoria por pares y prohibición de escritura directa sobre la rama principal. & Sí & Rama principal protegida, revisión obligatoria de todo cambio, incluidos los de infraestructura como código & Subdoc.4 §4.2 \\
\addlinespace[2pt]
RT-04.04 & Existirá trazabilidad completa entre requerimiento, incidencia, cambio de código, prueba ejecutada y despliegue realizado. & Sí & Trazabilidad apoyada en el catálogo de servicios versionado junto con el código & Subdoc.4 §4.2 \\
\addlinespace[2pt]
RT-04.05 & El flujo de integración continua ejecutará, como mínimo: compilación, pruebas unitarias, análisis estático de código, análisis de composición de software, escaneo de secretos y escaneo de imágenes de contenedor, con criterios de bloqueo automático del despliegue. & Sí, condicionado & Cadena DevSecOps con análisis estático, dinámico, de composición y de imágenes de contenedor, y bloqueo automático (herramientas por fijar en el Subdocumento 4 y el Formulario T-11) & Subdoc.4 §4.1.1 (Tabla 4.6) \\
\addlinespace[2pt]
RT-04.06 & Los despliegues serán automatizados y reproducibles, con reversión automatizada y sin intervención manual en el paso a producción. & Parcial & Despliegue automatizado con Terraform y canal de entrega continua hacia CLD-01 y SIT-01; producción exige aprobación explícita para Retail y para el Emisor, que es una excepción a la ausencia de intervención manual & Subdoc.4 §4.2 (Ambientes) \\
\addlinespace[2pt]
RT-04.07 & La estrategia de despliegue permitirá liberar sin interrupción del servicio: azul-verde, canario o despliegue progresivo. Se declarará cuál se emplea y se demostrará en Preproducción antes de cada paso a producción. & Sí & Despliegue azul-verde en CLD-01 y SIT-01; actualización progresiva de instancias en el ámbito Emisor & Subdoc.4 §4.2 \\
\addlinespace[2pt]
RT-04.08 & Toda configuración estará externalizada del artefacto y gestionada por ambiente. Un mismo artefacto deberá poder promoverse de QA a Preproducción y a Producción sin recompilación. & Sí & Gestor de configuración centralizado por ambiente, separado del artefacto de aplicación & Subdoc.4 §4.2 \\
\addlinespace[2pt]
RT-04.09 & Los secretos residirán en un gestor de secretos con rotación automática y auditoría de acceso. Queda prohibida toda credencial embebida en código, imágenes o archivos de configuración. & Sí & Azure Key Vault (aplicación) + módulo de seguridad de hardware (HSM) para las claves del Emisor, producto y emplazamiento por definir, con rotación declarada de claves y credenciales & Subdoc.4 §4.1, §4.2 \\
\addlinespace[2pt]
RT-04.10 & Las migraciones de esquema de base de datos serán versionadas, reversibles y ejecutadas de forma automatizada, con estrategia de compatibilidad que permita convivir dos versiones de la aplicación durante el despliegue. & Sí & Migraciones de esquema coherentes con la regla de evolución de contratos por adición sin sustracción & Subdoc.4 §4.2 \\
\addlinespace[2pt]
RT-04.11 & La cobertura de pruebas automatizadas del código de lógica de negocio será de al menos 70\%, con umbral bloqueante en el flujo de integración continua. & Sí & Cobertura mínima de 70\% declarada como criterio de aceptación técnica bloqueante en CI & Subdoc.4 §4.2 \\
\addlinespace[2pt]
RT-04.12 & El PROPONENTE declarará su frecuencia de despliegue objetivo, su tiempo desde el compromiso de código hasta producción, su tasa de cambios fallidos y su tiempo de restauración, y los medirá durante la Operación. & Sí & Los cuatro indicadores DORA (frecuencia de despliegue, lead time, tasa de cambios fallidos, tiempo de restauración) & Subdoc.4 §4.2 \\
\addlinespace[2pt]
RT-04.13 & Los ambientes no productivos se apagarán o reducirán fuera del horario de uso, con el ahorro reflejado en la estructura de costos. & Sí & Desarrollo y QA reducidos fuera de la ventana 08:00–20:00; Preproducción activa solo en ventanas de prueba & Subdoc.4 §4.2 \\
\addlinespace[2pt]
RT-04.14 & Se valorará la existencia de ambientes efímeros por rama o por incidencia, creados y destruidos automáticamente. & Sí & Ambientes efímeros por rama para los servicios sin estado de CLD-01, creados/destruidos por la canalización de CI & Subdoc.4 §4.2 \\
\addlinespace[2pt]
RT-05.01 & El PROPONENTE entregará el modelo de datos documentado y un diccionario de datos con el nombre, el tipo, el dominio de valores, la obligatoriedad, el propietario y la sensibilidad de cada atributo. & Sí & Modelo de datos por dominio (entidades, dueño, módulo, naturaleza del dato) con ejemplo de atributos versionados del Precio Vigente & Subdocumento 5, secciones 5.1 (Tablas 5.1-5.3) y 5.5 \\
\addlinespace[2pt]
RT-05.02 & El PROPONENTE justificará la selección del paradigma y del motor de persistencia: relacional o no relacional, garantías transaccionales, y la posición escogida entre consistencia y disponibilidad conforme al teorema CAP, para cada dominio de datos. & Sí & Selección de paradigma/motor por entidad con postura CAP explícita (CP retail transaccional, AP vitrina, CP financiero) & Subdocumento 5, sección 5.3, Tablas 5.4 y 5.5 \\
\addlinespace[2pt]
RT-05.03 & Toda operación de negocio será trazable: la solución permitirá reconstruir quién, qué, cuándo, desde qué dispositivo y con qué valores anteriores y posteriores, para cualquier registro y en cualquier momento del período de retención. & Sí & Registro de auditoría en cada ámbito: identidad, operación, fecha y hora, dispositivo o nodo, y valor anterior y posterior & Subdoc.5 §5.2.4 \\
\addlinespace[2pt]
RT-05.04 & La calidad de datos se gestionará conforme a ISO/IEC 25012, con validación en el punto de captura, indicadores de completitud, exactitud y consistencia, y tablero de calidad disponible para el CLIENTE. & Sí & Reglas de calidad de datos (completitud ``publicable'', clasificación de merma, inmutabilidad del precio) y tablero de calidad & Subdoc.4 §4.2 \textbar{} Subdocumento 5, sección 5.8.1 \\
\addlinespace[2pt]
RT-05.05 & El almacenamiento transaccional y el analítico estarán separados. Ninguna consulta analítica podrá degradar el desempeño de la operación. & Sí & Separación OLTP/OLAP, comunicación exclusivamente asíncrona vía bus de eventos, prohibición de lecturas analíticas sobre producción & Subdocumento 5, sección 5.10 \\
\addlinespace[2pt]
RT-05.06 & La solución permitirá exportar la totalidad de la información del CLIENTE en formatos abiertos y documentados, en cualquier momento del Contrato, sin costo adicional y sin intervención del ADJUDICATARIO. & Sí & Zona de exportación en el Data Lake de cada ámbito, en CSV, JSON y Apache Parquet con su diccionario de datos & Subdoc.5 §5.2.6 \\
\addlinespace[2pt]
RT-05.07 & El PROPONENTE declarará la política de retención, archivado y eliminación por cada dominio de datos, coherente con la normativa aplicable al caso, e implementará un procedimiento verificable de eliminación segura. & Sí & Política de retención por dominio de datos, archivado en frío y eliminación segura (borrado criptográfico/físico con certificado) & Subdocumento 5, secciones 5.8.2, 5.8.3 y 5.8.4 \\
\addlinespace[2pt]
RT-05.08 & Los datos personales se tratarán conforme al Artículo 85° de las Bases Administrativas, con seudonimización o cifrado a nivel de campo para las categorías sensibles que el caso identifique. & Sí & Cifrado de campo y tokenización de identificación del cliente, cartera, comportamiento de pago y evaluación crediticia & Subdoc.4 §4.1 \\
\addlinespace[2pt]
RT-05.09 & El PROPONENTE presentará una estrategia de gestión de datos maestros que evite la duplicación de entidades compartidas entre módulos y con sistemas externos. & Sí & Principio de propiedad del dato (Data Ownership) y Servicio de oferta comercial como dueño único de la entidad Catálogo & Subdocumento 5, sección 5.1 \\
\addlinespace[2pt]
RT-05.10 & Se valorará la implementación de un catálogo de datos con linaje automatizado, que permita rastrear el origen de cada indicador de negocio hasta su fuente. & Sí & Catálogo de datos con linaje automatizado (origen, transformaciones en Azure Data Lake Storage Gen2, motor analítico por decidir, consumo por tablero/modelo) & Subdocumento 5, sección 5.11.3 \\
\addlinespace[2pt]
RT-05.11 & El PROPONENTE presentará un plan de migración con alcance, origen, volumen, reglas de transformación, criterios de calidad, estrategia de ejecución y plan de reversión. & Sí & Estrategia de migración en cuatro fases (extracción, saneamiento, coexistencia/conciliación, corte) con circuit breaker y rollback & Subdocumento 5, sección 5.7 \\
\addlinespace[2pt]
RT-05.12 & La migración incluirá una etapa de perfilado y saneamiento previo, con informe de los defectos detectados en los datos de origen y la decisión adoptada sobre cada uno. & Sí & Fase 2 - Saneamiento: reconciliación de stock, deduplicación de SKU, reconstrucción de historial de precio, normalización de cartera & Subdocumento 5, sección 5.7 (Fase 2) \\
\addlinespace[2pt]
RT-05.13 & Se ejecutarán al menos dos ensayos completos de migración sobre Preproducción antes de la migración definitiva, con medición del tiempo total y del resultado de la conciliación. & Sí & Dos ensayos completos de migración en Preproducción antes de cada corte definitivo, con medición del tiempo total y conciliación & Subdoc.5 §5.3.4 (Tabla 5.14) \textbar{} Formulario T-18 \\
\addlinespace[2pt]
RT-05.14 & La conciliación posterior a la migración será cuantitativa y verificable: recuentos, sumas de control y muestreo dirigido. Toda diferencia deberá quedar explicada. & Sí & Conteos de control (row counts), checksums monetarias/de stock y conciliación cruzada diaria con circuit breaker & Subdocumento 5, sección 5.7 (Fase 3) \\
\addlinespace[2pt]
RT-05.15 & Los datos históricos que no se migren quedarán accesibles en un repositorio de consulta durante el período de retención que fije el caso. & Sí & Repositorio de consulta de solo lectura en el Data Lake de cada ámbito para los datos no migrados, con su plazo de conservación & Subdoc.5 §5.3.5 (Tabla 5.16) \\
\addlinespace[2pt]
RT-05.16 & Los servicios síncronos se documentarán en OpenAPI 3.1 y los flujos dirigidos por eventos en AsyncAPI 2.6 o superior. La documentación se generará desde el código y se mantendrá actualizada automáticamente. & Sí & Catálogo central de contratos versionado junto con el código, en OpenAPI 3.1 y AsyncAPI 2.6 & Subdoc.4 §4.1 \\
\addlinespace[2pt]
RT-05.17 & Los contratos de interfaz se versionarán semánticamente, con compatibilidad hacia atrás y política de obsolescencia con preaviso mínimo de seis meses. & Sí & Reglas de evolución de contratos (adición sin sustracción); versión antigua se retira solo con cero consumidores registrados & Subdoc.4 §4.1, §4.2 \\
\addlinespace[2pt]
RT-05.18 & La autenticación entre sistemas empleará OAuth 2.1 con credenciales de cliente o autenticación mutua TLS. Queda prohibida la autenticación por clave estática en la ruta de la dirección web. & Parcial & Identidades de servicio y credenciales de vida corta con Microsoft Entra ID (OIDC y OAuth); la versión 2.1 y la autenticación mutua TLS se confirman en el diseño detallado & Subdoc.4 §4.1 (Seguridad, continuidad y observabilidad) \\
\addlinespace[2pt]
RT-05.19 & Toda integración registrará la transacción de entrada y de salida, con identificador de correlación común que permita seguir una operación de negocio a través de todos los sistemas involucrados. & Sí & Referencia de correlación en cada hecho publicado y en el registro del Servicio de control de cruces & Subdoc.4 §4.1 (Recorridos y consistencia) \textbar{} Subdoc.5 §5.2.1 \\
\addlinespace[2pt]
RT-05.20 & Las integraciones con sistemas heredados o de terceros se aislarán mediante una capa anticorrupción, de modo que un cambio en el sistema externo no propague su modelo al núcleo de la solución. & Sí & Capa anticorrupción frente a cada sistema heredado que se conserva & Subdoc.4 §4.1 \\
\addlinespace[2pt]
RT-05.21 & El PROPONENTE declarará, por cada integración, el modo (síncrono o asíncrono), el volumen esperado, la ventana de disponibilidad del sistema contraparte y el comportamiento de la solución cuando ese sistema no responde. & Sí & Catálogo de interfaces internas y externas con modo, volumen, ventana y comportamiento ante no respuesta & Subdoc.4 §4.1, §4.2 \\
\addlinespace[2pt]
RT-05.22 & La solución soportará la carga y descarga masiva de información en formatos abiertos, con validación previa, informe de errores por registro y procesamiento parcial. & Parcial & Cargas masivas del saneamiento y de la migración con validación previa y conciliación; la carga y descarga masiva para el usuario final queda por declarar & Subdoc.5 §5.3.2 y §5.3.4 \\
\addlinespace[2pt]
RT-05.23 & Se emplearán los estándares sectoriales de intercambio que las Bases Técnicas del caso identifiquen. & Sí & Estándares sectoriales de intercambio: GS1 GTIN y GLN, EANCOM, EPCIS 2.0, OpenAPI 3.1, AsyncAPI y XML de documentos tributarios del SII & Subdoc.5 §5.2.6 (Tabla 5.10) \\
\addlinespace[2pt]
RT-05.24 & Se valorará la publicación de un portal de servicios para desarrolladores con documentación navegable, ambiente de pruebas y credenciales de prueba autoservidas. & No & Pendiente: se declara en el diseño detallado de la Etapa 1 & Subdoc.5 \\
\addlinespace[2pt]
RT-05.25 & La solución proveerá una capa analítica con tableros operacionales y de gestión, construidos sobre los indicadores que las Bases Técnicas del caso definan. & Sí & Capa de consumo (gold) con modelos dimensionales por dominio expuestos a tableros gerenciales & Subdocumento 5, sección 5.11.1 \\
\addlinespace[2pt]
RT-05.26 & Los tableros permitirán filtrar por período, unidad organizacional y dimensiones propias del caso, y profundizar desde el indicador agregado hasta la transacción de origen. & Parcial & Tableros sobre la capa de consumo de cada ámbito, con recorrido hasta el evento de origen; filtros por período y unidad por declarar & Subdoc.5 §5.2.2 \\
\addlinespace[2pt]
RT-05.27 & El CLIENTE podrá construir sus propios informes sin intervención del ADJUDICATARIO, mediante una herramienta de autoservicio con modelo semántico documentado. & Sí & BI de autoservicio sobre la capa de consumo (gold), con control de acceso basado en roles & Subdocumento 5, sección 5.11.1 \\
\addlinespace[2pt]
RT-05.28 & Todo informe será exportable en formatos abiertos y programable para envío automático por calendario. & Parcial & Informes sobre la capa de consumo de cada ámbito; exportación en formatos abiertos y envío programado por declarar & Subdoc.5 §5.2.2 \\
\addlinespace[2pt]
RT-05.29 & La latencia máxima entre la ocurrencia de una transacción y su disponibilidad en la capa analítica será la que fije el caso y, en su defecto, no superará las 4 horas. & Sí & Metas de latencia: disponibilidad publicada en 30 segundos, precio en 5 minutos, estado de pedido en tiempo real y posición financiera antes de las 07:00 & Subdoc.5 §5.2.2 (Tabla 5.6) \textbar{} Subdoc.4 §4.1 (Tabla 4.5, ADR-03) \\
\addlinespace[2pt]
RT-05.30 & Se valorará la incorporación de analítica predictiva pertinente al proceso del caso, con el modelo, sus variables, su métrica de desempeño y su plan de reentrenamiento documentados. & Parcial & Modelos predictivos de conteo y de clasificación de merma, condicionados a su validación; variables, métrica y plan de reentrenamiento en las fichas del Formulario T-19 & Subdoc.4 §4.1 \textbar{} Subdoc.13 §13.2 y Formulario T-19 \\
\addlinespace[2pt]
RT-06.01 & El espacio asignado será de uso exclusivo de la solución y estará aislado de otras dependencias del CLIENTE, con acceso independiente. & Sí & Recinto de uso exclusivo, aislado, con acceso independiente (tipología Sala técnica principal) & Subdoc.4 §4.2 \\
\addlinespace[2pt]
RT-06.02 & Los muros no estructurales del recinto contarán con blindaje perimetral; el PROPONENTE especificará el material y la resistencia. & Sí & Blindaje perimetral en muros no estructurales de la zona financiera y de la zona neutral & Subdoc.4 §4.2 \\
\addlinespace[2pt]
RT-06.03 & El PROPONENTE entregará el plano de distribución interna del recinto, con la separación de las zonas de generadores, baterías, climatización, servidores, comunicaciones, trabajo y respaldo. & Sí & Plano de distribución interna con cinco zonas declaradas & Subdoc.4 §4.2 \\
\addlinespace[2pt]
RT-06.04 & El piso técnico, la canalización, el cableado estructurado y el etiquetado se ejecutarán conforme a norma, con documentación de la certificación de cada enlace. & Sí & Piso técnico, canalización y cableado estructurado ejecutados y certificados conforme a norma & Subdoc.4 §4.2 \\
\addlinespace[2pt]
RT-06.05 & Los racks de servidores serán independientes de los racks de equipos de comunicación. Se declarará la ocupación proyectada de cada rack y su margen de crecimiento. & Parcial & Racks de servidores independientes de los de comunicaciones; ocupación proyectada y margen de crecimiento por declarar con el inventario & Subdoc.4 §4.3.1 \\
\addlinespace[2pt]
RT-06.06 & La obra civil de separación de las instalaciones es de cargo del CLIENTE; su especificación técnica y su coordinación son de cargo del PROPONENTE. & Sí & Especificación técnica y coordinación de la obra civil de separación de zonas a cargo del PROPONENTE & Subdoc.4 §4.2 \\
\addlinespace[2pt]
RT-06.07 & El suministro eléctrico de los equipos será ininterrumpido, con sistema de alimentación ininterrumpida dimensionado para una autonomía mínima de 30 minutos a plena carga. & Parcial & UPS con al menos 30 minutos de autonomía a plena carga; marca y modelo por definir en el Formulario T-11 & Subdoc.4 §4.3.1 \textbar{} Formulario T-11 \\
\addlinespace[2pt]
RT-06.08 & La capacidad de generación autónoma asegurará un rango mínimo de 24 horas continuas de operación, con estanque de combustible dimensionado y contrato de reabastecimiento declarado. & Sí & Grupo electrógeno con 24 h de autonomía, estanque dimensionado y contrato de reabastecimiento & Subdoc.4 §4.2 \\
\addlinespace[2pt]
RT-06.09 & La instalación eléctrica del recinto será independiente de la del resto del edificio y cumplirá la normativa eléctrica chilena vigente, incluida la NCh Elec. 2777 sobre sistemas de puesta a tierra. & Sí & Acometida propia del recinto, independiente del resto del edificio, conforme a NCh Elec. 2777 & Subdoc.4 §4.2 \\
\addlinespace[2pt]
RT-06.10 & Se efectuará revisión y medición semestral de las instalaciones eléctricas del recinto, con informe entregable al CLIENTE. & Sí & Revisión semestral programada en el calendario de mantención & Subdoc.4 §4.2 \\
\addlinespace[2pt]
RT-06.11 & El PROPONENTE declarará la carga eléctrica proyectada en kW, el factor de potencia y la eficiencia en el uso de la energía (PUE) estimada del recinto. & Parcial & Carga eléctrica y PUE del recinto por declarar con el inventario y el informe del recinto & Subdoc.4 §4.3.1 \\
\addlinespace[2pt]
RT-06.12 & Se valorará la redundancia de alimentación en configuración 2N o N+1 con doble acometida y transferencia automática. & Sí & Doble acometida con transferencia automática para la zona financiera & Subdoc.4 §4.2 \\
\addlinespace[2pt]
RT-06.13 & El recinto contará con climatización de precisión para operación continua, redundante en configuración N+1, con control de temperatura y de humedad relativa dentro de los rangos que recomienda el fabricante del equipamiento. & Sí & Climatización de precisión N+1 (referenciada en 4.2.4.4) & Subdoc.4 §4.2 \\
\addlinespace[2pt]
RT-06.14 & Se monitorearán en línea la temperatura, la humedad y la presencia de agua, con alertamiento integrado a la plataforma de observabilidad. & Parcial & Monitoreo del recinto integrado a la observabilidad (OpenTelemetry y Azure Monitor); sensores de temperatura, humedad y agua por especificar en el Formulario T-11 & Subdoc.4 §4.3.1 \textbar{} Formulario T-11 \\
\addlinespace[2pt]
RT-06.15 & El PROPONENTE declarará la estrategia de contención de pasillo frío o caliente y su efecto en la eficiencia energética. & Sí & Contención de pasillo frío/caliente referenciada en el cálculo de PUE & Subdoc.4 §4.2 \\
\addlinespace[2pt]
RT-06.16 & El recinto contará con detección temprana por aspiración de aire con tecnología láser, tipo AnaLASER o equivalente de prestaciones iguales o superiores. & Sí & Detección temprana por aspiración de aire (referenciada en 4.2.4.5) & Subdoc.4 §4.2 \\
\addlinespace[2pt]
RT-06.17 & La extinción será automática mediante agente limpio tipo FM-200 o equivalente, con aprobación UL e instalación conforme a norma NFPA. & Sí & Extinción automática por agente limpio (referenciada en 4.2.4.5) & Subdoc.4 §4.2 \\
\addlinespace[2pt]
RT-06.18 & Se proveerá un sistema secundario de extintores portátiles habilitados, con mantención y certificación vigentes. & Sí & Sistema secundario de extintores portátiles (referenciado en 4.2.4.5) & Subdoc.4 §4.2 \\
\addlinespace[2pt]
RT-06.19 & El sistema de detección y extinción se integrará al monitoreo en línea y notificará al NOC y a la contraparte del CLIENTE. & Sí & Integración del sistema de detección y extinción al monitoreo en línea & Subdoc.4 §4.2 \\
\addlinespace[2pt]
RT-06.20 & El ingreso al recinto se controlará mediante seguridad física y control de acceso biométrico basado principalmente en biometría facial, con AFIS como respaldo. Se admite proponer sistemas de mayor seguridad. & Sí & Control de acceso biométrico independiente por zona (referenciado en 4.2.4.2 y 4.2.4.6) & Subdoc.4 §4.2 \\
\addlinespace[2pt]
RT-06.21 & Todo ingreso y egreso quedará registrado en una bitácora auditable, con identificación de la persona, fecha, hora y motivo, conservada por el período de retención declarado. & Sí & Bitácora auditable de ingreso/egreso (referenciada en 4.2.4.6) & Subdoc.4 §4.2 \\
\addlinespace[2pt]
RT-06.22 & Entre el acceso principal y el término del pasillo de la zona de control se dispondrá un espacio para la atención de personas en proceso de enrolamiento. Se evaluará mejor la existencia de una estación de enrolamiento fuera de las instalaciones del recinto técnico. & Sí & Espacio de enrolamiento (referenciado en 4.2.4.6) & Subdoc.4 §4.2 \\
\addlinespace[2pt]
RT-06.23 & Al término del pasillo se instalará un acceso que impida el paso de más de una persona a la vez, con nueva verificación de identidad previa al ingreso. & Sí & Control de esclusa de una persona con verificación de identidad (referenciado en 4.2.4.6) & Subdoc.4 §4.2 \\
\addlinespace[2pt]
RT-06.24 & El recinto contará con videovigilancia y monitoreo IP, con imágenes en línea y disponibles para visualización de al menos los últimos 30 días. Las grabaciones anteriores se respaldarán en un medio secundario recuperable y auditable. & Sí & Videovigilancia y monitoreo IP con retención de 30 días (referenciado en 4.2.4.6) & Subdoc.4 §4.2 \\
\addlinespace[2pt]
RT-06.25 & El PROPONENTE declarará el procedimiento de acceso de terceros —fabricantes, mantenedores, auditores— con acompañamiento obligatorio y registro. & Sí & Procedimiento de acceso de terceros con acompañamiento y registro (referenciado en 4.2.4.6) & Subdoc.4 §4.2 \\
\addlinespace[2pt]
RT-06.26 & Se habilitará un servicio de custodia de medios de respaldo para el sitio primario, en un medio físico transportable a otro lugar cuando el CLIENTE lo determine. Se admite proponer una solución más segura y eficiente, debidamente justificada. & Sí & Servicio de custodia de medios de respaldo en formato físico transportable & Subdoc.4 §4.2 \\
\addlinespace[2pt]
RT-06.27 & El recinto de custodia cumplirá exigencias de luminosidad, humedad, ventilación y cualquier otro factor que pueda afectar la calidad y la disponibilidad de los medios. & Sí & Condiciones de luminosidad, humedad y ventilación adecuadas para la custodia & Subdoc.4 §4.2 \\
\addlinespace[2pt]
RT-06.28 & Se llevará un inventario de medios con rotación, verificación periódica de legibilidad y registro de todo movimiento de entrada y de salida. & Sí & Inventario de medios con rotación y verificación periódica de legibilidad & Subdoc.4 §4.2 \\
\addlinespace[2pt]
RT-06.29 & El PROPONENTE habilitará el espacio físico necesario para el personal encargado de la operación y administración de la plataforma, con estaciones de trabajo, telefonía, conexión a Internet y todo elemento que permita realizar la labor en condiciones adecuadas. & Sí & Espacio de operación con estaciones de trabajo, telefonía y conectividad, ≈20 m² & Subdoc.4 §4.2 \\
\addlinespace[2pt]
RT-06.30 & El espacio de operación estará separado de la sala de equipos y no requerirá el ingreso al recinto técnico para las labores habituales de operación. & Sí & Espacio de operación separado de la sala de equipos, sin necesidad de ingreso al recinto técnico & Subdoc.4 §4.2 \\
\addlinespace[2pt]
RT-06.31 & Las instalaciones sanitarias, las zonas de seguridad ante emergencia y las áreas exteriores existentes en el edificio del CLIENTE podrán utilizarse y no deben implementarse nuevamente. El PROPONENTE declarará de cuáles hará uso. & Sí & Uso de instalaciones sanitarias y zonas de emergencia ya existentes en el edificio de casa matriz & Subdoc.4 §4.2 \\
\addlinespace[2pt]
RT-06.32 & El acceso a las redes de comunicaciones estará provisto a través de rutas físicas distintas, con ingreso al edificio por puntos separados. & Sí & Rutas físicas distintas con ingreso al edificio por puntos separados & Subdoc.4 §4.2 \\
\addlinespace[2pt]
RT-06.33 & El PROPONENTE proveerá toda la conectividad, la seguridad y las canalizaciones, cañerías o ductos requeridos para cumplir los niveles de servicio comprometidos. & Sí & Provisión íntegra de conectividad, seguridad y canalizaciones & Subdoc.4 §4.2 \\
\addlinespace[2pt]
RT-06.34 & Se privilegiará al PROPONENTE que ofrezca o provea especificaciones nuevas o mejores que las aquí establecidas, debidamente fundamentadas. & No & Pendiente: se declara en el inventario del Formulario T-11 y el diseño detallado & Subdoc.4 §4.3 \textbar{} Formulario T-11 \\
\addlinespace[2pt]
RT-07.01 & El PROPONENTE declarará la modalidad escogida —activo-activo o activo-pasivo— y la justificará frente al costo, al RTO comprometido y a la complejidad operacional que introduce. & Sí & Modalidad activo-pasivo, justificada frente a costo, RTO y complejidad operacional & Subdoc.4 §4.2 \\
\addlinespace[2pt]
RT-07.02 & El sitio secundario estará emplazado a una distancia suficiente del principal para no verse afectado por el mismo evento de fuerza mayor. El PROPONENTE declarará la distancia y el análisis de amenazas comunes considerado. & Parcial & Recuperación en Azure Brazil South (CLD-02), en otro país y sin amenazas locales comunes con Chile Central; distancia y análisis formal antes de aprobar el diseño detallado & Subdoc.4 §4.3.1 y §4.3.2 \\
\addlinespace[2pt]
RT-07.03 & La replicación de datos será continua, con medición y alertamiento del retraso de replicación. & Sí & Replicación continua y asíncrona (réplica de Azure Database for PostgreSQL hacia CLD-02 y geo-replicación de Event Hubs) con medición y alertamiento de rezago & Subdoc.4 §4.2 \\
\addlinespace[2pt]
RT-07.04 & El objetivo de tiempo de recuperación (RTO) no superará 4 horas y el objetivo de punto de recuperación (RPO) no superará 15 minutos para los servicios críticos, salvo exigencia superior del caso. & Sí, condicionado & RTO ≤ 4 h y RPO ≤ 15 min medidos y demostrados en ensayo; copia de las claves del Emisor en CLD-02, sujeta a la disponibilidad regional y a la residencia que apruebe el Cliente & Subdoc.4 §4.2 (Tabla 4.12) \textbar{} §4.3.2 \\
\addlinespace[2pt]
RT-07.05 & El procedimiento de conmutación estará documentado, automatizado en la mayor medida posible y ejecutable por el personal del CLIENTE tras la transferencia de conocimiento. & Sí & Procedimiento de conmutación en seis pasos, con automatización de los pasos técnicos & Subdoc.4 §4.3.2 \\
\addlinespace[2pt]
RT-07.06 & Existirá un procedimiento de retorno al sitio principal igualmente documentado y probado, con reconciliación de los datos generados durante la contingencia. & Sí & Procedimiento de retorno a CLD-01 con reconciliación por drenaje idempotente & Subdoc.4 §4.2 \\
\addlinespace[2pt]
RT-07.07 & El plan de recuperación ante desastres se probará al menos dos veces al año mediante conmutación real, con informe de resultados, medición del RTO y del RPO efectivamente alcanzados y plan de corrección de las brechas detectadas. & Sí & Conmutación real dos veces al año, con medición de RTO/RPO efectivos e informe de brechas & Subdoc.4 §4.2 \\
\addlinespace[2pt]
RT-07.08 & Se valorará que la conmutación sea automática ante la detección de indisponibilidad, con criterio de disparo declarado y protección contra conmutación innecesaria. & Parcial & Conmutación semiautomática (detección y ejecución técnica automatizadas, decisión de disparo humana); la protección contra doble promoción se declara en el diseño detallado & Subdoc.4 §4.2 (Alta disponibilidad, respaldos, conexiones y puntos de falla) \textbar{} §4.3.2 \\
\addlinespace[2pt]
RT-07.09 & La política de respaldo seguirá el esquema 3-2-1-1-0: tres copias, en dos medios distintos, una fuera de sitio, una inmutable o fuera de línea y cero errores de verificación de restauración. & Sí & Política de respaldo 3-2-1-1-0 declarada por dominio de datos & Subdoc.4 §4.2 \\
\addlinespace[2pt]
RT-07.10 & Los respaldos estarán cifrados en reposo y en tránsito, con clave gestionada de forma independiente de la infraestructura respaldada. & Sí & Respaldos cifrados con clave gestionada de forma independiente (Azure Key Vault / módulo de seguridad de hardware del Emisor) & Subdoc.4 §4.2 \\
\addlinespace[2pt]
RT-07.11 & Las copias inmutables estarán protegidas contra borrado y contra modificación durante su período de retención, incluso frente a credenciales administrativas comprometidas. & Sí & Snapshot con retención bloqueada (object lock); política WORM inmutable por diseño & Subdoc.4 §4.2 \\
\addlinespace[2pt]
RT-07.12 & Se ejecutará y documentará una prueba de restauración al menos mensual, sobre una muestra representativa, con medición del tiempo efectivo de restauración. & Sí & Prueba de restauración mensual sobre muestra representativa de cada dominio & Subdoc.4 §4.2 \\
\addlinespace[2pt]
RT-07.13 & El PROPONENTE declarará, por cada dominio de datos, la frecuencia de respaldo, el período de retención y el tiempo estimado de restauración completa. & Sí & Política de respaldo declarada por dominio de datos (frecuencia, retención, tiempo de restauración) & Subdoc.4 §4.2 \\
\addlinespace[2pt]
RT-07.14 & Los respaldos permitirán la restauración granular: un registro, una tabla, un módulo o el sistema completo. & Sí & Restauración granular nativa en los cuatro dominios (respaldo lógico y snapshots) & Subdoc.4 §4.2 \\
\addlinespace[2pt]
RT-08.01 & El PROPONENTE especificará el equipamiento de cómputo con su marca, modelo de referencia, procesador, memoria, almacenamiento local, interfaces y consumo, junto con el cálculo de dimensionamiento que lo sustenta. & Parcial & Servidores de DC-01 con dimensionamiento del apartado Dimensionamiento y plan de capacidad; marca y modelo por definir en el Formulario T-11 & Subdoc.4 §4.2.1 (Tabla 4.15) \textbar{} Formulario T-11 \\
\addlinespace[2pt]
RT-08.02 & El almacenamiento será redundante, con tolerancia declarada a la falla de discos, control de errores y monitoreo predictivo de salud de los medios. & Sí & Arreglo NVMe RAID 10 / erasure coding con monitoreo predictivo S.M.A.R.T. / NVMe health telemetry & Subdoc.4 §4.2 \\
\addlinespace[2pt]
RT-08.03 & Los conmutadores de núcleo, los cortafuegos y los balanceadores de carga estarán en configuración de alta disponibilidad, sin punto único de falla. & Sí & Conmutadores, cortafuegos y balanceadores clase empresarial en configuración de alta disponibilidad activo-pasivo o clúster & Subdoc.4 §4.2 \\
\addlinespace[2pt]
RT-08.04 & Todo el equipamiento contará con fuentes de poder redundantes y conexión a circuitos eléctricos distintos. & Sí & Fuente de poder dual por servidor, cada una a un circuito eléctrico distinto & Subdoc.4 §4.2 \\
\addlinespace[2pt]
RT-08.05 & El PROPONENTE declarará el margen de crecimiento del dimensionamiento propuesto, expresado como porcentaje sobre la carga proyectada del caso, y el procedimiento de ampliación. & Parcial & Margen de crecimiento y procedimiento de ampliación por declarar con el inventario del Formulario T-11 & Subdoc.4 §4.2.1 \textbar{} Formulario T-11 \\
\addlinespace[2pt]
RT-08.06 & El equipamiento será nuevo, sin uso previo, con garantía de fábrica vigente desde la recepción conforme. & Sí & Equipamiento nuevo, sin uso previo, con garantía de fábrica vigente & Subdoc.4 §4.2 \\
\addlinespace[2pt]
RT-08.07 & El PROPONENTE especificará las estaciones de trabajo requeridas para la operación y la administración de la plataforma, en la cantidad que determine el dimensionamiento del caso, con monitores duales. & Sí & Estación de trabajo de clase empresarial con dos monitores, dimensionada sobre 46 puestos & Subdoc.4 §4.2 \\
\addlinespace[2pt]
RT-08.08 & Los puestos de trabajo cumplirán las condiciones ergonómicas de la NCh 2527 y los equipos contarán con certificación de eficiencia energética. & Sí & Mobiliario y equipos conformes a NCh 2527, con certificación ENERGY STAR o equivalente & Subdoc.4 §4.2 \\
\addlinespace[2pt]
RT-08.09 & Las estaciones estarán gestionadas de forma centralizada, con cifrado de disco, control de dispositivos extraíbles, antivirus con detección y respuesta y actualización automatizada. & Sí & Gestión centralizada con cifrado de disco completo, control de extraíbles, EDR y actualización automatizada & Subdoc.4 §4.2 \\
\addlinespace[2pt]
RT-08.10 & El PROPONENTE especificará cada dispositivo de terreno con marca, modelo de referencia, cantidad, características mínimas, accesorios, consumibles y costo unitario estimado, aun cuando su compra sea de cargo del CLIENTE. & Sí & Especificación completa de dispositivos de terreno (terminal móvil, lector, terminal de pago, impresora, lector de CD), costo remitido a la Oferta Económica & Subdoc.4 §4.2 \\
\addlinespace[2pt]
RT-08.11 & La especificación considerará las condiciones reales de uso del caso: intemperie, humedad, polvo, vibración, temperatura, uso con guantes, luminosidad y autonomía de batería requerida por turno. & Sí & Condiciones reales de uso declaradas por dispositivo (IP54, resistencia a caída, uso con guantes, batería de repuesto) & Subdoc.4 §4.2 \\
\addlinespace[2pt]
RT-08.12 & Los dispositivos declararán su grado de protección contra polvo y agua y su resistencia a caídas, coherentes con el entorno de operación. & Sí & Grado de protección IP54 y resistencia a caída desde 1,2 m declarados & Subdoc.4 §4.2 \\
\addlinespace[2pt]
RT-08.13 & El PROPONENTE indicará el ciclo de vida esperado de cada dispositivo, la disponibilidad de repuestos y el plan de reposición durante los 56 meses del Contrato. & Sí & Ciclo de vida (4-5 años dispositivos móviles, 5-7 servidores) con renovación proyectada al mes 40 & Subdoc.4 §4.2 \\
\addlinespace[2pt]
RT-08.14 & Los dispositivos se integrarán a la gestión centralizada de flota exigida en RT-03.18. & Sí & Incorporación de todos los dispositivos móviles y de terreno a la plataforma de gestión centralizada de flota & Subdoc.4 §4.2 \\
\addlinespace[2pt]
RT-08.15 & El PROPONENTE proveerá una unidad de cada tipo de dispositivo especificado para pruebas de aceptación por parte del CLIENTE, antes de la compra masiva. & Sí & Provisión de una unidad de cada tipo de dispositivo de terreno para prueba de aceptación & Subdoc.4 §4.2 \\
\addlinespace[2pt]
RT-08.16 & El PROPONENTE presentará el plan de ciclo de vida del equipamiento: recepción, puesta en servicio, mantención, actualización, retiro y disposición final. & Sí & Plan de ciclo de vida completo (recepción, puesta en servicio, mantención, actualización, retiro, disposición) & Subdoc.4 §4.2 \\
\addlinespace[2pt]
RT-08.17 & Todo medio de almacenamiento que salga de servicio será borrado de forma segura y verificable, con certificado de destrucción o de sanitización entregado al CLIENTE. & Sí & Borrado seguro conforme a NIST SP 800-88, con certificado de destrucción/sanitización & Subdoc.4 §4.2 \\
\addlinespace[2pt]
RT-08.18 & La disposición final de equipamiento electrónico se realizará con gestor autorizado, conforme a la normativa de residuos aplicable, con certificado de disposición. & Sí & Disposición final con gestor autorizado conforme a la Ley N° 20.920, con certificado por lote & Subdoc.4 §4.2 \\
\addlinespace[2pt]
RT-08.19 & Se valorará una estrategia de reacondicionamiento o de extensión de vida útil que reduzca el impacto ambiental, cuantificada en la propuesta. & Sí & Reacondicionamiento certificado del fabricante para el stock de repuestos de terminales móviles y lectores & Subdoc.4 §4.2 \\
\addlinespace[2pt]
RT-09.01 & El PROPONENTE presentará el cálculo de capacidad que sustenta su dimensionamiento, con los supuestos de usuarios concurrentes, transacciones por segundo, volumen de datos y crecimiento anual, tomados de la volumetría del caso. & Sí & Derivación de la volumetría de sistema (concurrencia, TPS, almacenamiento) a partir de la volumetría operacional del caso & Subdocumento 4, sección 4.2 (memoria de capacidad) \\
\addlinespace[2pt]
RT-09.02 & La solución soportará la concurrencia y el volumen de transacciones que fije el caso, y mantendrá los umbrales del numeral 9.1 bajo esa carga. & Sí & Dimensionamiento de cómputo, bus de eventos y motores de datos ajustado a los umbrales de latencia del caso & Subdoc.4 §4.2 \\
\addlinespace[2pt]
RT-09.03 & La solución soportará, sin rediseño, un crecimiento de al menos tres veces la volumetría inicial del caso en un horizonte de tres años. & Parcial & Dimensionamiento sobre la proyección a tres años del Caso (430 líneas de caja y pedidos del evento anual proyectados); la capacidad de triplicar la volumetría sin rediseño no se declara aún & Subdoc.4 §4.2 (Tabla 4.14) \\
\addlinespace[2pt]
RT-09.04 & El escalamiento de las capas de aplicación e integración será horizontal y automático, con tiempo de reacción declarado y sin pérdida de transacciones en curso. & Sí & Escalamiento automático por criterio de saturación (latencia p95, profundidad de cola) con umbrales por servicio & Subdoc.4 §4.1 \\
\addlinespace[2pt]
RT-09.05 & EL PROPONENTE identificará el componente que primero se convertirá en cuello de botella al crecer la carga y explicará cómo lo detectará y cómo lo resolverá. & Parcial & Enlace de respaldo de las tiendas que dependen de la red de un centro comercial y enlace único del CD-02 como primeros cuellos de botella & Subdoc.4 §4.2 (Tablas 4.11 y 4.14) \\
\addlinespace[2pt]
RT-09.06 & Se ejecutarán pruebas de carga sobre Preproducción con un volumen equivalente a 1,5 veces el peak declarado, y pruebas de estrés hasta identificar el punto de quiebre de la solución. & Sí & Pruebas de carga a 1,5x el peak proyectado, para verificar los supuestos de volumetría & Subdocumento 4, sección 4.2 (prueba de carga en Preproducción) \\
\addlinespace[2pt]
RT-09.07 & El informe de pruebas de carga incluirá la curva de tiempo de respuesta frente a carga, el punto de saturación, el consumo de recursos y el comportamiento durante y después del peak. & Parcial & Plan de pruebas de carga a 1,5 veces el peak; el informe con la curva de respuesta, el punto de saturación y el consumo se entrega con cada certificación & Formulario T-13 §1 \textbar{} Subdoc.9 \\
\addlinespace[2pt]
RT-09.08 & La solución degradará de forma controlada al superarse la capacidad: encolamiento, limitación de tasa y mensaje explícito a la persona usuaria, nunca error genérico ni pérdida silenciosa de transacciones. & Sí & Degradación controlada comercial (base tecnológica, observabilidad y degradación controlada): priorización de checkout sobre navegación, throttling de sesiones nuevas & Subdoc.4 §4.1, §4.2 \\
\addlinespace[2pt]
RT-09.09 & Se gestionará la capacidad durante la Operación con proyección trimestral de crecimiento, alertas anticipadas de agotamiento y propuesta de ajuste de dimensionamiento y de costo. & Sí & Curva de crecimiento de almacenamiento a 3 años, reevaluada en cada revisión trimestral de capacidad & Subdoc.4 §4.2 \\
\addlinespace[2pt]
RT-09.10 & Se valorará la existencia de pruebas de carga automatizadas ejecutadas de forma periódica en el flujo de integración continua, con detección de regresiones de desempeño. & No & Pendiente: se declara en el plan de capacidad y el plan de calidad & Subdoc.4 §4.2 \textbar{} Subdoc.9 \\
\addlinespace[2pt]
RT-10.01 & La solución alcanzará una disponibilidad mensual mínima de 99,9\% para los servicios clasificados como críticos, medida sobre la transacción de negocio de extremo a extremo. & Sí, condicionado & Diseño en tres zonas de disponibilidad con recuperación regional; la disponibilidad de 99,9 \% se demuestra con la medición extremo a extremo y no con el SLA de los servicios & Subdoc.4 §4.3.1 \\
\addlinespace[2pt]
RT-10.02 & El PROPONENTE clasificará cada servicio de la solución en crítico, alto, medio o bajo, justificando la clasificación con el impacto operacional de su indisponibilidad, y aplicará a cada uno el nivel de servicio correspondiente del Artículo 78° de las Bases Administrativas. & Sí & Clasificación de módulos en crítico/alto/medio/bajo con justificación de impacto & Subdoc.4 §4.2 \\
\addlinespace[2pt]
RT-10.03 & El plan de continuidad del negocio se elaborará conforme a ISO 22301, con análisis de impacto en el negocio, escenarios de contingencia, procedimientos manuales de respaldo y criterios de activación. & Sí & Plan de continuidad del negocio con tres escenarios de contingencia, activación y procedimiento manual de respaldo & Subdoc.4 §4.2 \\
\addlinespace[2pt]
RT-10.04 & La continuidad TIC se estructurará conforme a ISO/IEC 27031, articulada con el plan de recuperación ante desastres del Capítulo 7. & Sí & Continuidad TIC (ISO/IEC 27031) articulada con el plan de recuperación de sitio de la sección 4.2.5 & Subdoc.4 §4.2 \\
\addlinespace[2pt]
RT-10.05 & Los mantenimientos programados se ejecutarán fuera de la ventana operacional crítica que defina el caso, con aviso previo mínimo de diez días hábiles. & Parcial & Cinco ventanas de congelamiento declaradas y ningún mantenimiento programado dentro de ellas; el aviso previo mínimo de diez días hábiles se declara en el Subdocumento 10 & Subdoc.3 §3.2.3 \textbar{} Subdoc.4 §4.2 (Ambientes) \\
\addlinespace[2pt]
RT-10.06 & La solución permitirá desplegar cambios sin interrupción del servicio. Las ventanas de indisponibilidad programada serán excepcionales y deberán justificarse caso a caso. & Sí & Despliegue azul-verde / actualización progresiva sin ventana de indisponibilidad & Subdoc.4 §4.2 \\
\addlinespace[2pt]
RT-10.07 & Se ejecutarán pruebas de resiliencia mediante inyección controlada de fallas —caída de instancia, de zona, de dependencia externa, latencia elevada, saturación de disco— antes de cada paso a producción y al menos una vez por semestre durante la Operación. & Parcial & Falla de zona simulada y ensayo de conmutación; el plan de pruebas de resiliencia (ingeniería del caos) se completa en el Subdocumento 9 & Subdoc.4 §4.2 (Tabla 4.11) \textbar{} §4.3.2 \\
\addlinespace[2pt]
RT-10.08 & El PROPONENTE documentará, por cada dependencia externa, el comportamiento de la solución cuando esa dependencia no responde, responde con error o responde con lentitud. & Sí & Comportamiento ante no respuesta documentado por cada integración externa & Subdoc.4 §4.1, §4.2 \\
\addlinespace[2pt]
RT-10.09 & Se declarará un presupuesto de error por servicio crítico y su vinculación con el ritmo de despliegue de cambios. & Sí & Presupuesto de error mensual de 0,1\% sobre el nivel de servicio de 99,9\%, que detiene el despliegue de cambios al agotarse & Subdoc.4 §4.2 \\
\addlinespace[2pt]
RT-11.01 & La arquitectura de seguridad se basará en el modelo Zero Trust conforme a NIST SP 800-207: verificación explícita de cada solicitud, privilegio mínimo y presunción de compromiso. & Sí & Arquitectura de seguridad Zero Trust (NIST SP 800-207), denegación por omisión en la frontera de datos (Servicio de control de cruces (X-01)) & Subdoc.4 §4.1 \\
\addlinespace[2pt]
RT-11.02 & El PROPONENTE entregará un modelado de amenazas documentado por cada componente y por cada integración externa, con metodología declarada (STRIDE u otra), y lo actualizará ante cada cambio arquitectónico relevante. & Sí & Modelado de amenazas con metodología STRIDE por componente e integración externa & Subdoc.4 §4.1 \\
\addlinespace[2pt]
RT-11.03 & La información del CLIENTE se clasificará por nivel de sensibilidad, con controles diferenciados por nivel documentados en una matriz. & Sí & Matriz documentada de clasificación de datos por nivel de sensibilidad & Subdoc.4 §4.1 \\
\addlinespace[2pt]
RT-11.04 & Existirá un programa de gestión de vulnerabilidades con escaneo continuo y plazos máximos de remediación de 7 días corridos para vulnerabilidades críticas, 15 días para altas y 30 días para medias, contados desde su publicación o detección. & Sí & Escaneo continuo de vulnerabilidades con plazos de remediación de 7/15/30 días & Subdoc.4 §4.1 \\
\addlinespace[2pt]
RT-11.05 & El PROPONENTE mantendrá una matriz de controles de seguridad trazable a ISO/IEC 27001 e ISO/IEC 27002, indicando el control, su implementación concreta en la solución y la evidencia que lo acredita. & Sí & Matriz de controles trazable a ISO/IEC 27001:2022 e ISO/IEC 27002:2022 & Subdoc.4 §4.1 \\
\addlinespace[2pt]
RT-11.06 & Se aplicarán los controles de ISO/IEC 27017 para los servicios en nube y de ISO/IEC 27018 para el tratamiento de datos personales en nube. & Sí & Controles ISO/IEC 27017 (nube) e ISO/IEC 27018:2025 (datos personales en nube pública) & Subdoc.4 §4.1 \\
\addlinespace[2pt]
RT-11.07 & La publicación de servicios se realizará exclusivamente a través de la capa de borde, con red de distribución de contenidos, cortafuegos de aplicaciones web con reglas gestionadas y personalizadas, y protección contra denegación de servicio distribuida en capas 3, 4 y 7. & Sí & Publicación exclusiva por la capa de borde (CDN + WAF + DDoS L3/L4/L7) & Subdoc.4 §4.1 \\
\addlinespace[2pt]
RT-11.08 & El cifrado en tránsito empleará TLS 1.3, con prohibición expresa de TLS 1.0 y 1.1, conjuntos de cifrado modernos, HSTS con precarga y gestión automatizada de certificados con rotación y alerta anticipada de vencimiento. & Sí & TLS 1.3 con prohibición expresa de TLS 1.0/1.1 & Subdoc.4 §4.1 \\
\addlinespace[2pt]
RT-11.09 & La totalidad de los datos en reposo estará cifrada, con claves gestionadas en un servicio de gestión de claves o en un módulo de seguridad de hardware, política de rotación declarada y separación de funciones en la custodia de claves. & Sí & Cifrado en reposo de toda la información persistida; Azure Key Vault (nube) y módulo de seguridad de hardware (HSM) para las claves del Emisor, producto y emplazamiento por definir, con rotación declarada & Subdoc.4 §4.1 \\
\addlinespace[2pt]
RT-11.10 & Los datos de categoría sensible que el caso identifique se cifrarán adicionalmente a nivel de campo, de modo que el acceso a la base de datos no revele su contenido. & Sí & Cifrado de campo y tokenización de identificación de cliente, cartera, comportamiento de pago y evaluación crediticia & Subdoc.4 §4.1 \\
\addlinespace[2pt]
RT-11.11 & La puerta de enlace de servicios aplicará autenticación, autorización, cuotas, límites de tasa, validación de esquema e inspección de carga útil. & Sí, condicionado & Azure API Management Premium (ADR-06, condicionada) como puerta de entrada de APIs, con autenticación, autorización, cuotas y límites de tasa & Subdoc.4 §4.1.1 (Tabla 4.6) \\
\addlinespace[2pt]
RT-11.12 & Los puntos de entrada públicos contarán con protección contra bots y abuso automatizado, con reto progresivo que no degrade la accesibilidad ni bloquee a personas usuarias legítimas. & Sí & Azure Front Door Premium con firewall de aplicaciones web y protección contra bots en el borde de cada ámbito & Subdoc.4 §4.2 (Redes y segmentación) \\
\addlinespace[2pt]
RT-11.13 & El PROPONENTE declarará la superficie de exposición completa de la solución: cada nombre de dominio, puerto y servicio alcanzable desde fuera de la red del CLIENTE. & Sí & Declaración de superficie de exposición (dominios, puertos, servicios) con línea base al cierre de Etapa 1 & Subdoc.4 §4.1, §4.2 \\
\addlinespace[2pt]
RT-11.14 & Los eventos de seguridad se registrarán de forma centralizada e inalterable, con retención mínima de doce meses en línea y veinticuatro meses adicionales en archivo recuperable. & Sí & Registro centralizado e inalterable de eventos de seguridad, 12 meses en línea + 24 en archivo & Subdoc.4 §4.1 \\
\addlinespace[2pt]
RT-11.15 & Los eventos se correlacionarán en una plataforma SIEM, con casos de uso de detección definidos específicamente para el proceso de negocio del caso, y no sólo genéricos de infraestructura. & Parcial & SIEM (producto por definir en el Subdocumento 4 y el Formulario T-11) con casos de uso definidos para el negocio & Subdoc.4 §4.1, §4.2 \\
\addlinespace[2pt]
RT-11.16 & Se implementará detección y respuesta en puntos finales y en cargas de trabajo, tanto en nube como on-premise. & Parcial & Detección en puntos finales (producto por definir); EDR en estaciones de trabajo & Subdoc.4 §4.2 \\
\addlinespace[2pt]
RT-11.17 & El PROPONENTE dispondrá de un centro de operaciones de seguridad con cobertura 24x7, propio o subcontratado, y declarará su ubicación, dotación y procedimientos. & Parcial & Vigilancia con Microsoft Sentinel y Defender for Cloud; cobertura 24x7, ubicación y dotación del SOC por declarar en el Subdocumento 10 & Subdoc.4 §4.2 (Tabla 4.9) \\
\addlinespace[2pt]
RT-11.18 & El plan de respuesta a incidentes de seguridad definirá clasificación, cadena de escalamiento, plazos, responsables y protocolo de comunicación al CLIENTE dentro de las dos horas de detectado un incidente de severidad crítica. & Sí & Plazo de 2 horas para comunicación de incidente de severidad crítica & Subdoc.4 §4.1 \\
\addlinespace[2pt]
RT-11.19 & Toda brecha de seguridad o de datos personales se notificará al CLIENTE dentro de las 24 horas de su detección, con informe preliminar, y el análisis de causa raíz se entregará dentro de los cinco días hábiles siguientes. & Sí & Notificación de brecha dentro de 24 horas & Subdoc.4 §4.1 \\
\addlinespace[2pt]
RT-11.20 & Se ejecutarán pruebas de intrusión por un tercero independiente del ADJUDICATARIO, anualmente y antes de cada paso a producción, con entrega íntegra del informe al CLIENTE y plan de remediación con plazos. & Sí & Pruebas de intrusión anuales y antes de cada paso a producción & Subdoc.4 §4.1 \\
\addlinespace[2pt]
RT-11.21 & Se realizarán ejercicios de simulación de incidente con participación del CLIENTE al menos una vez al año durante la Operación. & No & Pendiente: se declara en el diseño detallado de seguridad & Subdoc.4 §4.1 (Seguridad, continuidad y observabilidad) \\
\addlinespace[2pt]
RT-11.22 & El flujo de integración continua incorporará análisis estático de código, análisis de composición de software, análisis dinámico y escaneo de imágenes de contenedor, con criterios de bloqueo automático del despliegue ante hallazgos críticos. & Sí & Cadena DevSecOps con análisis estático, dinámico, de composición y de imágenes de contenedor, y bloqueo automático (herramientas por fijar en el Subdocumento 4 y el Formulario T-11) & Subdoc.4 §4.1, §4.2 \\
\addlinespace[2pt]
RT-11.23 & Cada versión liberada se acompañará de su inventario de componentes de software en formato CycloneDX o SPDX, entregado al CLIENTE. & Parcial & SBOM en CycloneDX o SPDX por versión liberada (herramienta por fijar) & Subdoc.4 §4.1, §4.2 \\
\addlinespace[2pt]
RT-11.24 & Los artefactos se firmarán y su procedencia se verificará conforme a SLSA nivel 3 o superior. & Parcial & Firma y procedencia SLSA nivel 3 (herramienta de firma por fijar) & Subdoc.4 §4.1, §4.2 \\
\addlinespace[2pt]
RT-11.25 & Queda prohibido el uso de datos productivos reales en ambientes no productivos sin anonimización o seudonimización verificable. & Sí & Prohibición de datos productivos reales en ambientes no productivos; datos sintéticos en Desarrollo y QA & Subdoc.4 §4.1, §4.2 \\
\addlinespace[2pt]
RT-11.26 & El PROPONENTE declarará el proceso de aprobación de nuevas dependencias de terceros, incluyendo criterios de licencia, mantención activa y ausencia de vulnerabilidades conocidas. & No & Pendiente: se declara en el diseño detallado de seguridad & Subdoc.4 §4.1 (Seguridad, continuidad y observabilidad) \\
\addlinespace[2pt]
RT-11.27 & Las personas desarrolladoras no tendrán acceso interactivo directo al ambiente de producción. Todo acceso excepcional será temporal, aprobado, registrado y con sesión grabada. & Sí & Sin acceso interactivo directo a producción para developers & Subdoc.4 §4.1 \\
\addlinespace[2pt]
RT-11.28 & Se aplicará el marco OWASP SAMM o equivalente para medir y mejorar la madurez del proceso de desarrollo seguro, con evaluación inicial y reevaluación anual. & No & Pendiente: se declara en el diseño detallado de seguridad & Subdoc.4 §4.1 (Seguridad, continuidad y observabilidad) \\
\addlinespace[2pt]
RT-12.01 & La gestión de identidad será centralizada, con federación mediante OpenID Connect y OAuth 2.1, o SAML 2.0 cuando la integración con el CLIENTE lo requiera, e integración con el directorio corporativo del CLIENTE por LDAP o su equivalente en la nube. & Sí & Identidad y gestión de accesos (base tecnológica) implementada con Microsoft Entra ID federado con el directorio del Cliente, con identidades separadas por ámbito (Retail y Emisor) sin confianza transitiva, con protocolo OIDC y OAuth 2.0 & Subdoc.4 §4.1 \\
\addlinespace[2pt]
RT-12.02 & Existirá inicio de sesión único para todos los módulos de la solución, con cierre de sesión propagado a todos ellos. & Parcial & Inicio de sesión con Microsoft Entra ID federado con el directorio del Cliente; propagación del cierre de sesión a todos los servicios por declarar & Subdoc.4 §4.1.1 (Tabla 4.6) \\
\addlinespace[2pt]
RT-12.03 & La autenticación multifactor será obligatoria para personas usuarias administradoras, para todo acceso privilegiado y para todo acceso originado fuera de la red corporativa. & Parcial & Microsoft Entra ID con acceso privado y Bastion para la administración; la autenticación multifactor obligatoria se declara en el diseño detallado & Subdoc.4 §4.2 (Redes y segmentación) \\
\addlinespace[2pt]
RT-12.04 & Se soportarán factores resistentes a la suplantación de identidad, tipo FIDO2 o claves de acceso, al menos para los perfiles administradores. & No & Pendiente: se declara en el diseño detallado de identidad & Subdoc.4 §4.1.1 (Tabla 4.6) \\
\addlinespace[2pt]
RT-12.05 & El control de acceso será basado en roles, complementado con control basado en atributos donde el proceso lo exija, con matriz de segregación de funciones documentada y verificable. & Sí & Matriz de roles y permisos con segregación de funciones entre originación, aprobación, modificación de condiciones y cobranza (dominio financiero) & Subdoc.3 §3.3.2 \\
\addlinespace[2pt]
RT-12.06 & Los accesos privilegiados se gestionarán con elevación temporal a demanda, aprobación previa y grabación de sesión para las operaciones de mayor riesgo. & Parcial & Administración mediante Bastion y acceso privado sin exposición directa; elevación temporal con aprobación y grabación de sesión por declarar & Subdoc.4 §4.2 (Redes y segmentación) \\
\addlinespace[2pt]
RT-12.07 & La política de sesión declarará duración máxima, caducidad por inactividad, renovación de la credencial de sesión tras la autenticación, revocación inmediata y control de sesiones concurrentes. & No & Pendiente: se declara en el diseño detallado de identidad & Subdoc.4 §4.1.1 (Tabla 4.6) \\
\addlinespace[2pt]
RT-12.08 & Las credenciales de sesión serán firmadas y de vida breve, con credencial de refresco rotatoria. Queda prohibido transportar identificadores de sesión en la ruta de la dirección web. & Sí & IAM Microsoft Entra ID con OIDC y OAuth 2.0 ``con credenciales de vida corta'' & Subdoc.4 §4.1 \\
\addlinespace[2pt]
RT-12.09 & Se registrará auditoría completa del ciclo de vida de la identidad: creación, modificación, elevación, bloqueo y baja de cuentas, con no repudio y retención declarada. & No & Pendiente: se declara en el diseño detallado de identidad & Subdoc.4 §4.1.1 (Tabla 4.6) \\
\addlinespace[2pt]
RT-12.10 & El aprovisionamiento y el desaprovisionamiento estarán automatizados y ligados al ciclo de vida laboral, con baja efectiva en un plazo no superior a 24 horas desde la desvinculación. & Sí & Conector bidireccional IAM-RRHH (base tecnológica, identidad y gestión de accesos) que revoca accesos automáticamente por evento de finiquito, con SLA interno de revocación total ≤ 4 horas & Subdoc.3 §3.3.2 (plataforma común) \textbar{} Subdoc.4 §4.1 \\
\addlinespace[2pt]
RT-12.11 & El mecanismo de autenticación se adecuará al perfil operacional real descrito en el caso: entornos de terreno, uso con guantes, baja alfabetización digital, dispositivos compartidos por turno y ausencia de correo electrónico personal. & Sí & Terminal móvil compartido con traspaso rápido de sesión (hot-swapping) sin autenticación biométrica que añada tiempo al vendedor; equipamiento con grado IP54 y uso confirmado con guantes en centro de distribución & Subdoc.4 §4.2 \\
\addlinespace[2pt]
RT-12.12 & Las personas usuarias externas del CLIENTE —clientes, proveedores, pacientes, productores, según el caso— dispondrán de un mecanismo de registro, verificación de identidad y recuperación de acceso autoservido y seguro. & No & Pendiente: se declara en el diseño detallado de identidad & Subdoc.4 §4.1.1 (Tabla 4.6) \\
\addlinespace[2pt]
RT-12.13 & El PROPONENTE declarará el procedimiento de acceso de emergencia (cuenta de último recurso), su custodia, su control y su auditoría. & No & Pendiente: se declara en el diseño detallado de identidad & Subdoc.4 §4.1.1 (Tabla 4.6) \\
\addlinespace[2pt]
RT-13.01 & Todas las interfaces destinadas a personas usuarias cumplirán WCAG 2.2 nivel AA, verificado con herramientas automatizadas y con pruebas manuales, con informe de conformidad entregable. & No & Pendiente: se declara en el plan de calidad y el diseño de interfaz de la Etapa 1 & Subdoc.9 \\
\addlinespace[2pt]
RT-13.02 & El diseño será responsivo y funcionará correctamente en escritorio, tableta y teléfono, con puntos de quiebre coherentes y adaptación inteligente del contenido, no simple reducción. & No & Pendiente: se declara en el plan de calidad y el diseño de interfaz de la Etapa 1 & Subdoc.9 \\
\addlinespace[2pt]
RT-13.03 & El PROPONENTE ejecutará investigación con personas usuarias reales del CLIENTE, prototipado y pruebas de usabilidad antes de la construcción definitiva, y documentará los hallazgos y los cambios de diseño que produjeron. & No & Pendiente: se declara en el plan de calidad y el diseño de interfaz de la Etapa 1 & Subdoc.9 \\
\addlinespace[2pt]
RT-13.04 & Se comprometerán indicadores de usabilidad medibles: tiempo máximo de la transacción operacional crítica, número máximo de pasos, tasa de error tolerada y tiempo de aprendizaje esperado por perfil. & Sí & Requerimiento no funcional que fija el tiempo de capacitación operativa básica en interfaces de sala, caja y mesón financiero & Subdoc.3 §3.3.2 \\
\addlinespace[2pt]
RT-13.05 & Ninguna funcionalidad principal requerirá más de tres interacciones desde la pantalla de inicio del perfil correspondiente. & No & Pendiente: se declara en el plan de calidad y el diseño de interfaz de la Etapa 1 & Subdoc.9 \\
\addlinespace[2pt]
RT-13.06 & La solución entregará retroalimentación visual clara ante cada acción, manejará los errores de forma comprensible —indicando qué ocurrió y qué hacer— y evitará mensajes técnicos dirigidos a la persona usuaria final. & No & Pendiente: se declara en el plan de calidad y el diseño de interfaz de la Etapa 1 & Subdoc.9 \\
\addlinespace[2pt]
RT-13.07 & La solución soportará a personas usuarias con baja alfabetización digital: alto contraste, objetivos táctiles de al menos 44x44 píxeles, iconografía acompañada de texto, y flujos guiados paso a paso. & No & Pendiente: se declara en el plan de calidad y el diseño de interfaz de la Etapa 1 & Subdoc.9 \\
\addlinespace[2pt]
RT-13.08 & Cuando el caso lo requiera, las interfaces de terreno operarán con guantes, a la intemperie, con luminosidad variable y sin conexión. & Sí & Terminal móvil de consulta de existencia y lector de radiofrecuencia de centro de distribución, especificados con grado de protección IP54, pantalla legible con la luminosidad de sala y uso confirmado con guantes & Subdoc.4 §4.2 \\
\addlinespace[2pt]
RT-13.09 & Existirá un sistema de diseño documentado con paleta de colores acotada, jerarquía tipográfica de a lo más dos familias, iconografía coherente, retícula y componentes reutilizables. & No & Pendiente: se declara en el plan de calidad y el diseño de interfaz de la Etapa 1 & Subdoc.9 \\
\addlinespace[2pt]
RT-13.10 & El PROPONENTE declarará la matriz de navegadores y de versiones soportadas y su política de actualización. & No & Pendiente: se declara en el plan de calidad y el diseño de interfaz de la Etapa 1 & Subdoc.9 \\
\addlinespace[2pt]
RT-13.11 & La solución será navegable íntegramente por teclado, con orden de foco lógico y atajos para las operaciones frecuentes. & No & Pendiente: se declara en el plan de calidad y el diseño de interfaz de la Etapa 1 & Subdoc.9 \\
\addlinespace[2pt]
RT-13.12 & Se valorará el soporte de modo oscuro, de personalización de la interfaz por persona usuaria y de múltiples idiomas cuando el caso lo justifique. & No & Pendiente: se declara en el plan de calidad y el diseño de interfaz de la Etapa 1 & Subdoc.9 \\
\addlinespace[2pt]
RT-14.01 & La observabilidad será unificada para nube y on-premise, con métricas, registros y trazas distribuidas correlacionadas por un identificador único de transacción, instrumentadas conforme a OpenTelemetry. & Sí, condicionado & Observabilidad y degradación controlada (base tecnológica) instrumentada con OpenTelemetry y Azure Monitor & Subdoc.4 §4.1.1 (Tabla 4.6) \\
\addlinespace[2pt]
RT-14.02 & El CLIENTE dispondrá de acceso propio y permanente a los tableros operacionales y de negocio, con datos en tiempo real y capacidad de exportación. & No & Pendiente: se declara en los planes en operación (Subdocumento 11) y el nivel de servicio (Subdocumento 10) & Subdoc.10 y Subdoc.11 (por redactar) \\
\addlinespace[2pt]
RT-14.03 & Los indicadores de nivel de servicio se medirán sobre la experiencia real de la persona usuaria y no sobre pruebas sintéticas, sin perjuicio de que estas se empleen como complemento. & No & Pendiente: se declara en los planes en operación (Subdocumento 11) y el nivel de servicio (Subdocumento 10) & Subdoc.10 y Subdoc.11 (por redactar) \\
\addlinespace[2pt]
RT-14.04 & El alertamiento se basará en síntomas de negocio y no sólo en umbrales de infraestructura, con supresión de ruido, agrupación, escalamiento automático y turnos de disponibilidad declarados. & No & Pendiente: se declara en los planes en operación (Subdocumento 11) y el nivel de servicio (Subdocumento 10) & Subdoc.10 y Subdoc.11 (por redactar) \\
\addlinespace[2pt]
RT-14.05 & Existirá un libro de operación y una guía de resolución documentados para cada escenario de falla previsible, con automatización progresiva de las tareas repetitivas. & No & Pendiente: se declara en los planes en operación (Subdocumento 11) y el nivel de servicio (Subdocumento 10) & Subdoc.10 y Subdoc.11 (por redactar) \\
\addlinespace[2pt]
RT-14.06 & Todo incidente crítico dará lugar a un análisis de causa raíz obligatorio, con informe entregable dentro de cinco días hábiles y seguimiento de las acciones correctivas hasta su cierre. & No & Pendiente: se declara en los planes en operación (Subdocumento 11) y el nivel de servicio (Subdocumento 10) & Subdoc.10 y Subdoc.11 (por redactar) \\
\addlinespace[2pt]
RT-14.07 & Los registros de la solución no contendrán datos personales sensibles ni credenciales, y su acceso estará controlado y auditado. & No & Pendiente: se declara en los planes en operación (Subdocumento 11) y el nivel de servicio (Subdocumento 10) & Subdoc.10 y Subdoc.11 (por redactar) \\
\addlinespace[2pt]
RT-14.08 & El PROPONENTE declarará la retención de métricas, registros y trazas, y su costo asociado, distinguiendo el almacenamiento en línea del archivado. & No & Pendiente: se declara en los planes en operación (Subdocumento 11) y el nivel de servicio (Subdocumento 10) & Subdoc.10 y Subdoc.11 (por redactar) \\
\addlinespace[2pt]
RT-14.09 & Se valorará la detección proactiva de anomalías mediante análisis del comportamiento histórico, con alerta antes de que el incidente afecte a la operación. & No & Pendiente: se declara en los planes en operación (Subdocumento 11) y el nivel de servicio (Subdocumento 10) & Subdoc.10 y Subdoc.11 (por redactar) \\
\addlinespace[2pt]
RT-15.01 & El PROPONENTE dimensionará la infraestructura ajustada a la demanda real, evitando capacidad ociosa permanente, y declarará el factor de utilización proyectado. & No & Pendiente: se declara en la acreditación en el Subdocumento 12 y el Formulario T-8 & Subdoc.12 y Formulario T-8 (por redactar) \\
\addlinespace[2pt]
RT-15.02 & Los ambientes no productivos se apagarán o reducirán fuera del horario de uso. & Sí & Política de reducción de ambientes de Desarrollo y QA fuera de la ventana 08:00-20:00 de días hábiles, con Preproducción activa solo durante ventanas de prueba programadas & Subdoc.4 §4.2 \\
\addlinespace[2pt]
RT-15.03 & El PROPONENTE estimará la huella de carbono anual de la operación de la solución y declarará la metodología empleada. & No & Pendiente: se declara en la acreditación en el Subdocumento 12 y el Formulario T-8 & Subdoc.12 y Formulario T-8 (por redactar) \\
\addlinespace[2pt]
RT-15.04 & Se declarará la eficiencia en el uso de la energía (PUE) del recinto on-premise y la intensidad de carbono de la región de nube escogida. & Sí & Cálculo de PUE objetivo de 1,6 para el recinto on-premise (DC-01) & Subdoc.4 §4.2 \\
\addlinespace[2pt]
RT-15.05 & Se valorará la elección de regiones de nube con menor intensidad de carbono cuando la latencia y la regulación lo permitan, con el análisis comparativo correspondiente. & No & Pendiente: se declara en la acreditación en el Subdocumento 12 y el Formulario T-8 & Subdoc.12 y Formulario T-8 (por redactar) \\
\addlinespace[2pt]
RT-15.06 & Se valorará la definición de metas de reducción del consumo durante la Operación, con medición y reporte anual. & No & Pendiente: se declara en la acreditación en el Subdocumento 12 y el Formulario T-8 & Subdoc.12 y Formulario T-8 (por redactar) \\
\addlinespace[2pt]
RT-15.07 & Las certificaciones se acreditarán con copia del certificado vigente y con el código de verificación del organismo emisor cuando exista. & No & Pendiente: se declara en la acreditación en el Subdocumento 12 y el Formulario T-8 & Subdoc.12 y Formulario T-8 (por redactar) \\
\addlinespace[2pt]
RT-15.08 & Las personas certificadas formarán parte del equipo efectivamente asignado al PROYECTO, con dedicación declarada. No se aceptará acreditar personal que no participe. & Parcial & Equipo de jefaturas nominado en el Subdocumento 12 y el Formulario T-8; las certificaciones individuales se acreditan allí con sus copias & Subdoc.1 §1.5 \textbar{} Subdoc.12 y Formulario T-8 (por redactar) \\
\addlinespace[2pt]
RT-15.09 & El ADJUDICATARIO mantendrá vigentes estas certificaciones durante todo el período contractual y las repondrá ante la salida de una persona certificada, conforme al Artículo 76° de las Bases Administrativas. & No & Pendiente: se declara en la acreditación en el Subdocumento 12 y el Formulario T-8 & Subdoc.12 y Formulario T-8 (por redactar) \\
\addlinespace[2pt]
RT-16.01 & La solución dispondrá de un módulo de administración que permita al CLIENTE, sin intervención del ADJUDICATARIO, gestionar personas usuarias, roles, permisos, unidades organizacionales y sus jerarquías. & No & Pendiente: se declara en el diseño detallado de la Etapa 1 (administración, flujos, documentos y notificaciones) & Subdoc.3 §3.2.2 \\
\addlinespace[2pt]
RT-16.02 & Las reglas de negocio parametrizables —umbrales, plazos, montos, tolerancias, catálogos, listas de valores, textos de notificación— serán configurables desde la interfaz de administración, con control de versiones y registro de quién cambió qué y cuándo. & Sí & Múltiples requerimientos funcionales que exigen parametrización de umbrales, ventanas y catálogos de negocio (colchón de reserva, ventanas de conteo cíclico, ventanas de congelamiento, plazos de garantía), junto con registro de auditoría de cambios & Subdocumento 2, Anexos A.1 y A.2 (p. ej. RF-131 a RF-134, RF-184, RF-195 y Tabla A.2.7, RNF-15) \\
\addlinespace[2pt]
RT-16.03 & Todo cambio de parámetro con impacto operacional requerirá aprobación de un segundo perfil y quedará registrado con su justificación. & No & Pendiente: se declara en el diseño detallado de la Etapa 1 (administración, flujos, documentos y notificaciones) & Subdoc.3 §3.2.2 \\
\addlinespace[2pt]
RT-16.04 & El PROPONENTE declarará expresamente qué elementos son parametrizables y cuáles requieren desarrollo. Presentar como parametrizable lo que exige desarrollo se evaluará como observación grave. & Sí & Catálogo atómico de requerimientos (versión vigente 3.1: 227 funcionales y 72 no funcionales) que identifica explícitamente, requerimiento por requerimiento, cuáles son parametrizables & Subdocumento 2, Anexos A.1 y A.2 (listados de requerimientos funcionales y no funcionales) \\
\addlinespace[2pt]
RT-16.05 & Existirá un ambiente de simulación que permita probar el efecto de un cambio de parámetro antes de aplicarlo a producción. & No & Pendiente: se declara en el diseño detallado de la Etapa 1 (administración, flujos, documentos y notificaciones) & Subdoc.3 §3.2.2 \\
\addlinespace[2pt]
RT-16.06 & Toda operación que cree, modifique o elimine información quedará registrada con identificación de la persona o del sistema que la ejecutó, fecha y hora con zona horaria, origen, valores anteriores y valores posteriores. & Sí & Requerimiento no funcional de auditabilidad general: toda acción relevante ejecutada en el sistema queda registrada con usuario individual, acción e instante & Subdoc.3 §3.3.2 \\
\addlinespace[2pt]
RT-16.07 & El registro de auditoría será inalterable y no podrá ser modificado ni eliminado por ningún perfil, incluido el administrador de la plataforma. & Sí & Bitácora de gobernanza y cruces (Servicio de control de cruces (X-01)), almacén de solo anexión (append-only) con valor forense & Subdocumento 5, sección 5.1.3 ``Entidades transversales'' \\
\addlinespace[2pt]
RT-16.08 & El CLIENTE podrá consultar y exportar la auditoría desde la interfaz, con filtros por persona, período, entidad y tipo de operación, sin requerir acceso a la base de datos. & No & Pendiente: se declara en el diseño detallado de la Etapa 1 (administración, flujos, documentos y notificaciones) & Subdoc.3 §3.2.2 \\
\addlinespace[2pt]
RT-16.09 & Las consultas a información sensible quedarán registradas, no sólo las modificaciones. & Sí & Servicio de control de cruces (X-01): todo acceso a la cartera de crédito, comportamiento de pago y antecedentes de evaluación crediticia queda registrado (usuario, dato, finalidad, instante) & Subdoc.3 §3.3.2 (Frente Filial emisora, Frente Frontera) \textbar{} Subdoc.4 §4.1 (Tabla 4.2) \\
\addlinespace[2pt]
RT-16.10 & El período de retención de la auditoría será el que fije el caso y, en su defecto, no inferior a cinco años. & Sí & Política de retención de datos gobernada centralizada por dominio (documentos tributarios y de venta: 6 años; antecedentes de crédito: plazo del crédito + 6 años; precio: 3 años), fijada explícitamente por el caso & Subdocumento 5, sección 5.8.2 ``Retención'' (Tabla 5.X) \\
\addlinespace[2pt]
RT-16.11 & La solución soportará flujos de trabajo con estados, transiciones, responsables, plazos, escalamiento automático por vencimiento y delegación por ausencia. & No & Pendiente: se declara en el diseño detallado de la Etapa 1 (administración, flujos, documentos y notificaciones) & Subdoc.3 §3.2.2 \\
\addlinespace[2pt]
RT-16.12 & Los flujos serán configurables por el CLIENTE sin desarrollo, al menos en lo relativo a responsables, plazos y niveles de aprobación. & No & Pendiente: se declara en el diseño detallado de la Etapa 1 (administración, flujos, documentos y notificaciones) & Subdoc.3 §3.2.2 \\
\addlinespace[2pt]
RT-16.13 & Toda solicitud pendiente será visible para su responsable en una bandeja de tareas unificada, con priorización y alerta de vencimiento. & No & Pendiente: se declara en el diseño detallado de la Etapa 1 (administración, flujos, documentos y notificaciones) & Subdoc.3 §3.2.2 \\
\addlinespace[2pt]
RT-16.14 & El motor de reglas permitirá definir condiciones de negocio evaluables sin recompilación, con trazabilidad de qué regla se aplicó a cada transacción. & No & Pendiente: se declara en el diseño detallado de la Etapa 1 (administración, flujos, documentos y notificaciones) & Subdoc.3 §3.2.2 \\
\addlinespace[2pt]
RT-16.15 & La solución gestionará documentos con versionado, metadatos, control de acceso, búsqueda por contenido y por metadato, y previsualización sin descarga. & No & Pendiente: se declara en el diseño detallado de la Etapa 1 (administración, flujos, documentos y notificaciones) & Subdoc.3 §3.2.2 \\
\addlinespace[2pt]
RT-16.16 & Los documentos se almacenarán cifrados, con verificación de integridad y retención conforme a la política del caso. & Sí & Expediente de evidencia (Servicio de evidencia financiera), custodiado bajo política WORM, sellado criptográfico SHA-256 encadenado y almacenamiento cifrado & Subdocumento 5, sección 5.6 ``Modelo de evidencia de consentimiento'' \\
\addlinespace[2pt]
RT-16.17 & La solución soportará firma electrónica conforme a la Ley N° 19.799, con firma avanzada para los actos que el caso lo requiera, y verificación de validez del certificado al momento de la firma. & Sí & Firma electrónica avanzada (Servicio de evidencia financiera) integrada con prestador acreditado por la autoridad chilena (Acepta, E-Sign o TOC Biometrics), conforme a Ley N° 19.799, con verificación de validez del certificado al firmar & Subdoc.4 §4.1 \\
\addlinespace[2pt]
RT-16.18 & Se generará el sello de tiempo y se conservará la evidencia de firma que permita verificar el documento con posterioridad al vencimiento del certificado. & Sí & Expediente de consentimiento sellado criptográficamente (SHA-256 encadenado) y custodiado bajo WORM, con retención del plazo del crédito más 6 años & Subdoc.4 §4.1 \textbar{} Subdocumento 5, sección 5.6 ``Modelo de evidencia de consentimiento'' \\
\addlinespace[2pt]
RT-16.19 & La solución generará documentos a partir de plantillas administrables por el CLIENTE, con datos de la transacción y salida en formato abierto. & No & Pendiente: se declara en el diseño detallado de la Etapa 1 (administración, flujos, documentos y notificaciones) & Subdoc.3 §3.2.2 \\
\addlinespace[2pt]
RT-16.20 & La solución enviará notificaciones por al menos tres canales: correo electrónico, notificación en la aplicación y mensajería instantánea o SMS, según lo que el caso requiera. & No & Pendiente: se declara en el diseño detallado de la Etapa 1 (administración, flujos, documentos y notificaciones) & Subdoc.3 §3.2.2 \\
\addlinespace[2pt]
RT-16.21 & Las plantillas de notificación serán administrables por el CLIENTE, con variables de la transacción, y versionadas. & No & Pendiente: se declara en el diseño detallado de la Etapa 1 (administración, flujos, documentos y notificaciones) & Subdoc.3 §3.2.2 \\
\addlinespace[2pt]
RT-16.22 & Cada persona usuaria podrá configurar sus preferencias de canal y de frecuencia, respetando las notificaciones que el CLIENTE defina como obligatorias. & No & Pendiente: se declara en el diseño detallado de la Etapa 1 (administración, flujos, documentos y notificaciones) & Subdoc.3 §3.2.2 \\
\addlinespace[2pt]
RT-16.23 & El envío será asíncrono, con reintento ante falla, control de duplicados y registro de entrega, apertura y error por cada mensaje. & No & Pendiente: se declara en el diseño detallado de la Etapa 1 (administración, flujos, documentos y notificaciones) & Subdoc.3 §3.2.2 \\
\addlinespace[2pt]
RT-16.24 & El PROPONENTE declarará el proveedor de cada canal, su costo unitario, su volumen proyectado y el tratamiento del costo variable en la Oferta Económica. & No & Pendiente: se declara en el diseño detallado de la Etapa 1 (administración, flujos, documentos y notificaciones) & Subdoc.3 §3.2.2 \\
\addlinespace[2pt]
RT-16.25 & Las notificaciones respetarán la normativa de comunicaciones comerciales y permitirán la baja cuando corresponda. & No & Pendiente: se declara en el diseño detallado de la Etapa 1 (administración, flujos, documentos y notificaciones) & Subdoc.3 §3.2.2 \\
\addlinespace[2pt]
RT-16.26 & Se valorará la integración con canales conversacionales que permitan a la persona usuaria responder y ejecutar acciones desde el propio canal. & No & Pendiente: se declara en el diseño detallado de la Etapa 1 (administración, flujos, documentos y notificaciones) & Subdoc.3 §3.2.2 \\
\addlinespace[2pt]
RT-16.27 & La solución dispondrá de búsqueda global con indexación de texto completo, tolerancia a errores de escritura, filtros facetados y respeto del control de acceso de la persona que busca. & Parcial & Motor de búsqueda y catálogo (producto por definir), búsqueda facetada y reindexación incremental & Subdoc.4 §4.1.1 (Tabla 4.6) \\
\addlinespace[2pt]
RT-16.28 & Los listados serán ordenables, filtrables, paginados y exportables en formatos abiertos, con el filtro aplicado reflejado en la exportación. & No & Pendiente: se declara en el diseño detallado de la Etapa 1 (administración, flujos, documentos y notificaciones) & Subdoc.3 §3.2.2 \\
\addlinespace[2pt]
RT-16.29 & Las exportaciones de gran volumen se procesarán de forma asíncrona, con notificación al completarse y sin bloquear la sesión. & No & Pendiente: se declara en el diseño detallado de la Etapa 1 (administración, flujos, documentos y notificaciones) & Subdoc.3 §3.2.2 \\
\addlinespace[2pt]
RT-16.30 & Toda exportación de información sensible quedará registrada en la auditoría, con identificación de quién exportó qué y cuándo. & No & Pendiente: se declara en el diseño detallado de la Etapa 1 (administración, flujos, documentos y notificaciones) & Subdoc.3 §3.2.2 \\
\addlinespace[2pt]
RT-16.31 & La solución dispondrá de un portal público con la información que el caso determine, accesible sin autenticación, con los mismos estándares de accesibilidad y desempeño que el resto de la plataforma. & Sí & Canal digital (comercio electrónico vigente): expone catálogo público con precio y disponibilidad y consulta de estado de pedido sin autenticación, servido tras CDN & Subdoc.3 §3.3.2 (Frente Retail) \textbar{} Subdoc.4 §4.1 (Tabla 4.2) \\
\addlinespace[2pt]
RT-16.32 & Las personas usuarias externas dispondrán de autoatención para las consultas de mayor frecuencia, evitando el contacto telefónico para operaciones simples. & Sí & Canal digital (comercio electrónico vigente): consulta de estado de pedido y notificaciones transaccionales sin necesidad de contacto reiterado a la compañía & Subdoc.3 §3.3.2 (Frente Retail) \\
\addlinespace[2pt]
RT-16.33 & El PROPONENTE estimará la reducción esperada del volumen de atención asistida por efecto de la autoatención y comprometerá el indicador. & No & Pendiente: se declara en el diseño detallado de la Etapa 1 (administración, flujos, documentos y notificaciones) & Subdoc.3 §3.2.2 \\
\addlinespace[2pt]
RT-16.34 & El portal público resistirá picos de tráfico sin degradar los servicios transaccionales internos, mediante aislamiento de recursos y caché. & Sí & Portal del cliente servido tras CDN, absorbiendo el pico del evento anual mayoritariamente en la capa de borde & Subdoc.4 §4.1 \\
\addlinespace[2pt]
RT-17.01 & La solución proveerá una aplicación móvil para los perfiles operacionales que el caso identifique, con funcionamiento sin conexión y sincronización diferida cuando la operación en terreno lo requiera. & Sí & Aplicación web progresiva (PWA) de Punto de Venta con service workers e IndexedDB, replicada en el nodo de tienda para operar 8 horas sin conexión & Subdoc.4 §4.1 \\
\addlinespace[2pt]
RT-17.02 & El PROPONENTE declarará si la aplicación es nativa, híbrida o web progresiva, y justificará la decisión frente al requisito de operación desconectada, al acceso a periféricos y al costo de mantención. & Sí & Decisión explícita por PWA, justificada por la necesidad de ejecutar indistintamente en las 640 terminales compartidas sin instalación por dispositivo & Subdoc.4 §4.2 \\
\addlinespace[2pt]
RT-17.03 & La aplicación soportará las versiones de sistema operativo móvil vigentes y las dos anteriores, con política de actualización declarada. & No & Pendiente: se declara en el diseño de la aplicación de campo de la Etapa 1 & Subdoc.4 §4.1 \\
\addlinespace[2pt]
RT-17.04 & La aplicación se distribuirá por las tiendas oficiales o mediante gestión de flota corporativa, con firma de la aplicación y verificación de integridad. & No & Pendiente: se declara en el diseño de la aplicación de campo de la Etapa 1 & Subdoc.4 §4.1 \\
\addlinespace[2pt]
RT-17.05 & La información almacenada en el dispositivo estará cifrada, con borrado remoto y bloqueo ante pérdida o desvinculación de la persona usuaria. & No & Pendiente: se declara en el diseño de la aplicación de campo de la Etapa 1 & Subdoc.4 §4.1 \\
\addlinespace[2pt]
RT-17.06 & La aplicación integrará los periféricos que el caso requiera: cámara, lector de códigos, NFC, GPS, impresora de etiquetas, balanza o báscula. & Sí & Terminal móvil compartido con lector de código integrado; integración de la red de videovigilancia (1.840 cámaras) como periférico de red & Subdoc.4 §4.2 \\
\addlinespace[2pt]
RT-17.07 & El consumo de datos móviles y de batería se optimizará y se declarará el consumo estimado por turno de trabajo. & No & Pendiente: se declara en el diseño de la aplicación de campo de la Etapa 1 & Subdoc.4 §4.1 \\
\addlinespace[2pt]
RT-17.08 & Se valorará la disponibilidad de una versión de la interfaz para dispositivos de bajo costo o de generaciones anteriores, ampliando la cobertura de personas usuarias. & No & Pendiente: se declara en el diseño de la aplicación de campo de la Etapa 1 & Subdoc.4 §4.1 \\
\addlinespace[2pt]
RT-18.01 & El PROPONENTE declarará cada componente de inteligencia artificial: propósito, modelo o servicio empleado, proveedor, versión y ubicación de procesamiento de los datos. & Parcial & Componentes de inteligencia artificial declarados en las fichas de innovación (propósito, modelo o servicio y ubicación) y en los modelos predictivos del Subdocumento 4; proveedor y versión por completar & Subdoc.13 \textbar{} Formulario T-19 (campos 4 y 7) \\
\addlinespace[2pt]
RT-18.02 & Se garantizará contractualmente que los datos del CLIENTE no serán utilizados para entrenar modelos de terceros, salvo autorización expresa y escrita. & Parcial & El proveedor no entrena con datos de Ancoa en la búsqueda multimodal; la cláusula contractual general se declara en el Subdocumento 12 & Formulario T-19 (ficha 3) \textbar{} Subdoc.12 (por redactar) \\
\addlinespace[2pt]
RT-18.03 & Se documentarán los límites de uso, los casos en que el resultado requiere validación humana previa y el procedimiento de supervisión. & No & Pendiente: límites de uso, validación humana previa y supervisión de los componentes de inteligencia artificial, por declarar en el Subdocumento 13 y en el diseño detallado & Subdoc.13 \textbar{} Formulario T-19 \\
\addlinespace[2pt]
RT-18.04 & Los riesgos de sesgo, alucinación, fuga de información y uso indebido se evaluarán y mitigarán conforme al NIST AI Risk Management Framework 1.0 y a la norma ISO/IEC 42001. & Parcial & Evaluación de riesgos de sesgo, fuga y uso indebido con el NIST AI Risk Management Framework en las innovaciones con IA & Subdoc.13 \\
\addlinespace[2pt]
RT-18.05 & Toda interacción relevante con un componente de inteligencia artificial quedará registrada para efectos de auditoría y trazabilidad, incluyendo la entrada, la salida y la decisión humana posterior. & No & Pendiente: registro de las interacciones relevantes con componentes de inteligencia artificial, por declarar en el Subdocumento 13 y en el diseño detallado & Subdoc.13 \textbar{} Formulario T-19 \\
\addlinespace[2pt]
RT-18.06 & El componente podrá desactivarse sin comprometer la operación del resto de la solución, y existirá un procedimiento manual de respaldo para la función que automatiza. & No & Pendiente: desactivación del componente sin afectar al resto de la solución y procedimiento manual de respaldo, por declarar en el Subdocumento 13 y en el diseño detallado & Subdoc.13 \textbar{} Formulario T-19 \\
\addlinespace[2pt]
RT-18.07 & Cuando el resultado se presente a una persona usuaria, se indicará expresamente que fue generado o sugerido de forma automática, y su nivel de confianza cuando el modelo lo provea. & No & Pendiente: indicación al usuario de que el resultado es automático y de su nivel de confianza, por declarar en el Subdocumento 13 y en el diseño detallado & Subdoc.13 \textbar{} Formulario T-19 \\
\addlinespace[2pt]
RT-18.08 & Los modelos predictivos declararán sus variables de entrada, su métrica de desempeño, su línea base y su plan de reentrenamiento y de detección de deriva. & No & Pendiente: variables, métrica, línea base y plan de reentrenamiento de los modelos predictivos, por declarar en el Subdocumento 13 y en el diseño detallado & Subdoc.13 \textbar{} Formulario T-19 \\
\addlinespace[2pt]
RT-18.09 & La responsabilidad por los resultados de los componentes de inteligencia artificial recae íntegramente en el ADJUDICATARIO, conforme al Artículo 86° de las Bases Administrativas. & No & Pendiente: asunción de responsabilidad por los resultados de la inteligencia artificial, por declarar en el Subdocumento 13 y en el diseño detallado & Subdoc.13 \textbar{} Formulario T-19 \\
\addlinespace[2pt]
RT-18.10 & Se valorará el uso de automatización robótica de procesos o de agentes para tareas repetitivas de back office, con el ahorro de horas cuantificado y reflejado en la Oferta Económica. & No & Pendiente: automatización robótica de procesos o agentes para tareas de back office, por declarar en el Subdocumento 13 y en el diseño detallado & Subdoc.13 \textbar{} Formulario T-19 \\
\addlinespace[2pt]
RT-19.01 & El ADJUDICATARIO constituirá una oficina de gestión del proyecto con metodología declarada, basada en el PMBOK del Project Management Institute e integrando prácticas ágiles donde el PROPONENTE lo justifique. & Sí & Oficina de gestión del proyecto con metodología PMBOK adaptada con enfoque ágil & Formulario T-9 §T-9.1 \textbar{} Subdoc.6 §6.1 \\
\addlinespace[2pt]
RT-19.02 & El Jefe de Proyecto tendrá dedicación exclusiva durante toda la fase de implementación y facultades para comprometer al ADJUDICATARIO en materias de ejecución. & Sí & Jefe de Proyecto con dedicación exclusiva al frente de un ``equipo experto con dedicación exclusiva al proyecto'' & Subdocumento 1, sección 1.5 ``Estructura para Proyecto'' (rol); Subdocumento 12 y Formulario T-8 (nombre, certificación y facultades) \\
\addlinespace[2pt]
RT-19.03 & Existirá un procedimiento formal de control de cambios conforme al Artículo 72° de las Bases Administrativas, con registro de cambios, análisis de impacto y aprobación previa a la ejecución. & Sí & Procedimiento de control integrado de cambios con registro, análisis de impacto y decisión & Formulario T-9 §T-9.3 (Procedimiento de control integrado de cambios) \\
\addlinespace[2pt]
RT-19.04 & La gestión del riesgo seguirá la norma ISO 31000, con registro de riesgos vivo, revisión en cada Comité de Proyecto y planes de mitigación con responsable, plazo y disparador. & Sí & Gestión del riesgo conforme a ISO 31000, con registro vivo y revisión en el Comité de Proyecto & Subdoc.8 \textbar{} Formulario T-16 \textbar{} Formulario T-9 §T-9.1 \\
\addlinespace[2pt]
RT-19.05 & El ADJUDICATARIO habilitará un espacio colaborativo accesible al CLIENTE, con la documentación del proyecto, los entregables, las actas, el registro de riesgos y el registro de cambios siempre actualizados. & Sí & Sistema de información del proyecto (SharePoint, Jira y Slack) con documentación, actas y registros accesibles al Cliente & Formulario T-9 §T-9.5 (Sistema de información para la dirección del proyecto) \\
\addlinespace[2pt]
RT-19.06 & EL ADJUDICATARIO entregará un informe mensual de avance con estado del cronograma, avance físico y financiero, entregables del período, desviaciones, riesgos, incidencias y compromisos del período siguiente. & Sí & Informe mensual de avance con cronograma, avance, entregables, desviaciones y riesgos & Formulario T-9 §T-9.5 (Contenidos mínimos) \\
\addlinespace[2pt]
RT-19.07 & El avance se medirá con valor ganado, reportando el índice de desempeño del cronograma y del costo, y no mediante declaración cualitativa de porcentaje de avance. & Sí & Medición del desempeño con valor ganado: índices de desempeño del cronograma y del costo & Formulario T-9 §T-9.3 (Medición del desempeño con valor ganado) \\
\addlinespace[2pt]
RT-19.08 & El CLIENTE dispondrá de un tablero de estado del proyecto actualizado, accesible en cualquier momento. & Sí & Tablero de estado del proyecto en el sistema de información, accesible al Cliente & Formulario T-9 §T-9.5 \\
\addlinespace[2pt]
RT-19.09 & Las actas de todos los comités del Artículo 71 de las Bases Administrativas se levantarán dentro de los dos días hábiles siguientes y registrarán acuerdos, responsables y plazos. & Sí & Actas de los comités emitidas dentro de los dos días hábiles siguientes, con acuerdos y responsables & Formulario T-9 §T-9.2 (Instancias de gobierno) \textbar{} §T-9.5 (Matriz de comunicaciones) \\
\addlinespace[2pt]
RT-19.10 & Toda desviación superior al 10\% en un hito se comunicará dentro de los cinco días hábiles de detectada, con plan de recuperación. & Sí & Aviso de toda desviación mayor al 10 \% en un hito dentro de los cinco días hábiles, con plan de recuperación & Formulario T-9 §T-9.5 (Matriz de comunicaciones) \\
\addlinespace[2pt]
RT-20.01 & El PROPONENTE presentará un plan de implantación por sitio y por perfil, coherente con el cronograma obligatorio del Artículo 17° de las Bases Administrativas. & Parcial & Plan por sitio y por perfil (14 tiendas en centros comerciales, 8 a la calle, tienda insignia y dos centros de distribución); la asignación a oleadas se fija con el inventario de sitios en H1 & Formulario T-18 (Plan de implantación por sitio y por perfil) \\
\addlinespace[2pt]
RT-20.02 & La estrategia de paso a producción será gradual y reversible. Se declarará el criterio de avance entre olas y el procedimiento de reversión, con su tiempo de ejecución. & Sí & Criterio de avance entre oleadas y tiempo de reversión por componente, comprobado en el ensayo de retorno & Formulario T-18 (Tablas 4 y 5) \\
\addlinespace[2pt]
RT-20.03 & Durante la marcha blanca, la solución convivirá con la operación vigente del CLIENTE, con conciliación diaria y sin doble digitación no declarada. & Sí & Condiciones de cierre de marcha blanca y estrategia de migración de cartera con conciliación diaria automatizada de saldos entre plataforma legada y nueva & Subdoc.3 §3.3.2; Anexo A.4.1 (SUP-25) \\
\addlinespace[2pt]
RT-20.04 & El PROPONENTE definirá los indicadores diarios que se medirán durante la marcha blanca y sus umbrales de avance, conforme al Artículo 17.3 de las Bases Administrativas. & Sí & Indicadores diarios de la marcha blanca y umbrales de avance entre oleadas & Formulario T-18 (Tabla 4) \\
\addlinespace[2pt]
RT-20.05 & Se dispondrá de acompañamiento en terreno durante las primeras semanas de cada ola, con dotación declarada y decreciente según la curva de adopción. & Parcial & Acompañamiento en el puesto durante las primeras semanas de cada oleada con dotación decreciente; dotación y duración por declarar con la nivelación del Formulario T-15 & Formulario T-18 (Acompañamiento en terreno y estabilización) \\
\addlinespace[2pt]
RT-20.06 & Se establecerá un período de estabilización con atención reforzada tras cada paso a producción, con dotación y duración declaradas y sin costo adicional. & Parcial & Período de estabilización con atención reforzada tras cada paso a producción (meses 16 y 21), sin costo adicional; dotación y duración por declarar con el Formulario T-15 & Formulario T-18 (Acompañamiento en terreno y estabilización) \\
\addlinespace[2pt]
RT-20.07 & Existirá una definición de terminado acordada con el CLIENTE, aplicable a cada entregable, que incluya código, pruebas, documentación, seguridad y despliegue. & Sí & Definición de terminado acordada con el Cliente para cada entregable: código, pruebas, documentación, seguridad y despliegue & Subdoc.9 \textbar{} Formulario T-9 §T-9.3 \\
\addlinespace[2pt]
RT-20.08 & El protocolo de aceptación de cada hito se formalizará conforme al Formulario T-17, con criterios objetivos y verificables. & Sí & Protocolo de aceptación por hito con criterios objetivos y verificables & Formulario T-17 \\
\addlinespace[2pt]
RT-21.01 & El ADJUDICATARIO dispondrá de un centro de operaciones de red con cobertura 24x7x365, propio o subcontratado, y declarará su ubicación, dotación por turno y procedimientos. & No & Pendiente: la cobertura 24x7x365, su ubicación y la dotación por turno se declaran en el Subdocumento 10 & Subdoc.10 (por redactar) \\
\addlinespace[2pt]
RT-21.02 & Se designará un gerente de servicio dedicado como contraparte permanente del CLIENTE durante la fase de Operación. & No & Pendiente: se declara en el Subdocumento 10 (servicios de operación y niveles de servicio) & Subdoc.10 (por redactar) \\
\addlinespace[2pt]
RT-21.03 & El equipo de operación contará con especialistas por tecnología, nominados y con dedicación declarada. & No & Pendiente: se declara en el Subdocumento 10 (servicios de operación y niveles de servicio) & Subdoc.10 (por redactar) \\
\addlinespace[2pt]
RT-21.04 & Los procedimientos de operación estarán documentados y serán ejecutables por el personal del CLIENTE tras la transferencia de conocimiento. & No & Pendiente: se declara en el Subdocumento 10 (servicios de operación y niveles de servicio) & Subdoc.10 (por redactar) \\
\addlinespace[2pt]
RT-21.05 & El PROPONENTE dispondrá de un centro de atención adecuado para soportar integralmente las consultas, propio o subcontratado con un operador de nivel acreditado. & No & Pendiente: se declara en el Subdocumento 10 (servicios de operación y niveles de servicio) & Subdoc.10 (por redactar) \\
\addlinespace[2pt]
RT-21.06 & El centro de atención cumplirá un tiempo de atención al usuario final de 80\% antes de 20 segundos y una resolución al primer contacto de al menos 70\%. La tasa de abandono no superará el 5\%. & No & Pendiente: se declara en el Subdocumento 10 (servicios de operación y niveles de servicio) & Subdoc.10 (por redactar) \\
\addlinespace[2pt]
RT-21.07 & El horario mínimo de atención será de 8:00 a 20:00 en días hábiles, ampliado a 24x7 para los incidentes de severidad crítica y para la ventana operacional que defina el caso. & Sí & Atención en horario comercial extendido (09:00 a 23:00) para sucursales y centros de distribución, y cobertura 24x7x365 para canal digital y servicios financieros críticos & Subdoc.4 §4.2 \\
\addlinespace[2pt]
RT-21.08 & La atención cubrirá orientación funcional de distintos grados de complejidad, desde preguntas simples hasta situaciones complejas, y las preguntas técnicas más frecuentes sobre la aplicación y su entorno de uso. & No & Pendiente: se declara en el Subdocumento 10 (servicios de operación y niveles de servicio) & Subdoc.10 (por redactar) \\
\addlinespace[2pt]
RT-21.09 & Se mantendrá un registro histórico de las actividades de soporte que permita monitorear y gestionar los principales requerimientos, con análisis de tendencia mensual. & No & Pendiente: se declara en el Subdocumento 10 (servicios de operación y niveles de servicio) & Subdoc.10 (por redactar) \\
\addlinespace[2pt]
RT-21.10 & Existirá un proceso de aprendizaje del centro de soporte que acumule conocimiento sobre las complejidades de la operación y las necesidades de las personas usuarias, reflejado en la base de conocimiento. & No & Pendiente: se declara en el Subdocumento 10 (servicios de operación y niveles de servicio) & Subdoc.10 (por redactar) \\
\addlinespace[2pt]
RT-21.11 & Se proveerán servicios de capacitación en línea en modalidad de autoformación, con registro de avance para efectos de monitoreo. & No & Pendiente: se declara en el Subdocumento 10 (servicios de operación y niveles de servicio) & Subdoc.10 (por redactar) \\
\addlinespace[2pt]
RT-21.12 & Se habilitarán espacios de interacción entre las entidades y personas usuarias del sistema, que permitan compartir experiencias de uso. & No & Pendiente: se declara en el Subdocumento 10 (servicios de operación y niveles de servicio) & Subdoc.10 (por redactar) \\
\addlinespace[2pt]
RT-21.13 & Los indicadores clave del proceso de soporte estarán abiertos al CLIENTE, y los indicadores específicos serán accesibles según nivel de autorización. & No & Pendiente: se declara en el Subdocumento 10 (servicios de operación y niveles de servicio) & Subdoc.10 (por redactar) \\
\addlinespace[2pt]
RT-21.14 & El PROPONENTE dimensionará la dotación del centro de atención con fundamento cuantitativo, empleando teoría de colas o el modelo Erlang C, a partir del volumen de contactos proyectado del caso. & No & Pendiente: se declara en el Subdocumento 10 (servicios de operación y niveles de servicio) & Subdoc.10 (por redactar) \\
\addlinespace[2pt]
RT-21.15 & Existirá un canal único de registro de incidentes y solicitudes, con número de ticket, clasificación por severidad y seguimiento del ciclo de vida completo hasta el cierre conforme. & No & Pendiente: se declara en el Subdocumento 10 (servicios de operación y niveles de servicio) & Subdoc.10 (por redactar) \\
\addlinespace[2pt]
RT-21.16 & Cuando el caso comprenda sitios alejados, los especialistas de niveles 2 y 3 deberán trasladarse cuando la resolución lo requiera, por el medio más rápido disponible y con disponibilidad de tiempo suficiente, sin alterar la atención normal del resto de los sitios. El costo del traslado está incluido en la oferta. & No & Pendiente: se declara en el Subdocumento 10 (servicios de operación y niveles de servicio) & Subdoc.10 (por redactar) \\
\addlinespace[2pt]
RT-21.17 & El cierre de un ticket requerirá confirmación de la persona usuaria o transcurso del plazo de confirmación automática declarado. & No & Pendiente: se declara en el Subdocumento 10 (servicios de operación y niveles de servicio) & Subdoc.10 (por redactar) \\
\addlinespace[2pt]
RT-21.18 & El ADJUDICATARIO reportará mensualmente el cumplimiento de los niveles de servicio de atención, con el detalle por severidad y el análisis de los incumplimientos. & No & Pendiente: se declara en el Subdocumento 10 (servicios de operación y niveles de servicio) & Subdoc.10 (por redactar) \\
\addlinespace[2pt]
RT-21.19 & El PROPONENTE declarará el tamaño de la bolsa anual de horas de mantención evolutiva, su composición por perfil y el procedimiento de solicitud, estimación, aprobación y liquidación. & No & Pendiente: se declara en el Subdocumento 10 (servicios de operación y niveles de servicio) & Subdoc.10 (por redactar) \\
\addlinespace[2pt]
RT-21.20 & La mantención evolutiva se someterá al mismo estándar de calidad, pruebas y seguridad que el desarrollo original. & No & Pendiente: se declara en el Subdocumento 10 (servicios de operación y niveles de servicio) & Subdoc.10 (por redactar) \\
\addlinespace[2pt]
RT-21.21 & El PROPONENTE mantendrá un registro de deuda técnica cuantificado y destinará una fracción declarada de la capacidad de la Operación a reducirla. & No & Pendiente: se declara en el Subdocumento 10 (servicios de operación y niveles de servicio) & Subdoc.10 (por redactar) \\
\addlinespace[2pt]
RT-21.22 & Las actualizaciones de versión de los componentes de base se planificarán anualmente, con ventana acordada y plan de reversión. & No & Pendiente: se declara en el Subdocumento 10 (servicios de operación y niveles de servicio) & Subdoc.10 (por redactar) \\
\addlinespace[2pt]
RT-22.01 & El plan de capacitación se estructurará por perfil: personas usuarias finales, usuarias avanzadas, administradoras, equipo técnico y soporte de niveles 1 y 2. & Sí & Plan de capacitación por perfil: personal de tienda, dueños de producto, TI, repositores externos y soporte de niveles 1 y 2 & Formulario T-9 §T-9.4 (Estrategias de involucramiento por dimensión) \\
\addlinespace[2pt]
RT-22.02 & Las modalidades incluirán capacitación presencial en cada sitio de operación, sesiones en línea sincrónicas, autoformación y acompañamiento en puesto de trabajo durante la marcha blanca. & Sí & Capacitación presencial, sincrónica, autoformación y acompañamiento en el puesto & Formulario T-9 §T-9.4 (Estrategias de involucramiento por dimensión) \\
\addlinespace[2pt]
RT-22.03 & Todo el material de capacitación se entregará en español, en formato editable y de propiedad del CLIENTE: manuales por perfil, guías rápidas, preguntas frecuentes, videos tutoriales y base de conocimiento consultable. & Sí & Material editable, en español y de propiedad de Ancoa & Formulario T-9 §T-9.4 (Estrategias de involucramiento por dimensión) \\
\addlinespace[2pt]
RT-22.04 & La capacitación no podrá afectar la operación del CLIENTE: se programará por turnos y en horarios acordados, considerando la estacionalidad y la ventana operacional del caso. & Parcial & Formación programada por perfil y turno, y sin intervenciones dentro de los congelamientos; la estacionalidad por área se declara en el plan de capacitación & Formulario T-9 §T-9.4 \textbar{} Formulario T-18 \\
\addlinespace[2pt]
RT-22.05 & Las personas usuarias administradoras y el equipo técnico del CLIENTE serán evaluados y certificados como condición para el cierre de cada marcha blanca. & Sí & Evaluación y certificación del personal como condición de cierre de cada marcha blanca (paquete 1.13.3) & Formulario T-18 (Condiciones de cierre de cada marcha blanca) \textbar{} Formulario T-14 1.13.3 \\
\addlinespace[2pt]
RT-22.06 & Durante la Operación se ejecutarán al menos dos jornadas anuales de actualización y se capacitará al personal nuevo del CLIENTE sin costo adicional. & Parcial & Dos jornadas anuales de actualización en operación; la capacitación del personal nuevo sin costo adicional se declara en el Subdocumento 11 & Formulario T-9 §T-9.4 \textbar{} Subdoc.11 (por redactar) \\
\addlinespace[2pt]
RT-22.07 & El ADJUDICATARIO proveerá acompañamiento y mentoría al equipo técnico del CLIENTE durante al menos seis meses posteriores al paso a producción de la Etapa 2. & Sí & Mentoría técnica de al menos seis meses tras el paso a producción de la Etapa 2 & Formulario T-9 §T-9.4 (Estrategias de involucramiento por dimensión) \\
\addlinespace[2pt]
RT-22.08 & La base de conocimiento será mantenida y actualizada durante todo el Contrato, con métricas de uso y de utilidad percibida. & Parcial & Base de conocimiento de la mesa de servicio, revisada mensualmente; métricas de uso y de utilidad percibida por declarar & Subdoc.1 §1.4 \textbar{} Subdoc.11 (por redactar) \\
\addlinespace[2pt]
RT-22.09 & Se valorará la existencia de un ambiente permanente de entrenamiento, con datos ficticios, disponible para la práctica sin riesgo. & Parcial & Ambiente de entrenamiento permanente en el plan de capacitación; datos ficticios por declarar & Formulario T-9 §T-9.4 \\
\addlinespace[2pt]
RT-23.01 & Información institucional: descripción detallada del giro principal, historia y trayectoria de la empresa con al menos tres años de antigüedad comprobable, misión, visión y valores, certificaciones y acreditaciones vigentes, y presencia geográfica y oficinas. & Parcial & Sitio web corporativo declarado, que según el proponente aloja esta información institucional (giro, historia de 10 años, misión, visión, valores, certificaciones) & Subdocumento 1, Anexo A.1 ``Presencia digital y web corporativa'' y secciones 1.1, 1.3 y 1.5 \\
\addlinespace[2pt]
RT-23.02 & Experiencia y casos de éxito: portafolio de proyectos similares de los últimos cinco años, casos documentados con métricas verificables, testimonios de clientes con autorización de publicación, industrias atendidas con énfasis en la del caso asignado, y volumen de transacciones o de personas usuarias gestionadas. & Sí & Portafolio de tres proyectos de misión crítica con métricas verificables, declarado como disponible en la plataforma oficial & Subdocumento 1, sección 1.4 ``Experiencia y Certificaciones'' y Anexo A.1 \\
\addlinespace[2pt]
RT-23.03 & Equipo profesional: organigrama del equipo directivo, perfiles de socios y directores, currículos resumidos del equipo técnico clave y certificaciones profesionales del personal. & Sí & Organigrama y roles estratégicos del equipo de jefaturas, declarados como parte de la información disponible en el sitio corporativo; currículos y certificaciones del personal en el Formulario T-8 & Subdocumento 1, secciones 1.2 y 1.5; Anexo A.1; Subdocumento 12 \\
\addlinespace[2pt]
RT-23.04 & Capacidades técnicas: servicios y soluciones ofrecidas, conjunto tecnológico dominado, alianzas tecnológicas con proveedores de nube y de plataforma, metodologías de trabajo certificadas, e infraestructura y capacidad instalada. & Sí & Capacidades instaladas (fábrica de software, célula de datos, célula de ciberseguridad, mesa de servicio) y alianza tecnológica con Microsoft Azure & Subdocumento 1, secciones 1.2 y 1.5; Anexo A.1 \\
\addlinespace[2pt]
RT-23.05 & El sitio será accesible conforme a WCAG 2.2 nivel AA, responsivo y con certificado TLS válido y vigente. & No & Pendiente: el sitio web del proponente & Subdoc.1 Anexo A.1 \\
\addlinespace[2pt]
RT-23.06 & El sitio se mantendrá activo y disponible durante todo el proceso de licitación, desde el registro del participante hasta la adjudicación. & No & Pendiente: el sitio web del proponente & Subdoc.1 Anexo A.1 \\
\addlinespace[2pt]
RT-23.07 & Material educativo digital: seminarios en línea, artículos técnicos o centro de recursos y documentación. & No & Pendiente: el sitio web del proponente & Subdoc.1 Anexo A.1 \\
\addlinespace[2pt]
RT-23.08 & Demostración en línea, recorrido virtual de la solución o portal de soporte para clientes. & No & Pendiente: el sitio web del proponente & Subdoc.1 Anexo A.1 \\
\addlinespace[2pt]
RT-23.09 & Métricas de disponibilidad y desempeño de los servicios del proponente publicadas en tiempo real. & No & Pendiente: el sitio web del proponente & Subdoc.1 Anexo A.1 \\
\addlinespace[2pt]
RT-23.10 & Calculadora de retorno de la inversión o herramienta de estimación pertinente a la industria del caso. & No & Pendiente: el sitio web del proponente & Subdoc.1 Anexo A.1 \\
\addlinespace[2pt]
RT-24.01 & Todos los integrantes clave del equipo deberán aparecer individualmente presentándose, en toma de medio cuerpo, mirando directamente a la cámara, con vestimenta formal o de negocio informal, y con una duración mínima de aparición de 10 segundos por persona. & No & Pendiente: el video de presentación del equipo, que se entrega como pieza separada del Sobre 2 & Pieza audiovisual del Sobre 2 (por producir) \\
\addlinespace[2pt]
RT-24.02 & Cada vez que aparezca un integrante deberá mostrarse en pantalla su nombre completo, su cargo en el proyecto, sus años de experiencia y su especialización relevante. & No & Pendiente: el video de presentación del equipo, que se entrega como pieza separada del Sobre 2 & Pieza audiovisual del Sobre 2 (por producir) \\
\addlinespace[2pt]
RT-24.03 & Las certificaciones principales de cada integrante podrán mostrarse en pantalla. & No & Pendiente: el video de presentación del equipo, que se entrega como pieza separada del Sobre 2 & Pieza audiovisual del Sobre 2 (por producir) \\
\addlinespace[2pt]
RT-24.04 & El fondo será el mismo para todas las tomas de personas, con iluminación consistente y encuadre similar para todos los participantes. & No & Pendiente: el video de presentación del equipo, que se entrega como pieza separada del Sobre 2 & Pieza audiovisual del Sobre 2 (por producir) \\
\addlinespace[2pt]
RT-24.05 & La paleta de colores corporativa, la tipografía y el uso de la marca serán consistentes en todo el video, con logotipo visible en todas las escenas y datos de contacto en el encabezado o pie. & No & Pendiente: el video de presentación del equipo, que se entrega como pieza separada del Sobre 2 & Pieza audiovisual del Sobre 2 (por producir) \\
\addlinespace[2pt]
RT-24.06 & Los niveles de volumen estarán normalizados, con la misma calidad de grabación y sin variaciones bruscas de sonido. & No & Pendiente: el video de presentación del equipo, que se entrega como pieza separada del Sobre 2 & Pieza audiovisual del Sobre 2 (por producir) \\
\addlinespace[2pt]
RT-25.01 & Navegación intuitiva sin necesidad de manual, con un máximo de tres interacciones para alcanzar cualquier funcionalidad principal. & No & Pendiente: el prototipo navegable, que se entrega como pieza separada del Sobre 2 & Prototipo del Sobre 2 (por producir) \\
\addlinespace[2pt]
RT-25.02 & Retroalimentación visual clara ante cada acción, estados de carga y transición fluidos, y manejo elegante de los errores. & No & Pendiente: el prototipo navegable, que se entrega como pieza separada del Sobre 2 & Prototipo del Sobre 2 (por producir) \\
\addlinespace[2pt]
RT-25.03 & Cumplimiento de WCAG 2.2 nivel AA: contraste adecuado, tamaños de fuente legibles, objetivos táctiles de al menos 44x44 píxeles, navegación completa por teclado, y textos alternativos y etiquetas de accesibilidad. & No & Pendiente: el prototipo navegable, que se entrega como pieza separada del Sobre 2 & Prototipo del Sobre 2 (por producir) \\
\addlinespace[2pt]
RT-25.04 & Diseño responsivo demostrado en escritorio (1920x1080 y 1366x768), tableta (768 x 1024 en vertical y horizontal) y teléfono (375x812), con puntos de quiebre coherentes. & No & Pendiente: el prototipo navegable, que se entrega como pieza separada del Sobre 2 & Prototipo del Sobre 2 (por producir) \\
\addlinespace[2pt]
RT-25.05 & Sistema de diseño documentado: paleta de a lo más cinco colores principales, a lo más dos familias tipográficas jerarquizadas, iconografía coherente, retícula y espaciado consistentes, y componentes reutilizables. & No & Pendiente: el prototipo navegable, que se entrega como pieza separada del Sobre 2 & Prototipo del Sobre 2 (por producir) \\
\addlinespace[2pt]
RT-25.06 & El prototipo se entregará como enlace navegable en línea, sin instalación ni complementos, compatible con los navegadores modernos vigentes. & No & Pendiente: el prototipo navegable, que se entrega como pieza separada del Sobre 2 & Prototipo del Sobre 2 (por producir) \\
\addlinespace[2pt]
RT-25.07 & El acceso al prototipo estará garantizado por un mínimo de seis meses desde su entrega. & No & Pendiente: el prototipo navegable, que se entrega como pieza separada del Sobre 2 & Prototipo del Sobre 2 (por producir) \\
\addlinespace[2pt]
RT-25.08 & El prototipo permitirá navegación completa entre pantallas, simulación de ingreso de datos, transiciones y animaciones básicas, estados de posado y activo, simulación de carga de archivos, elementos desplegables funcionales y simulación de validaciones. & No & Pendiente: el prototipo navegable, que se entrega como pieza separada del Sobre 2 & Prototipo del Sobre 2 (por producir) \\
\addlinespace[2pt]
RT-25.09 & Se entregará la arquitectura de información del prototipo: mapa del sitio completo, diagrama de navegación, taxonomía y nomenclatura, estructura de menús y jerarquía de la información. & No & Pendiente: el prototipo navegable, que se entrega como pieza separada del Sobre 2 & Prototipo del Sobre 2 (por producir) \\
\addlinespace[2pt]
RT-25.10 & Se valorará la posibilidad de exportar el prototipo a un documento interactivo para revisión sin conexión. & No & Pendiente: el prototipo navegable, que se entrega como pieza separada del Sobre 2 & Prototipo del Sobre 2 (por producir) \\
\addlinespace[2pt]
RT-26.01 & Cada innovación se ubicará explícitamente en la arquitectura: qué capa la contiene, qué componentes la implementan y qué interfaces consume o expone. & Sí & Las 5 fichas del Formulario T-19 declaran explícitamente el campo ``Dónde se inserta en la arquitectura'' (capa, zona y módulo) para cada una de las cinco innovaciones & Subdocumento 13, Anexos A.1 a A.5 (campo ``Dónde se inserta en la arquitectura'') \\
\addlinespace[2pt]
RT-26.02 & Cada innovación identificará los paquetes de la estructura de descomposición del trabajo que la ejecutan y el mes del cronograma en que se materializa. & Sí & Paquetes 1.10.1 a 1.10.5 de la estructura de desglose del trabajo y mes de materialización en cada ficha & Formulario T-19 (campos 8 y 9) \textbar{} Formulario T-14 (rama 1.10) \\
\addlinespace[2pt]
RT-26.03 & Las innovaciones de base tecnológica declararán el nivel de madurez de la tecnología con la escala utilizada y citarán las fuentes en norma APA 7.ª edición. & Sí & Cada una de las cinco fichas declara ``Nivel de madurez y escala utilizada'' (escala TRL) y ``Fuentes citadas (APA 7.a ed.)'' con referencias reales & Subdocumento 13, Anexos A.1 a A.5 (campos ``Nivel de madurez y escala utilizada'' y ``Fuentes citadas'') \\
\addlinespace[2pt]
RT-26.04 & Cada innovación declarará su riesgo de adopción, su probabilidad, su impacto, la estrategia de mitigación y el plan de contingencia si no rinde lo esperado. & Sí & Riesgo de adopción, probabilidad, impacto, mitigación y contingencia por innovación & Formulario T-19 (campos 15 a 17) \\
\addlinespace[2pt]
RT-26.05 & Cada innovación declarará su indicador de verificación con línea base, meta y momento de medición, y su impacto en inversión, costo operacional y beneficio esperado. & Sí & Indicador de verificación con línea base, meta y momento de medición, e impacto en inversión, costo operacional y beneficio & Formulario T-19 (campos 10 a 14) \\
\addlinespace[2pt]
RT-26.06 & Las innovaciones que incorporen inteligencia artificial cumplirán íntegramente el Capítulo 18 de este documento. & Parcial & Las innovaciones con IA (conteo predictivo y búsqueda multimodal) cumplen el Capítulo 18 en lo declarado en las fichas; la cobertura completa depende de RT-18.01 a RT-18.10 & Subdoc.13 \textbar{} Formulario T-19 \\
\addlinespace[2pt]
RT-26.07 & Las innovaciones que modifiquen la arquitectura de seguridad requerirán su propio modelado de amenazas. & Parcial & Modelado de amenazas propio para lockers, búsqueda multimodal y Retail Media antes de integrar & Formulario T-19 (campo 16, fichas 1, 3 y 4) \\
\addlinespace[2pt]
RT-26.08 & Se valorará que al menos una innovación sea verificable durante la marcha blanca de la Etapa 1, es decir, que su beneficio pueda medirse antes del mes 16. & Sí & Conteo cíclico predictivo evaluado en sombra en los meses 13 a 15, antes del paso a producción del mes 16 & Formulario T-19 (ficha 2) \textbar{} Subdoc.13 §13.2 \\
\addlinespace[2pt]
\end{longtable}

\ossFinalPage

\end{document}
